% Produire une lettre officielle : la lettre de motivation — Français 1re Tech — Cours signé OnBuch+
% Programme officiel MINESEC. Compiler avec : tectonic N12-produire-lettre-officielle-lettre-motivation.tex
\newcommand{\DOCMATIERE}{Français}
\newcommand{\DOCNIVEAU}{1re Tech}
\newcommand{\DOCMODULE}{Français – classe de Première technique}
\newcommand{\DOCLECON}{N12}
\newcommand{\DOCTITRE}{Produire une lettre officielle : la lettre de motivation}
% OnBuch+ — préambule « cours signés » ENRICHI (compilé avec tectonic / XeLaTeX)
% Système visuel « Pop » d'OnBuch+ : fond crème (jamais blanc), bordures encre
% épaisses, ombres « dures » décalées, titres Archivo Black en capitales.
% Version MATHÉMATIQUES : graphes (pgfplots/tikz), boîtes pédagogiques
% (définition, propriété, méthode, attention, le savais-tu, exercice résolu,
% exercice, TP, bilan), page de garde et signature OnBuch+ ancrée en bas.
%
% Macros attendues dans main.tex AVANT \input{preamble} :
%   \DOCMATIERE  \DOCNIVEAU  \DOCMODULE  \DOCLECON (numéro)  \DOCTITRE
\documentclass[11pt]{article}

\usepackage{fontspec}
\usepackage[french]{babel}
\usepackage[a4paper,margin=2cm,top=3.4cm,bottom=2.8cm,headheight=1.4cm,footskip=1.3cm]{geometry}
\usepackage{amsmath,amssymb,amsfonts}
\usepackage{mathtools} % \xmapsto, \coloneqq... souvent utilisés par les modèles
\usepackage{cancel} % simplifications barrées (bilans de forces)
% \abs{z} (module, valeur absolue) et \norm{u} (norme), utilisés par les modèles
\providecommand{\abs}{}\let\abs\relax\DeclarePairedDelimiter{\abs}{\lvert}{\rvert}
\providecommand{\norm}{}\let\norm\relax\DeclarePairedDelimiter{\norm}{\lVert}{\rVert}
\usepackage[table]{xcolor} % [table] : fournit aussi \rowcolors (alternance de lignes)
\usepackage{pagecolor}
\usepackage{tikz}
\usetikzlibrary{arrows.meta,positioning,calc,shapes.geometric,shapes.misc,shapes.arrows,decorations.pathmorphing,decorations.pathreplacing,decorations.markings,patterns,shadows,mindmap,trees,fit,backgrounds,matrix,intersections,angles,quotes}
\usepackage{pgfplots}
\usepackage[version=4]{mhchem} % équations nucléaires \ce{^{235}_{92}U}
\usepackage[european,siunitx]{circuitikz} % schémas électriques
\pgfplotsset{compat=1.18}
\usepgfplotslibrary{fillbetween,groupplots} % aires entre courbes (calcul intégral)
\usepackage{enumitem}
\usepackage{fancyhdr}
% [most] charge toutes les bibliothèques tcolorbox (listings, minted, raster,
% documentation...) et gonfle inutilement la mémoire TeX sur un document long
% (~300 boîtes) : on ne charge que ce qui est réellement utilisé.
\usepackage[skins,breakable]{tcolorbox}
\usepackage{graphicx}
\usepackage{colortbl}
\usepackage{booktabs}
\usepackage{tabularx}
\usepackage{multirow}
\usepackage{diagbox} % case d'en-tête diagonale des tableaux à double entrée
% Colonnes extensibles centrée / gauche / droite (C, L, R), souvent utilisées par les modèles
\newcolumntype{C}{>{\centering\arraybackslash}X}
\newcolumntype{L}{>{\raggedright\arraybackslash}X}
\newcolumntype{R}{>{\raggedleft\arraybackslash}X}
\usepackage{array}
\usepackage{multicol}
\usepackage{siunitx}
% parse-numbers=false : accepte aussi \SI{2,5 \times 10^{-3}}{...}
\sisetup{locale=FR,output-decimal-marker={,},group-minimum-digits=4,parse-numbers=false}
\DeclareSIUnit\ohm{\ensuremath{\symup{\Omega}}} % U+2126 absent de la police
\DeclareSIUnit\division{div} % réglages d'oscilloscope (ms/div, V/div)
\DeclareSIUnit\tr{tr} % tours (vitesse de rotation, ex. \SI{1500}{\tr\per\minute})
\DeclareSIUnit\molar{M} % concentration molaire (ex. \SI{3}{\molar}, \SI{1}{\micro\molar})
\DeclareSIUnit\an{an} % année (le modèle confond parfois avec \year, un compteur TeX) — ex. \SI{5}{\kilo\meter\per\an}
\DeclareSIUnit\FCFA{FCFA} % franc CFA (ex. \SI{500}{\FCFA\per\kilo\gram})
\DeclareSIUnit\degreeBrix{\textdegree Brix} % teneur en sucre d'un sirop/jus (réfractométrie)
\DeclareSIUnit\month{mois} % le modèle utilise parfois \month, un compteur TeX, comme unité de durée
\DeclareSIUnit\microliter{\micro\liter} % le modèle écrit parfois \microliter en un seul token
\DeclareSIUnit\million{millions} % dénombrement (ex. \SI{15}{\million\per\mL} spermatozoïdes)
\usepackage{qrcode}
% \degree : commande fréquemment utilisée par les modèles pour un angle ou une
% température en dehors de \SI{}{} (non fournie par les paquets chargés).
\providecommand{\degree}{^{\circ}}
% \qed (fin de démonstration), utilisé par les modèles sans amsthm
\providecommand{\qed}{\hfill$\square$}
% \barre : double barre des valeurs interdites dans un tableau de variations (tkz-tab)
\providecommand{\barre}{\ensuremath{\|}}
% environnement proof (amsthm non chargé) utilisé par les modèles
\ifdefined\proof\else\newenvironment{proof}[1][Démonstration]{\par\noindent\textit{#1.}\ }{\qed\par}\fi
\usepackage{lastpage}
\usepackage{hyperref}
\hypersetup{hidelinks}

% ============================================================
% Palette « Pop » OnBuch+ — mêmes hex que la classe P (plus_ui.dart)
% ============================================================
\definecolor{poporange}{HTML}{F59321}
\definecolor{popdark}{HTML}{E07A0C}
\definecolor{popcream}{HTML}{FFF6E8}
\definecolor{popink}{HTML}{1C1714}
\definecolor{poppurple}{HTML}{7A5AE0}
\definecolor{popgreen}{HTML}{2E9E6A}
\definecolor{poppink}{HTML}{E0457B}
\definecolor{popblue}{HTML}{2F7DE1}
\definecolor{popgold}{HTML}{C98A12}
\definecolor{popmuted}{HTML}{8A7F76}
\definecolor{popline}{HTML}{EADFD2}
\definecolor{popgray}{HTML}{8A7F76}
\colorlet{popgrey}{popgray}
% Alias tolérants (le modèle nomme parfois les couleurs différemment)
\colorlet{popwhite}{white}
\colorlet{popblack}{popink}
\colorlet{popred}{poppink}
\colorlet{popyellow}{popgold}
% Teintes claires (fonds de tableaux/encadrés) — le modèle nomme parfois
% les couleurs avec un suffixe L (ex. popblueL) sans les avoir définies.
\colorlet{poporangeL}{poporange!20}
\colorlet{popdarkL}{popdark!20}
\colorlet{poppurpleL}{poppurple!20}
\colorlet{popgreenL}{popgreen!20}
\colorlet{poppinkL}{poppink!20}
\colorlet{popblueL}{popblue!20}
\colorlet{popgoldL}{popgold!20}
\colorlet{popredL}{popred!20}
\colorlet{popgrayL}{popgray!20}
% Teintes claires pour fonds de tableaux / zones
\colorlet{popblueL}{popblue!10!white}
\colorlet{poporangeL}{poporange!14!white}
\colorlet{popgreenL}{popgreen!12!white}
\colorlet{poppurpleL}{poppurple!12!white}
\colorlet{poppinkL}{poppink!12!white}
\colorlet{popgoldL}{popgold!14!white}

\pagecolor{popcream}

% ============================================================
% Typographie
% ============================================================
\defaultfontfeatures{Ligatures=TeX}
\setmainfont{PlusJakartaSans-Medium.ttf}[Path=../../../fonts/,BoldFont=PlusJakartaSans-Bold.ttf]
% Maths en sans-serif (Fira Math) avec chiffres et lettres droites de Plus
% Jakarta, pour que les formules se fondent dans le texte.
\usepackage[math-style=ISO,bold-style=ISO]{unicode-math}
\setmathfont{FiraMath-Regular.otf}[Path=../../../fonts/]
\setmathfont{PlusJakartaSans-Medium.ttf}[Path=../../../fonts/,range={up/num,up/{Latin,latin},\mathup}]
\usepackage{icomma} % virgule décimale sans espace en mode math
% Caractères mathématiques Unicode absents des polices de texte (Plus Jakarta,
% Archivo Black) : rendus par leur équivalent mathématique, en texte comme en
% formule — sinon ils s'affichent en carré vide (ex. « Arithmétique dans ℤ »).
\usepackage{newunicodechar}
\newunicodechar{ℝ}{\ensuremath{\mathbb{R}}}
\newunicodechar{ℤ}{\ensuremath{\mathbb{Z}}}
\newunicodechar{ℂ}{\ensuremath{\mathbb{C}}}
\newunicodechar{ℕ}{\ensuremath{\mathbb{N}}}
\newunicodechar{ℚ}{\ensuremath{\mathbb{Q}}}
\newunicodechar{𝔻}{\ensuremath{\mathbb{D}}}
\newunicodechar{α}{\ensuremath{\alpha}}
\newunicodechar{β}{\ensuremath{\beta}}
\newunicodechar{γ}{\ensuremath{\gamma}}
\newunicodechar{δ}{\ensuremath{\delta}}
\newunicodechar{ε}{\ensuremath{\varepsilon}}
\newunicodechar{θ}{\ensuremath{\theta}}
\newunicodechar{λ}{\ensuremath{\lambda}}
\newunicodechar{μ}{\ensuremath{\mu}}
\newunicodechar{σ}{\ensuremath{\sigma}}
\newunicodechar{τ}{\ensuremath{\tau}}
\newunicodechar{φ}{\ensuremath{\varphi}}
\newunicodechar{ψ}{\ensuremath{\psi}}
\newunicodechar{ω}{\ensuremath{\omega}}
\newunicodechar{Σ}{\ensuremath{\Sigma}}
% majuscules produites par \MakeUppercase dans les titres (α-aminés -> Α)
\newunicodechar{Α}{\ensuremath{\alpha}}
\newunicodechar{Β}{\ensuremath{\beta}}
\newunicodechar{Γ}{\ensuremath{\Gamma}}
\newunicodechar{Δ}{\ensuremath{\Delta}}
\newunicodechar{Ε}{\ensuremath{\varepsilon}}
\newunicodechar{Θ}{\ensuremath{\Theta}}
\newunicodechar{Λ}{\ensuremath{\Lambda}}
\newunicodechar{Μ}{\ensuremath{\mu}}
\newunicodechar{Τ}{\ensuremath{\tau}}
\newunicodechar{Φ}{\ensuremath{\Phi}}
\newunicodechar{Ψ}{\ensuremath{\Psi}}
\newunicodechar{Ω}{\ensuremath{\Omega}}
\newunicodechar{↦}{\ensuremath{\mapsto}}
\newunicodechar{∈}{\ensuremath{\in}}
\newunicodechar{∉}{\ensuremath{\notin}}
\newunicodechar{∧}{\ensuremath{\wedge}}
\newunicodechar{⇔}{\ensuremath{\Leftrightarrow}}
\newunicodechar{⟺}{\ensuremath{\Longleftrightarrow}}
\newunicodechar{⇒}{\ensuremath{\Rightarrow}}
\newunicodechar{⟹}{\ensuremath{\Longrightarrow}}
\newunicodechar{∘}{\ensuremath{\circ}}
\newunicodechar{∪}{\ensuremath{\cup}}
\newunicodechar{∩}{\ensuremath{\cap}}
\newunicodechar{≡}{\ensuremath{\equiv}}
\newunicodechar{⊂}{\ensuremath{\subset}}
\newunicodechar{⊥}{\ensuremath{\perp}}
\newunicodechar{∃}{\ensuremath{\exists}}
\newunicodechar{∀}{\ensuremath{\forall}}
\newunicodechar{∛}{\ensuremath{\sqrt[3]{\,}}}
\newunicodechar{∜}{\ensuremath{\sqrt[4]{\,}}}
\newunicodechar{⃗}{\ensuremath{{}^{\to}}}
\newunicodechar{ⁿ}{\ifmmode^{n}\else\textsuperscript{\ensuremath{n}}\fi}
\newunicodechar{ˣ}{\ifmmode^{x}\else\textsuperscript{\ensuremath{x}}\fi}
\newunicodechar{ᵇ}{\ifmmode^{b}\else\textsuperscript{\ensuremath{b}}\fi}
\newunicodechar{ʳ}{\ifmmode^{r}\else\textsuperscript{\ensuremath{r}}\fi}
\newunicodechar{ᵘ}{\ifmmode^{u}\else\textsuperscript{\ensuremath{u}}\fi}
\newunicodechar{ʸ}{\ifmmode^{y}\else\textsuperscript{\ensuremath{y}}\fi}
\newunicodechar{ᵃ}{\ifmmode^{a}\else\textsuperscript{\ensuremath{a}}\fi}
\newunicodechar{ᵉ}{\ifmmode^{e}\else\textsuperscript{\ensuremath{e}}\fi}
\newunicodechar{ᵏ}{\ifmmode^{k}\else\textsuperscript{\ensuremath{k}}\fi}
\newunicodechar{ᵖ}{\ifmmode^{p}\else\textsuperscript{\ensuremath{p}}\fi}
\newunicodechar{ᴺ}{\ifmmode^{N}\else\textsuperscript{\ensuremath{N}}\fi}
\newunicodechar{ᶜ}{\ifmmode^{c}\else\textsuperscript{\ensuremath{c}}\fi}
\newunicodechar{ʰ}{\ifmmode^{h}\else\textsuperscript{\ensuremath{h}}\fi}
\newunicodechar{ᵠ}{\ifmmode^{q}\else\textsuperscript{\ensuremath{q}}\fi}
\newunicodechar{ₙ}{\ifmmode_{n}\else\textsubscript{\ensuremath{n}}\fi}
\newunicodechar{ₐ}{\ifmmode_{a}\else\textsubscript{\ensuremath{a}}\fi}
\newunicodechar{ᵢ}{\ifmmode_{i}\else\textsubscript{\ensuremath{i}}\fi}
\newunicodechar{ᵣ}{\ifmmode_{r}\else\textsubscript{\ensuremath{r}}\fi}
\newunicodechar{ₖ}{\ifmmode_{k}\else\textsubscript{\ensuremath{k}}\fi}
\newunicodechar{ᵦ}{\ifmmode_{\beta}\else\textsubscript{\ensuremath{\beta}}\fi}
\newunicodechar{ₓ}{\ifmmode_{x}\else\textsubscript{\ensuremath{x}}\fi}
\newfontfamily\popdisplay{ArchivoBlack-Regular.ttf}[Path=../../../fonts/]
\newfontfamily\poplogo{SpaceGrotesk-Bold.ttf}[Path=../../../fonts/]
\newfontfamily\popsemi{PlusJakartaSans-SemiBold.ttf}[Path=../../../fonts/]
\newfontfamily\popxbold{PlusJakartaSans-ExtraBold.ttf}[Path=../../../fonts/]
\newfontfamily\popxboldls{PlusJakartaSans-ExtraBold.ttf}[Path=../../../fonts/,LetterSpace=3.0]

\setlist[enumerate]{leftmargin=1.7em,itemsep=5pt,topsep=4pt}
\setlist[itemize]{leftmargin=1.7em,itemsep=4pt,topsep=4pt,label={\color{poporange}\textbullet}}
\setlength{\parindent}{0pt}
\setlength{\parskip}{5pt}
\renewcommand{\arraystretch}{1.25}

% pgfplots : style Pop par défaut
\pgfplotsset{
  every axis/.append style={axis line style={popink,line width=1pt},tick style={popink},
    grid style={popline,line width=0.5pt},label style={font=\small},tick label style={font=\footnotesize}},
  popcurve/.style={poporange,line width=1.8pt,no markers,smooth},
}
% Même style disponible dans un \draw TikZ simple (hors pgfplots)
\tikzset{popcurve/.style={poporange,line width=1.8pt}}
\tikzset{popgrid/.style={popline,very thin}}
\colorlet{popcurve}{poporange} % le nom du style est parfois utilisé comme couleur

% ============================================================
% Pill / squiggle
% ============================================================
\newcommand{\poppill}[3]{%
  \tikz[baseline=-0.55ex]\node[draw=popink,line width=1.2pt,rounded corners=2.6pt,%
    fill=#2,inner xsep=6.5pt,inner ysep=3pt]{%
    \popxboldls\fontsize{7}{7}\selectfont\color{#3}\MakeUppercase{#1}};%
}
\newcommand{\popsquiggle}[1]{%
  \tikz[baseline=-0.4ex]{
    \draw[#1,line width=2.6pt,line cap=round]
      plot [smooth] coordinates {(0,0.09) (0.34,0.01) (0.68,0.21) (1.02,0.01) (1.36,0.10)};
  }%
}
% Mot-clé surligné (marqueur orange sous le texte)
\newcommand{\cle}[1]{{\popxbold\color{popdark}#1}}

% ============================================================
% En-tête / pied
% ============================================================
\pagestyle{fancy}
\fancyhf{}
\renewcommand{\headrulewidth}{0pt}
\renewcommand{\footrulewidth}{0pt}
\lhead{\raisebox{-0.15ex}{\poplogo\fontsize{13}{13}\selectfont\color{popink}OnBuch\color{poporange}+}}
\rhead{\poppill{\DOCMATIERE\ \textbullet\ \DOCNIVEAU\ \textbullet\ Leçon \DOCLECON}{white}{popink}}
\lfoot{\popxbold\fontsize{7.6}{7.6}\selectfont\color{popmuted}\MakeUppercase{Cours signé OnBuch+}\ \textbullet\ {\popsemi\DOCTITRE}}
\rfoot{\tikz[baseline=-0.6ex]\node[rounded corners=8.5pt,fill=popink,minimum height=17pt,minimum width=17pt,inner xsep=5pt,inner sep=0]{\popxbold\fontsize{8}{8}\selectfont\color{popcream}\thepage\,/\,\pageref*{LastPage}};}

% ============================================================
% Titres
% ============================================================
\newcommand{\chaptitre}[2]{%
  \begin{tcolorbox}[enhanced,colback=poporange,colframe=popink,boxrule=2.6pt,arc=11pt,
      drop shadow=popink,left=15pt,right=15pt,top=12pt,bottom=11pt,before skip=3pt,after skip=18pt]
    \poppill{#2}{popink}{popcream}\par\vspace{9pt}
    {\raggedright\hyphenpenalty=10000\exhyphenpenalty=10000\popdisplay\fontsize{22}{24}\selectfont\color{popcream}\MakeUppercase{#1}\par}
  \end{tcolorbox}
}
% Section numérotée : pastille orange avec le numéro + titre Archivo
\newcounter{popsec}
\newcommand{\coursec}[1]{%
  \stepcounter{popsec}\setcounter{popsub}{0}%
  \par\vspace{16pt}\noindent
  \begin{minipage}{\linewidth}
  \tikz[baseline=-0.7ex]\node[draw=popink,line width=1.6pt,fill=poporange,rounded corners=4pt,minimum size=22pt,inner sep=0]{\popdisplay\fontsize{12}{12}\selectfont\color{popcream}\thepopsec};%
  \hspace{8pt}{\popdisplay\fontsize{15}{17}\selectfont\color{popink}\MakeUppercase{#1}}\par
  \vspace{4pt}\noindent{\color{poporange}\rule{36pt}{3.6pt}}
  \end{minipage}\par\nopagebreak\vspace{8pt}
}
\newcounter{popsub}
\newcommand{\courssub}[1]{%
  \stepcounter{popsub}%
  \par\vspace{9pt}\noindent{\popxbold\fontsize{11.5}{13}\selectfont\color{popink}{\color{poporange}\thepopsec.\thepopsub}\hspace{6pt}#1}\par\nopagebreak\vspace{4pt}
}

% ============================================================
% Boîtes pédagogiques — même recette : fond blanc, bordure épaisse colorée,
% ombre dure encre, badge plein. Titre optionnel [#1].
% ============================================================
\tcbset{popbase/.style={breakable,enhanced,colback=white,boxrule=2.3pt,arc=9pt,
  shadow={3pt}{-3pt}{0pt}{popink},left=14pt,right=14pt,top=11pt,bottom=11pt,before skip=11pt,after skip=15pt,
  fontupper=\popsemi\fontsize{10.6}{14.5}\selectfont\color{popink},
  fontlower=\popsemi\fontsize{10.6}{14.5}\selectfont\color{popink}}}
% Coche dessinée (le glyphe ✓ n'existe pas dans les polices du document)
\newcommand{\popcheck}[1]{\tikz[baseline=-0.5ex]\draw[#1,line width=1.6pt,line cap=round,line join=round](-0.13,0.0)--(-0.04,-0.09)--(0.13,0.1);}
\providecommand{\coche}{\popcheck{popgreen}}
\providecommand{\resultat}[1]{\par\smallskip\noindent\fcolorbox{popgreen}{popgreenL}{\parbox{\dimexpr\linewidth-2\fboxsep-2\fboxrule\relax}{\textbf{Résultat.} #1}}\par\smallskip}
\newcommand{\popboxhead}[3]{\poppill{#1}{#2}{white}\ifx\relax#3\relax\else\hspace{6pt}{\popxbold\fontsize{10.6}{12}\selectfont\color{popink}#3}\fi\par\vspace{7pt}}

\newtcolorbox{aretenir}[1][]{popbase,colframe=popgreen,before upper={\popboxhead{À retenir}{popgreen}{#1}}}
\newtcolorbox{exemplebox}[1][]{popbase,colframe=poporange,before upper={\popboxhead{Exemple}{poporange}{#1}}}
\newtcolorbox{definition}[1][]{popbase,colframe=popblue,before upper={\popboxhead{Définition}{popblue}{#1}}}
\newtcolorbox{propriete}[1][]{popbase,colframe=poppurple,before upper={\popboxhead{Propriété}{poppurple}{#1}}}
\newtcolorbox{methode}[1][]{popbase,colframe=popgold,before upper={\popboxhead{Méthode}{popgold}{#1}}}
\newtcolorbox{attention}[1][]{popbase,colframe=poppink,before upper={\popboxhead{Attention}{poppink}{#1}}}
\newtcolorbox{experience}[1][]{popbase,colframe=popblue,colback=popblueL,before upper={\popboxhead{Expérience}{popblue}{#1}}}
% « Le savais-tu ? » : carte violette pleine (contexte, culture, Cameroun)
\newtcolorbox{savaistu}[1][]{popbase,colback=poppurple,colframe=popink,
  fontupper=\popsemi\fontsize{10.4}{14.2}\selectfont\color{white},
  before upper={\poppill{Le savais-tu\,?}{white}{poppurple}\ifx\relax#1\relax\else\hspace{6pt}{\popxbold\color{white}#1}\fi\par\vspace{7pt}}}
% Exercice résolu : énoncé en haut, \tcblower puis la solution sur fond crème
\newcounter{popexo}\newcounter{popexr}
\newtcolorbox{exoresolu}[1][]{popbase,colframe=poporange,colbacklower=popcream,
  segmentation style={popink,line width=1.2pt,dashed},
  before upper={\stepcounter{popexr}\popboxhead{Exercice résolu \thepopexr}{poporange}{#1}},
  before lower={\poppill{Solution}{popink}{popcream}\par\vspace{6pt}}}
% Exercice (non résolu en place) — la correction est dans la partie Corrigés
\newtcolorbox{exercice}[1][]{popbase,colframe=popink,
  before upper={\stepcounter{popexo}\popboxhead{Exercice \thepopexo}{popink}{#1}}}
% Corrigé
\newtcolorbox{corrige}[1][]{popbase,colframe=popgreen,colback=popgreenL,
  before upper={\popboxhead{Corrigé}{popgreen}{#1}}}
% Objectifs / prérequis / bilan
\newtcolorbox{objectifs}{popbase,colframe=popink,colback=poporangeL,
  before upper={\popboxhead{Objectifs de la leçon}{popink}{}}}
\newtcolorbox{prerequis}{popbase,colframe=popink,colback=popblueL,
  before upper={\popboxhead{Prérequis}{popblue}{}}}
\newtcolorbox{bilan}{popbase,colframe=popink,colback=white,boxrule=2.8pt,
  before upper={\popboxhead{Fiche bilan}{poporange}{L'essentiel à connaître pour l'examen}}}
% Figure encadrée avec légende
\newtcolorbox{popfigure}[1][]{enhanced,breakable=false,colback=white,colframe=popline,boxrule=1.4pt,arc=8pt,
  left=10pt,right=10pt,top=10pt,bottom=8pt,before skip=10pt,after skip=12pt,halign=center,
  lower separated=false,halign lower=center,
  fontlower=\popsemi\fontsize{9}{11.5}\selectfont\color{popmuted},
  #1}
\newcommand{\legende}[1]{\par\vspace{6pt}{\popsemi\fontsize{9}{11.5}\selectfont\color{popmuted}#1}}

% ============================================================
% Signature OnBuch+
% ============================================================
\newcommand{\poplogoinline}[1]{{\poplogo\fontsize{#1}{#1}\selectfont\color{popink}OnBuch\color{poporange}+}}
\newcommand{\signatureonbuch}{%
  \begin{tcolorbox}[enhanced,colback=popink,colframe=popink,boxrule=2.2pt,arc=10pt,
      drop shadow=poporange,left=14pt,right=14pt,top=10pt,bottom=10pt,before skip=0pt,after skip=0pt]
    \tikz[baseline=-0.7ex]\node[circle,fill=popgreen,draw=popcream,line width=1.3pt,minimum size=22pt,inner sep=0]{\popcheck{white}};%
    \hspace{9pt}\begin{minipage}[c]{0.62\linewidth}
      {\popxbold\fontsize{12}{13}\selectfont\color{popcream}Signé\ \textbullet\ {\poplogo OnBuch\color{poporange}+}}\par\vspace{2pt}
      {\raggedright\popsemi\fontsize{8}{10.5}\selectfont\color{popline}\MakeUppercase{Équipe pédagogique OnBuch+}\\\MakeUppercase{Conforme au programme officiel MINESEC}\par}
    \end{minipage}\hfill
    {\poplogo\fontsize{22}{22}\selectfont\color{popcream}OnBuch\color{poporange}+}
  \end{tcolorbox}%
}

% ============================================================
% Page de garde
% #1 titre · #2 matière · #3 niveau · #4 module · #5 n° leçon · #6 durée ·
% #7 description / objectifs (texte ou itemize)
% ============================================================
\newcommand{\pagegarde}[7]{%
  \thispagestyle{empty}
  \newgeometry{a4paper,margin=1.7cm,top=1.6cm,bottom=1.4cm}%
  \begin{tikzpicture}[remember picture,overlay]
    \fill[poporange!18] ([xshift=-3.2cm,yshift=-3.6cm]current page.north east) circle (4.6cm);
    \fill[poppurple!14] ([xshift=2.4cm,yshift=7.6cm]current page.south west) circle (3.8cm);
  \end{tikzpicture}%
  {\poplogo\fontsize{30}{30}\selectfont\color{popink}OnBuch\color{poporange}+}\hfill
  \raisebox{0.6ex}{\poppill{Cours signé \textbullet\ Premium}{popink}{popcream}}\par
  \vspace{1pt}\popsquiggle{poppurple}\par
  \vspace{8pt}
  \begin{tcolorbox}[enhanced,colback=poporange,colframe=popink,boxrule=2.8pt,arc=13pt,
      drop shadow=popink,left=18pt,right=18pt,top=12pt,bottom=13pt,after skip=10pt]
    \poppill{#2 \textbullet\ #3}{popink}{popcream}\hspace{5pt}\poppill{#4}{white}{popink}\par\vspace{8pt}
    {\popdisplay\fontsize{14}{14}\selectfont\color{popink}LEÇON #5}\par\vspace{4pt}
    {\raggedright\hyphenpenalty=10000\exhyphenpenalty=10000\popdisplay\fontsize{25}{27}\selectfont\color{popcream}\MakeUppercase{#1}\par}
    \vspace{8pt}
    {\popxbold\fontsize{9.5}{9.5}\selectfont\color{popink}\MakeUppercase{Durée conseillée : #6}}
  \end{tcolorbox}
  \begin{tcolorbox}[enhanced,colback=white,colframe=popink,boxrule=2.2pt,arc=10pt,
      drop shadow=popink,left=16pt,right=16pt,top=10pt,bottom=10pt,after skip=8pt]
    \poppill{Dans cette leçon}{poppurple}{white}\par\vspace{5pt}
    \popsemi\fontsize{9.3}{12.6}\selectfont\color{popink}#7
  \end{tcolorbox}
  \noindent\begin{minipage}[t]{0.487\linewidth}
    \begin{tcolorbox}[enhanced,colback=popgreen,colframe=popink,boxrule=2.2pt,arc=10pt,
        drop shadow=popink,left=11pt,right=11pt,top=10pt,bottom=10pt,height=3.1cm,valign=center]
      \tcbox[colback=white,colframe=popink,boxrule=1.3pt,arc=5pt,boxsep=0pt,nobeforeafter,
        left=3pt,right=3pt,top=3pt,bottom=3pt,valign=center]{\qrcode[height=1.9cm]{https://play.google.com/store/apps/details?id=com.luvvix.onbuch&hl=fr}}%
      \hspace{8pt}\begin{minipage}[c]{0.5\linewidth}
        {\popdisplay\fontsize{11}{12.5}\selectfont\color{white}\MakeUppercase{Télécharge OnBuch}}\par\vspace{4pt}
        {\raggedright\popsemi\fontsize{8.4}{11}\selectfont\color{white}Gratuit \textbullet\ Google Play\\Tuteur IA, annales, concours, résultats\par}
      \end{minipage}
    \end{tcolorbox}
  \end{minipage}\hfill
  \begin{minipage}[t]{0.487\linewidth}
    \begin{tcolorbox}[enhanced,colback=poppurple,colframe=popink,boxrule=2.2pt,arc=10pt,
        drop shadow=popink,left=11pt,right=11pt,top=10pt,bottom=10pt,height=3.1cm,valign=center]
      \tikz[baseline=-0.7ex]\node[circle,fill=white,draw=popink,line width=1.3pt,minimum size=1.6cm,inner sep=0]{\popdisplay\fontsize{24}{24}\selectfont\color{poppurple}+};%
      \hspace{8pt}\begin{minipage}[c]{0.6\linewidth}
        {\popdisplay\fontsize{11}{12.5}\selectfont\color{white}\MakeUppercase{Passe à OnBuch+}}\par\vspace{4pt}
        {\raggedright\popsemi\fontsize{8.4}{11}\selectfont\color{white}Tous les cours signés, Tuteur IA illimité, fiches PDF — onglet OnBuch+ de l'app\par}
      \end{minipage}
    \end{tcolorbox}
  \end{minipage}
  \vspace{8pt}
  \popcontenu
  \vfill
  \signatureonbuch
  \restoregeometry
  \newpage
}

% Rangée « Ce cours contient » (page de garde)
\newcommand{\poptile}[3]{% #1 couleur #2 gros texte #3 légende
  \begin{tcolorbox}[enhanced,colback=#1,colframe=popink,boxrule=2pt,arc=9pt,shadow={2.5pt}{-2.5pt}{0pt}{popink},
      left=6pt,right=6pt,top=9pt,bottom=9pt,halign=center,valign=center,height=1.75cm,width=\linewidth,nobeforeafter]
    {\popdisplay\fontsize{12.5}{13}\selectfont\color{white}\MakeUppercase{#2}}\par\vspace{3pt}
    {\centering\popsemi\fontsize{7.8}{9.5}\selectfont\color{white}#3\par}
  \end{tcolorbox}}
\newcommand{\popcontenu}{%
  \par\noindent{\popxboldls\fontsize{7.5}{7.5}\selectfont\color{popmuted}CE COURS SIGNÉ CONTIENT}\par\vspace{7pt}
  \noindent\begin{minipage}[t]{0.235\linewidth}\poptile{popblue}{Cours}{complet, illustré, pas à pas}\end{minipage}\hfill
  \begin{minipage}[t]{0.235\linewidth}\poptile{poporange}{Méthodes}{et exercices résolus}\end{minipage}\hfill
  \begin{minipage}[t]{0.235\linewidth}\poptile{poppink}{Exercices}{type examen + corrigés}\end{minipage}\hfill
  \begin{minipage}[t]{0.235\linewidth}\poptile{popgreen}{Fiche bilan}{l'essentiel à retenir}\end{minipage}\par}

% Sommaire
\newtcolorbox{sommaire}{enhanced,colback=white,colframe=popink,boxrule=2.2pt,arc=10pt,
  drop shadow=popink,left=18pt,right=18pt,top=13pt,bottom=13pt,before skip=6pt,after skip=20pt,
  before upper={\poppill{Sommaire}{poppurple}{white}\par\vspace{9pt}\popsemi\fontsize{10.6}{14.5}\selectfont\color{popink}}}

% Signature de fin de cours — poussée tout en bas de la dernière page
\newcommand{\findecours}{%
  \par\vfill\nopagebreak
  \begin{center}\popsquiggle{poporange}\end{center}\vspace{4pt}
  \signatureonbuch
}
\AtBeginDocument{\def\llbracket{\mathopen{[\mkern-3.5mu[}}\def\rrbracket{\mathclose{]\mkern-3.5mu]}}} % intervalles d'entiers (glyphe ⟦ absent de la police)
% Glyphes absents de Fira Math : redéfinis à partir de glyphes disponibles
\AtBeginDocument{%
  \def\setminus{\mathbin{\backslash}}\let\smallsetminus\setminus
  \def\vdots{\mathord{\vbox{\baselineskip3.2pt\lineskiplimit0pt\kern4pt\hbox{.}\hbox{.}\hbox{.}}}}%
  \def\ddots{\mathinner{\mkern1mu\raise6.5pt\hbox{.}\mkern2mu\raise3.6pt\hbox{.}\mkern2mu\raise0.7pt\hbox{.}\mkern1mu}}%
  \def\triangle{\mathord{\scalebox{0.9}{$\symup{\Delta}$}}}%
  \def\nabla{\mathord{\raisebox{\depth}{\scalebox{1}[-1]{$\symup{\Delta}$}}}}%
  \def\checkmark{\popcheck{popdark}}%
  \def\star{\mathbin{\ast}}\def\bigstar{\mathord{\ast}}%
  \def\diamond{\mathbin{\scalebox{0.6}{$\Diamond$}}}%
  \def\bigcup{\mathop{\mathchoice{\raisebox{-0.3em}{\scalebox{1.7}{$\cup$}}}{\scalebox{1.25}{$\cup$}}{\cup}{\cup}}\displaylimits}%
  \def\bigcap{\mathop{\mathchoice{\raisebox{-0.3em}{\scalebox{1.7}{$\cap$}}}{\scalebox{1.25}{$\cap$}}{\cap}{\cap}}\displaylimits}%
  \def\lesseqqgtr{\mathrel{\lessgtr}}\def\bigoplus{\mathop{\oplus}\displaylimits}%
}
\providecommand{\ssi}{si et seulement si} % abréviation fréquente
\AtBeginDocument{\providecommand{\wideparen}{\overparen}} % arc de cercle

%% FIGFIX-BEGIN v3 — correctifs globaux des figures (docs/onbuch-generation/tools/figfix)
\usepackage{adjustbox}\usepackage{etoolbox}
% toute figure TikZ/pgfplots plus large que la ligne est réduite à la largeur disponible
\BeforeBeginEnvironment{tikzpicture}{\begin{adjustbox}{max width=\linewidth}}
\AfterEndEnvironment{tikzpicture}{\end{adjustbox}}
% légendes pgfplots lisibles par-dessus les courbes
\pgfplotsset{every axis/.append style={legend style={fill=white,fill opacity=0.92,text opacity=1,draw=black!15,inner sep=3pt}}}
% étiquettes d'arêtes (syntaxe edge["..."]) avec fond blanc
\tikzset{every edge quotes/.append style={fill=white,inner sep=1.2pt}}
% bulles de cartes mentales : texte à la ligne dans la bulle, bulle un peu plus grande
\tikzset{every concept/.append style={text width=2.0cm,align=center,minimum size=2.3cm,inner sep=2pt}}
%% FIGFIX-END
\begin{document}

\pagegarde{Produire une lettre officielle : la lettre de motivation}{Français}{1re Tech}{Français – classe de Première technique}{N12}{21 séances (fiche de progression harmonisée)}{Cette leçon outille l'élève de Première technique pour lire, analyser et produire une lettre de motivation conforme aux normes de la correspondance officielle. De l'observation de modèles authentiques à la rédaction complète (introduction, corps, conclusion), elle prépare directement à l'épreuve de production de texte du Probatoire et du Baccalauréat technique.
\vspace{4pt}
\begin{multicols}{2}\small
\begin{itemize}
\item À la fin de cette leçon, je sais identifier la structure externe et interne d'une lettre de motivation à partir de modèles authentiques.
\item À la fin de cette leçon, je sais distinguer lettre de motivation, lettre administrative et curriculum vitae.
\item À la fin de cette leçon, je sais rédiger l'en-tête, l'objet et les références d'une lettre officielle selon les normes.
\item À la fin de cette leçon, je sais rédiger une introduction accrocheuse et une conclusion polie et incitative.
\item À la fin de cette leçon, je sais construire le corps de la lettre en articulant formation, expérience, qualités et motivation.
\item À la fin de cette leçon, je sais employer les figures de style (gradation, antithèse, métonymie) pour rendre mon propos persuasif.
\end{itemize}\end{multicols}}

\begin{sommaire}
\begin{multicols}{2}
\begin{enumerate}
\item Découverte : lire des lettres de motivation authentiques (textes iconiques 1 à 5)
\item La structure externe de la lettre de motivation
\item La structure interne : introduction et conclusion
\item Le corps de la lettre : convaincre avec style
\item Production guidée : rédiger une lettre de motivation complète
\item Compte-rendu, évaluation et remédiation
\item Activité d'intégration
\item Méthodes et exercices résolus
\item Exercices
\item Corrigés des exercices
\item Fiche bilan
\end{enumerate}
\end{multicols}
\end{sommaire}

\begin{objectifs}
\begin{itemize}
\item À la fin de cette leçon, je sais identifier la structure externe et interne d'une lettre de motivation à partir de modèles authentiques.
\item À la fin de cette leçon, je sais distinguer lettre de motivation, lettre administrative et curriculum vitae.
\item À la fin de cette leçon, je sais rédiger l'en-tête, l'objet et les références d'une lettre officielle selon les normes.
\item À la fin de cette leçon, je sais rédiger une introduction accrocheuse et une conclusion polie et incitative.
\item À la fin de cette leçon, je sais construire le corps de la lettre en articulant formation, expérience, qualités et motivation.
\item À la fin de cette leçon, je sais employer les figures de style (gradation, antithèse, métonymie) pour rendre mon propos persuasif.
\item À la fin de cette leçon, je sais analyser une image ou une affiche (communication par l'image, affiches sécuritaires) et l'exploiter dans un écrit.
\item À la fin de cette leçon, je sais produire, réviser et améliorer une lettre de motivation complète à partir d'une situation donnée.
\item À la fin de cette leçon, je sais justifier mes choix de rédaction et appliquer ces compétences à des situations concrètes de la vie courante.
\end{itemize}
\end{objectifs}

\begin{prerequis}
\begin{itemize}
\item La correspondance administrative vue en classe de Seconde (lettre de demande, requête).
\item Les types de textes (narratif, descriptif, argumentatif) vus en début d'année.
\item Les connecteurs logiques et l'organisation du paragraphe argumentatif.
\item Les formules de politesse et les registres de langue (familier, courant, soutenu).
\item La relecture et la correction orthographique (accords, homophones).
\end{itemize}
\end{prerequis}

\newpage

\coursec{Découverte : lire des lettres de motivation authentiques (textes iconiques 1 à 5)}

\courssub{Situation d'accroche et problématique}

Imaginez Aminatou, jeune diplômée d'un lycée technique de Douala, spécialité comptabilité. Un matin, elle feuillette \emph{Cameroon Tribune} et repère une offre : une PME de Bonabéri cherche une assistante comptable. Pressée, elle rédige sa lettre en cinq minutes sur une feuille arrachée : pas d'en-tête, pas de référence à l'annonce, des fautes d'orthographe, une formule de politesse familière (« \emph{Salut, je veux le job} »). Résultat : sa candidature finit à la corbeille sans être lue.

Son amie Clarisse, moins brillante sur les bulletins, voit la même annonce. Elle prend le temps : elle soigne la présentation, adresse sa lettre au directeur par son nom, cite la référence de l'annonce, met en avant son stage à la SABC, sa maîtrise de Sage Compta, sa rigueur. Elle joint son CV. Trois jours plus tard, le téléphone sonne : Clarisse est convoquée en entretien.

\medskip
\textbf{Problématique :} Qu'est-ce qui fait la différence entre une lettre qui ouvre des portes et une lettre qui finit à la corbeille ? Comment produire une lettre de motivation capable de « vendre » un profil auprès d'un employeur camerounais ?

\courssub{Les cinq textes iconiques : observation guidée}

Nous allons analyser cinq modèles authentiques (textes iconiques n°1 à 5) correspondant à des situations réelles du marché de l'emploi et de la formation au Cameroun. Pour chacun, nous appliquons la méthode d'analyse des textes iconiques.

\begin{methode}[Analyser un texte iconique (lettre de motivation)]
\textbf{Étape 1 -- Décrire :} Que voit-on ? (Mise en page, en-tête, police, paragraphes, signature, pièces jointes mentionnées.)
\textbf{Étape 2 -- Repérer la structure :} Où sont les éléments obligatoires ? (Expéditeur, destinataire, date, objet, formules d'appel et de politesse, corps structuré en trois parties : introduction, développement, conclusion.)
\textbf{Étape 3 -- Interpréter :} Quel effet sur le lecteur ? (Ton respectueux, arguments concrets, adaptation au poste, clarté, incitation à lire le CV et à convoquer.)
\end{methode}

\begin{exemplebox}[Texte iconique n°1 -- Candidature à un stage à la SABC (Douala)]
\textbf{Contexte :} Élève de Première G2 (Gestion), stage de vacances de deux mois.
\textbf{Extrait commenté :}
\begin{verbatim}
Douala, le 15 juin 2025

À l'attention de Monsieur le Directeur des Ressources Humaines
Société Anonyme des Brasseries du Cameroun (SABC)
B.P. 4070, Douala

Objet : Demande de stage de vacances (juillet--août 2025) -- Réf. : Annonce n° 2025/SABC/STG/045

Monsieur le Directeur,

Actuellement élève en classe de Première G2 au Lycée Technique de New-Bell, je souhaite effectuer un stage de vacances au sein de votre direction comptable du 1er juillet au 31 août 2025.

Votre entreprise, leader du secteur agroalimentaire au Cameroun, représente pour moi un cadre d'apprentissage idéal. Lors de mon stage d'initiation de l'an dernier à la Trésorerie Générale de Douala, j'ai pu manipuler le logiciel Sage Compta et participer à la saisie des écritures journalières. Rigoureuse, ponctuelle et dotée d'un bon esprit d'équipe, je suis prête à mettre mon enthousiasme et mes connaissances au service de votre service comptabilité.

Je vous prie d'agréer, Monsieur le Directeur, l'expression de ma considération distinguée.

[Signature]
Aminatou NGONO
Tél. : 6XX XX XX XX -- Email : aminatou.ngono@email.cm

Pièce jointe : Curriculum Vitae
\end{verbatim}
\textbf{Analyse :} En-tête complet (expéditeur, destinataire, date, objet précis avec référence). Introduction : qui je suis + quoi + quand. Développement : pourquoi la SABC + preuve de compétence (stage antérieur, logiciel) + qualités personnelles. Conclusion : formule administrative standard + signature + coordonnées + mention CV. Ton formel, respectueux, confiant sans arrogance.
\end{exemplebox}

\begin{exemplebox}[Texte iconique n°2 -- Candidature à un emploi à la Mairie de Bafoussam]
\textbf{Contexte :} Diplômé BTS Assistant de Direction, réponse à un avis de recrutement pour un poste de secrétaire de mairie.
\textbf{Points saillants :} Objet mentionnant le numéro de l'avis (N° 012/Mairie/BFS/2025). Introduction rappelant le diplôme et la spécialité. Développement : connaissance de l'environnement communal (stages en mairie d'arrondissement), maîtrise des outils bureautiques (Word, Excel, Outlook), sens du service public, discrétion. Conclusion avec disponibilité pour entretien. Mention « Pièces jointes : CV, copies des diplômes, extrait d'acte de naissance, certificat de nationalité ».
\end{exemplebox}

\begin{exemplebox}[Texte iconique n°3 -- Candidature à une formation à l'ENAM (École Nationale d'Administration et de Magistrature)]
\textbf{Contexte :} Fonctionnaire catégorie C (secrétaire administratif) candidat au concours professionnel cycle A (attaché d'administration).
\textbf{Spécificités :} Destinataire : « Monsieur le Directeur Général de l'ENAM ». Objet : « Demande d'inscription au concours professionnel cycle A, session 2025 -- Réf. : Arrêté n° 00012/MINFOPRA/ENAM ». Corps : rappel de la situation administrative actuelle (grade, ancienneté, affectation), motivation pour l'évolution de carrière, projets pour l'administration. Ton très administratif, respectueux de la hiérarchie.
\end{exemplebox}

\begin{exemplebox}[Texte iconique n°4 -- Candidature à un recrutement à Camtel (Yaoundé)]
\textbf{Contexte :} Jeune ingénieur télécoms (DUT Génie Électrique), réponse à une campagne de recrutement « Jeunes Talents Camtel 2025 ».
\textbf{Originalité :} Lettre plus « moderne » : objet par email « Candidature Jeunes Talents -- Ingénieur Réseaux -- Jean KAMGA ». Corps structuré en \textbf{trois arguments numérotés} : 1) Compétences techniques (configuration routeurs Cisco, fibre optique, projet fin d'études sur la 5G), 2) Expérience terrain (stage de 6 mois à MTN Cameroon, maintenance sites relais), 3) Adhésion aux valeurs Camtel (service public, innovation, intégrité). Conclusion proactive : « Je me tiens à votre disposition pour un entretien technique dès la semaine prochaine. »
\end{exemplebox}

\begin{exemplebox}[Texte iconique n°5 -- Candidature à un concours professionnel (Ministère de la Santé)]
\textbf{Contexte :} Infirmier diplômé d'État (IDE) avec 5 ans d'ancienneté, candidat au concours d'accès au grade d'Infirmier Principal.
\textbf{Structure :} Très codifiée. Référence à l'arrêté d'ouverture du concours. Rappel précis de la situation administrative (matricule, date de recrutement, dernier poste). Motivation : volonté d'assumer des responsabilités d'encadrement, projet d'amélioration de la prise en charge dans son district de santé. Formule de politesse administrative longue (« Je vous prie d'agréer, Monsieur le Ministre, l'assurance de ma très haute considération. »).
\end{exemplebox}

\courssub{Repérage des constantes : ce que toutes les lettres partagent}

Après l'observation des cinq modèles, cinq constantes émergent. Elles constituent le « squelette » de toute lettre de motivation réussie au Cameroun.

\begin{aretenir}[Les cinq constantes de la lettre de motivation]
\begin{enumerate}
\item \textbf{Destinataire identifié avec précision :} Nom, fonction, structure, adresse postale (ou email). Pas de « Monsieur le Directeur » vague si le nom est connu.
\item \textbf{Objet précis et référencé :} Nature de la demande (stage, emploi, formation, concours) + référence à l'annonce/arrêté/avis (numéro, date, source).
\item \textbf{Ton respectueux et formel :} Vouvoiement, formules d'appel et de politesse administratives adaptées au rang du destinataire.
\item \textbf{Mise en valeur du candidat (le « moi » argumenté) :} Diplômes, expériences, compétences techniques, qualités humaines, \emph{en lien direct avec le poste visé}.
\item \textbf{Appel à l'action (entretien) + mentions obligatoires :} Disponibilité, coordonnées, signature, liste des pièces jointes (CV en premier).
\end{enumerate}
\end{aretenir}

\courssub{Distinction cruciale : lettre de motivation, lettre administrative, demande d'emploi, CV}

Beaucoup d'élèves (et même d'adultes) confondent ces écrits. Le tableau comparatif ci-dessous fixe les frontières.

\begin{popfigure}
\begin{center}
\begin{tabularx}{\linewidth}{|l|X|X|X|}
\hline
\rowcolor{popblueL}
\textbf{Critère} & \textbf{Lettre de motivation} & \textbf{Lettre administrative} & \textbf{Curriculum Vitae (CV)} \\
\hline
Destinataire & Une personne identifiée (DRH, Directeur, Ministre) & Une administration / service (souvent impersonnel) & Le même destinataire que la lettre (accompagne la lettre) \\
\hline
But principal & \textbf{Convaincre} : prouver l'adéquation profil/poste & \textbf{Informer / Demander} : exposer une situation, solliciter un droit & \textbf{Informer / Énumérer} : présenter le parcours factuel (diplômes, stages, compétences) \\
\hline
Ton / Style & Persuasif, personnel, argumenté, « je » assumé & Impersonnel, objectif, formel, norme administrative & Neutre, factuel, télégraphique (pas de phrases complètes), « je » absent \\
\hline
Structure & 3 parties : intro (accroche), corps (arguments), conclusion (appel entretien) & En-tête, objet, exposé des faits, conclusion administrative, visa/signature & Rubriques : État civil, Formation, Expériences, Compétences, Langues, Centres d'intérêt \\
\hline
Pièces jointes & CV + copies diplômes + attestations + extraits acte naissance... & Selon la demande (pièces justificatives) & Aucune (c'est lui-même une pièce jointe) \\
\hline
Longueur & 1 page maximum (200--300 mots) & Variable (souvent 1 page) & 1 à 2 pages maximum \\
\hline
\end{tabularx}
\end{center}
\legende{Tableau comparatif : lettre de motivation vs lettre administrative vs CV. La lettre de motivation est le seul écrit \textbf{persuasif} ; le CV est \textbf{factuel} ; la lettre administrative est \textbf{procédurale}.}
\end{popfigure}

\begin{attention}[Piège fréquent : confondre lettre de motivation et CV]
\textbf{Erreur :} Recopier son CV en phrases dans la lettre (« J'ai eu mon BAC en 2022, puis j'ai fait un BTS en 2024, puis un stage ici... »).
\textbf{Correction :} Le CV \textbf{énumère} (liste chronologique, faits bruts). La lettre \textbf{argumente} (sélectionne 2--3 expériences clés, explique \emph{quoi} on a appris, \emph{comment} cela sert le poste visé). La lettre donne du sens au CV.
\end{attention}

\courssub{Analyse des affiches d'indications sécuritaires (textes iconiques) : deux écrits utilitaires à visée incitative}

Le programme prévoit l'étude d'affiches de sécurité (texte iconique n°6). Bien que différentes par le support et le public, l'affiche de sécurité et la lettre de motivation partagent une nature commune : ce sont des \textbf{écrits utilitaires à visée incitative}.

\begin{center}
\begin{tabularx}{\linewidth}{|l|X|X|}
\hline
\rowcolor{popblueL}
\textbf{Caractéristique} & \textbf{Affiche de sécurité} & \textbf{Lettre de motivation} \\
\hline
\cle{Émetteur} & Employeur / Service HSE / Administration & Candidat (personne physique) \\
\cle{Destinataire} & Collectivité large (employés, public, usagers) & Personne unique identifiée (décideur) \\
\cle{Visée} & \textbf{Inciter à un comportement} (porter casque, ne pas fumer, issue de secours) & \textbf{Inciter à une décision} (lire le CV, convoquer en entretien) \\
\cle{Structure} & \textbf{Consigne} (impératif) + \textbf{Image} (pictogramme) + \textbf{Slogan} (mémorisable) & \textbf{Introduction} (accroche) + \textbf{Corps} (arguments) + \textbf{Conclusion} (appel à l'action) \\
\cle{Registre de langue} & Impératif, injonctif, visuel & Argumentatif, persuasif, textuel \\
\cle{Exemple} & \emph{« Port du casque obligatoire »} + pictogramme casque + \emph{« Ta vie, ton casque »} & \emph{« Fort de 3 ans en maintenance industrielle, je souhaite... »} + \emph{« Disponible pour entretien »} \\
\hline
\end{tabularx}
\end{center}

\begin{savaistu}[Le pouvoir de l'écrit utilitaire au Cameroun]
Dans les entreprises camerounaises (SABC, Camtel, CDC, ENEO, mairies), l'affichage sécuritaire est une obligation légale (Code du travail, décret n° 92/074 du 13 avril 1992). Une affiche mal conçue (texte trop long, image absente, slogan incompréhensible) ne protège personne. De même, une lettre de motivation mal structurée ne « protège » pas le candidat : elle ne déclenche pas l'entretien. Dans les deux cas, \textbf{la forme conditionne l'efficacité}.
\end{savaistu}

\courssub{Définition et fonctions de la lettre de motivation}

\begin{definition}[Lettre de motivation]
La \cle{lettre de motivation} est une \textbf{lettre officielle} par laquelle un candidat (demandeur d'emploi, de stage, de formation ou de promotion) \textbf{sollicite} un poste, un stage ou une place en formation en \textbf{mettant en valeur son profil et sa motivation} afin de \textbf{convaincre le destinataire} de l'inviter à un entretien.
\end{definition}

\begin{propriete}[Trois fonctions indissociables]
\begin{enumerate}
\item \textbf{Fonction informative :} Informer le destinataire de l'identité du candidat, de l'objet précis de la demande, de la source de l'information (annonce, arrêté, connaissance).
\item \textbf{Fonction persuasive (convaincre) :} Démontrer l'adéquation entre le profil du candidat (compétences, expériences, qualités) et les exigences du poste/formation. C'est le cœur de la lettre : \emph{argumentation}.
\item \textbf{Fonction incitative (déclencher l'entretien) :} Pousser le recruteur à l'action : lire le CV joint, puis téléphoner pour fixer un rendez-vous. La conclusion doit être un \cle{call to action} clair.
\end{enumerate}
\end{propriete}

\courssub{Vocabulaire spécifique du recrutement}

Maîtriser ce lexique, c'est parler le langage des recruteurs camerounais.

\begin{center}
\begin{tabularx}{\linewidth}{|l|X|}
\hline
\rowcolor{popblueL}
\textbf{Terme} & \textbf{Définition / Contexte camerounais} \\
\hline
\cle{Candidature} & Ensemble des documents envoyés : lettre de motivation + CV + pièces justificatives. \\
\cle{Poste / Emploi} & Fonction précise (ex. : « Assistant comptable », « Technicien de maintenance »). \\
\cle{Recruteur / DRH} & Personne (Directeur des Ressources Humaines, Chef de service, Maire, Ministre) qui décide. \\
\cle{Pièces jointes / Pièces du dossier} & CV, copies certifiées conformes des diplômes, attestations de stage/travail, extrait d'acte de naissance, certificat de nationalité, casier judiciaire (bulletin n°3), photos d'identité. \\
\cle{Entretien} & Échange oral (souvent en panel au Cameroun : DRH + chef technique + DG) pour valider les compétences et la motivation. \\
\cle{Présélection / Shortlist} & Première sélection sur dossier (lettre + CV) avant l'entretien. \\
\cle{Référence / Réf. de l'annonce} & Numéro ou date de l'avis (ex. : « Ref. 2025/SABC/STG/045 » ou « Avis n° 012/MINFOPRA »). Indispensable dans l'objet. \\
\cle{Formule de politesse / Formule d'appel} & Codes rituels : « Monsieur le Directeur, » (appel) ; « Je vous prie d'agréer, Monsieur le Directeur, l'expression de ma considération distinguée. » (politesse). \\
\hline
\end{tabularx}
\end{center}

\courssub{Schéma : le parcours d'une candidature (de l'offre à l'embauche)}

Comprendre où s'insère la lettre permet de mesurer son poids stratégique.

\begin{popfigure}
\begin{center}
\begin{tikzpicture}[
  node distance=1.2cm and 1.5cm,
  box/.style={rectangle, draw=popblue, fill=popblueL, rounded corners, minimum width=2.8cm, minimum height=1.2cm, align=center, font=\small},
  arrow/.style={->, >=Stealth, thick, popblue},
  decision/.style={diamond, draw=poporange, fill=poporangeL, aspect=2, minimum width=3cm, minimum height=1.2cm, align=center, font=\small},
  start/.style={ellipse, draw=popgreen, fill=popgreenL, minimum width=2.5cm, minimum height=1cm, align=center, font=\small},
]
% Nœuds
\node[start, align=center] (offre){Offre d'emploi /\\ Avis de concours /\\ Besoin de stage};
\node[box, right=of offre, align=center] (prepa){Préparation :\\ Analyse de l'offre\\ + Connaissance de soi\\ + Recherche structure};
\node[box, right=of prepa, align=center] (redaction){Rédaction :\\ Lettre de motivation\\ + CV actualisé\\ + Pièces jointes};
\node[box, right=of redaction, align=center] (envoi){Envoi / Dépôt :\\ Email / Physique /\\ Plateforme en ligne};
\node[decision, right=of envoi, align=center] (preselection){Présélection\\ sur dossier ?};
\node[box, right=of preselection, align=center] (entretien){Entretien :\\ Tests techniques\\ + Évaluation soft skills\\ + Négociation};
\node[decision, right=of entretien] (decision) {Décision finale ?};
\node[box, right=of decision, fill=popgreenL, draw=popgreen, align=center] (embauche){Embauche /\\ Admission /\\ Signature contrat};
\node[box, below=of decision, fill=poppinkL, draw=poppink, align=center] (echec){Échec :\\ Demander feedback\\ + Rebondir};
% Flèches
\draw[arrow] (offre) -- (prepa) node[midway, above, font=\tiny] {1. Repérer};
\draw[arrow] (prepa) -- (redaction) node[midway, above, font=\tiny] {2. Construire};
\draw[arrow] (redaction) -- (envoi) node[midway, above, font=\tiny] {3. Envoyer};
\draw[arrow] (envoi) -- (preselection) node[midway, above, font=\tiny] {4. Attendre};
\draw[arrow] (preselection) -- node[midway, above, font=\tiny] {OUI} (entretien);
\draw[arrow] (preselection) -- node[midway, right, font=\tiny] {NON} ++(0,-1.5) -| (echec);
\draw[arrow] (entretien) -- (decision) node[midway, above, font=\tiny] {5. Convaincre};
\draw[arrow] (decision) -- node[midway, above, font=\tiny] {OUI} (embauche);
\draw[arrow] (decision) -- node[midway, right, font=\tiny] {NON} (echec);
% Annotation importante
\node[above=0.3cm of redaction, font=\small\bfseries, text=popdark] (crucial) {\cle{La lettre + le CV = Clé de l'étape 4}};
\draw[popdark, dashed, thick] (crucial.south) -- (redaction.north);
\end{tikzpicture}
\end{center}
\legende{Parcours d'une candidature : la lettre de motivation (étape 3) est le \textbf{seul document persuasif} lu par le recruteur avant l'entretien. Si elle échoue, le CV, aussi bon soit-il, ne sera pas ouvert.}
\end{popfigure}

\courssub{Exercice résolu : analyse comparative de deux lettres}

\begin{exoresolu}[Deux lettres pour le même poste : laquelle passe ?]
\textbf{Situation :} Poste d'agent de saisie à la Mairie de Yaoundé IV. Deux candidates : \textbf{Marie} (BTS Secrétariat, 22 ans) et \textbf{Rose} (BAC G2, 24 ans, 2 ans expérience saisie à la CNPS).

\textbf{Lettre de Marie (extrait) :}
\begin{quote}
« Monsieur le Maire, je vous écris pour demander un travail. J'ai mon BTS Secrétariat l'année dernière. Je sais taper vite. Je suis sérieuse. Donnez-moi une chance. Merci. Marie. Tél : 6XX... »
\end{quote}

\textbf{Lettre de Rose (extrait) :}
\begin{quote}
« Yaoundé, le 10 janvier 2026

À l'attention de Monsieur le Maire de la Commune de Yaoundé IV
B.P. 334, Yaoundé

Objet : Candidature au poste d'agent de saisie -- Réf. Avis n° 045/MYIV/2025

Monsieur le Maire,

Titulaire d'un BAC G2 et forte de deux années d'expérience en saisie et traitement de données à la Caisse Nationale de Prévoyance Sociale (CNPS) -- direction des prestations familiales --, je souhaite mettre mes compétences au service de votre commune.

Au quotidien, j'assure la saisie de 150 dossiers/jour sous Excel et Access avec un taux d'erreur inférieur à 0,5\%. Ma rigueur, ma discrétion (données sensibles) et ma maîtrise des procédures administratives (archivage, classement) me permettent d'être opérationnelle immédiatement.

Disponible pour un entretien à votre convenance, je vous prie d'agréer, Monsieur le Maire, l'assurance de ma considération distinguée.

Rose MEKONDE
Tél. : 6XX XX XX XX -- Email : rose.mekonde@email.cm

Pièces jointes : CV, copie BAC, attestation CNPS, extrait acte de naissance, certificat nationalité.
\end{quote}
\tcblower
\textbf{Analyse comparative (grille d'évaluation) :}

\begin{center}
\begin{tabularx}{\linewidth}{|l|X|X|}
\hline
\rowcolor{popblueL}
\textbf{Critère} & \textbf{Marie} & \textbf{Rose} \\
\hline
En-tête complet & \textcolor{red}{Non} (pas de date, pas d'adresse destinataire) & \textcolor{popgreen}{Oui} (norme administrative) \\
\hline
Objet référencé & \textcolor{red}{Non} & \textcolor{popgreen}{Oui} (numéro avis) \\
\hline
Formule d'appel & \textcolor{orange}{Incomplète} & \textcolor{popgreen}{Correcte} \\
\hline
Introduction (qui/quoi) & \textcolor{orange}{Floue} & \textcolor{popgreen}{Claire : diplôme + expérience + poste} \\
\hline
Corps (arguments) & \textcolor{red}{Vague (« tape vite »)} & \textcolor{popgreen}{Chiffrés (150 dossiers/jour, 0,5\% erreur), logiciels, soft skills} \\
\hline
Adéquation poste & \textcolor{red}{Non démontrée} & \textcolor{popgreen}{Preuve d'opérationnalité immédiate} \\
\hline
Conclusion + appel entretien & \textcolor{red}{Absente} & \textcolor{popgreen}{Proactive + formule complète} \\
\hline
Pièces jointes listées & \textcolor{red}{Non} & \textcolor{popgreen}{Oui (dossier complet)} \\
\hline
\end{tabularx}
\end{center}

\textbf{Verdict :} Rose obtient l'entretien. Marie non. \textbf{La différence ne vient pas du diplôme (BTS vs BAC) mais de la qualité de l'argumentation et du respect des codes.}

\textbf{Leçon :} Une lettre de motivation n'est pas une supplique (« donnez-moi une chance »), c'est une \textbf{preuve écrite de valeur ajoutée}.
\end{exoresolu}

\courssub{Synthèse de la découverte : ce qu'il faut retenir avant d'écrire}

\begin{aretenir}[Check-list avant de rédiger sa lettre]
\begin{enumerate}
\item \textbf{J'ai lu et compris l'offre / l'avis / l'arrêté} (référence notée).
\item \textbf{Je connais le destinataire exact} (nom, fonction, adresse).
\item \textbf{J'ai identifié 2--3 exigences clés du poste} (compétences techniques, savoir-être, diplômes).
\item \textbf{J'ai sélectionné dans mon parcours 2--3 preuves concrètes} (stages, projets, résultats chiffrés) qui répondent à ces exigences.
\item \textbf{Je connais les formules d'appel et de politesse} adaptées au rang du destinataire.
\item \textbf{Mon CV est prêt, à jour, et cohérent avec la lettre.}
\item \textbf{J'ai rassemblé les pièces jointes exigées} (originaux + copies).
\end{enumerate}
\end{aretenir}

\begin{methode}[Préparer sa lettre en 4 étapes (méthode « PARC »)]
\textbf{P -- Prospecter :} Lire l'offre, noter la référence, identifier le destinataire, lister les mots-clés du poste.
\textbf{A -- Analyser son profil :} Faire le tableau « Exigences du poste $\leftrightarrow$ Mes preuves (diplômes, stages, résultats, qualités) ».
\textbf{R -- Rédiger la structure :} Appliquer le canevas en 3 parties (Intro -- Corps -- Conclusion) avec les formules rituelles.
\textbf{C -- Contrôler et finaliser :} Relire (orthographe, grammaire, ton), vérifier l'en-tête, l'objet, la signature, les pièces jointes, le format (PDF si email, papier qualité si dépôt physique).
\end{methode}

\medskip
La section découverte s'achève ici. Vous avez maintenant \textbf{vu} (textes iconiques), \textbf{comparé} (tableau), \textbf{modélisé} (parcours candidature), \textbf{défini} (fonctions), \textbf{vocabularisé} (lexique recrutement) et \textbf{évalué} (exercice résolu). La prochaine section abordera la \textbf{structure externe} : la mise en page normée, les formules d'appel et de politesse, l'en-tête parfait.

\coursec{La structure externe de la lettre de motivation}

Après avoir découvert, dans la section précédente, des lettres de motivation authentiques (textes iconiques 1 à 5) et analysé leur ton, leur vocabulaire et leur finalité, nous devons maintenant nous concentrer sur l'\emph{architecture visible} de la lettre. Avant même de lire le premier mot de votre argumentation, le recruteur scanne la page en quelques secondes : la disposition des coordonnées, la clarté de l'objet, la propreté de la mise en page. Cette « première impression visuelle » est souvent décisive. Comme le dit l'adage professionnel : \emph{« On n'a jamais une deuxième chance de faire une bonne première impression. »} Maîtriser la structure externe, c'est garantir que votre candidature passe le filtre du regard rapide pour accéder à la lecture approfondie.

\courssub{Les coordonnées de l'expéditeur et du destinataire : l'identification rigoureuse}

La lettre officielle obéit à un code spatial strict, hérité de la correspondance administrative et commerciale. En haut de page, deux blocs s'opposent symétriquement : l'expéditeur (vous) à gauche, le destinataire (l'employeur) à droite. Cette disposition n'est pas arbitraire ; elle permet un tri et un archivage immédiats.

\begin{definition}[Coordonnées de l'expéditeur (en-tête gauche)]
Elles identifient sans ambiguïté le candidat. Elles comprennent, dans l'ordre :
\begin{enumerate}
    \item \textbf{Nom et Prénom} (Nom en majuscules pour une lecture rapide : \textsc{MBARGA} Jean).
    \item \textbf{Adresse postale complète} (numéro, rue, quartier, ville, code postal si existant, pays).
    \item \textbf{Téléphone} (format international recommandé : +237 6xx xx xx xx).
    \item \textbf{Courriel professionnel} (évitez les adresses fantaisistes ; privilégiez \texttt{prenom.nom@domaine.com}).
\end{enumerate}
\end{definition}

\begin{definition}[Coordonnées du destinataire (en-tête droit)]
Elles désignent précisément la personne ou le service qui recevra la lettre. Elles comprennent :
\begin{enumerate}
    \item \textbf{Titre civil et fonction} (ex. : Monsieur le Directeur Général, Madame la Directrice des Ressources Humaines).
    \item \textbf{Nom de l'entreprise ou de l'organisme} (ex. : Cameroon Development Corporation).
    \item \textbf{Adresse de l'entreprise} (BP, ville).
\end{enumerate}
\end{definition}

\begin{attention}[Erreur fréquente : l'inversion ou l'oubli]
Ne placez jamais les coordonnées du destinataire à gauche ni les vôtres à droite. Ne zappez pas le titre de civilité (Monsieur, Madame) suivi de la fonction : écrire seulement « Direction des Ressources Humaines » est trop impersonnel et peut faire penser à un publipostage non ciblé.
\end{attention}

\begin{exemplebox}[En-tête complet pour une candidature à la CDC (Limbé)]
\begin{flushleft}
\texttt{\small
\textsc{MBARGA} Jean\\
Quartier Molyko, Rue 12\\
BP 456 Buea\\
Tel : +237 6 77 88 99 00\\
Email : jean.mbarga@email.cm\\
\hfill Monsieur le Directeur Général\\
\hfill Cameroon Development Corporation (CDC)\\
\hfill BP 180 Limbé\\
\hfill République du Cameroun
}
\end{flushleft}
\emph{Observation :} L'alignement à gauche pour l'expéditeur, à droite pour le destinataire. L'espace central reste libre pour l'objet et le corps de la lettre.
\end{exemplebox}

\begin{popfigure}
\begin{tikzpicture}[
    box/.style={draw=popblue, thick, rounded corners=3pt, fill=popblue!5, inner sep=4pt, font=\small},
    arrow/.style={->, >=Stealth, thick, poporange},
    label/.style={font=\tiny, text=popdark, anchor=west},
    title/.style={font=\bfseries\small, text=popblue}
]
% Page frame
\draw[popline, thick] (0,0) rectangle (16,10.5);
% Expéditeur (Left)
\node[box, anchor=north west, align=left] (exp) at (0.5,10) {
    \textbf{EXPÉDITEUR (Gauche)}\\
    \textsc{MBARGA} Jean\\
    Quartier Molyko, Rue 12\\
    BP 456 Buea\\
    Tel : +237 6 77 88 99 00\\
    Email : jean.mbarga@email.cm
};
% Destinataire (Right)
\node[box, anchor=north east, align=left] (dest) at (15.5,10) {
    \textbf{DESTINATAIRE (Droite)}\\
    Monsieur le Directeur Général\\
    Cameroon Development Corporation\\
    BP 180 Limbé\\
    République du Cameroun
};
% Lieu et Date (Right, below dest)
\node[box, anchor=north east, align=left] (lieu) at (15.5,6.8) {
    \textbf{LIEU ET DATE}\\
    Limbé, le 15 mars 2026
};
% Objet (Centered)
\node[box, anchor=north, align=center, fill=poporange!10, draw=poporange] (obj) at (8,6.8) {
    \textbf{OBJET}\\
    \underline{Candidature au poste de Technicien de Maintenance}\\
    \textit{Réf. : Annonce n°45 du 10/03/2026}
};
% Formule d'appel (Left)
\node[box, anchor=north west, align=left] (appel) at (0.5,5.8) {
    \textbf{FORMULE D'APPEL}\\
    Monsieur le Directeur Général,
};
% Corps (Placeholder)
\draw[popline, dashed] (0.5,5.4) -- (15.5,5.4);
\node[label] at (0.5,5.2) {\textit{Corps de la lettre (Introduction, Développement, Conclusion)}};
\draw[popline, dashed] (0.5,2.5) -- (15.5,2.5);
% Formule de politesse (Left)
\node[box, anchor=south west, align=left] (pol) at (0.5,2.3) {
    \textbf{FORMULE DE POLITESSE}\\
    Je vous prie d'agréer, Monsieur le Directeur Général,\\
    l'expression de ma considération distinguée.
};
% Signature (Right)
\node[box, anchor=south east, align=left] (sign) at (15.5,1.5) {
    \textbf{SIGNATURE}\\
    \textit{(Espace pour signature manuscrite)}\\
    \textsc{MBARGA} Jean
};
% Pièces jointes (Left bottom)
\node[box, anchor=south west, align=left] (pj) at (0.5,1.5) {
    \textbf{PIÈCES JOINTES}\\
    - CV détaillé\\
    - Copies des diplômes\\
    - Attestations de travail
};
% Arrows and labels for annotation
\draw[arrow] (exp.north west) -- ++(-0.5,0.5) node[label, anchor=south east] {1. Identité complète};
\draw[arrow] (dest.north east) -- ++(0.5,0.5) node[label, anchor=south west] {2. Cible précise};
\draw[arrow] (lieu.north east) -- ++(0.5,0.3) node[label, anchor=south west] {3. Datation};
\draw[arrow] (obj.north) -- ++(0,0.5) node[label, anchor=south] {4. Objet + Réf. (Gras/Souligné)};
\draw[arrow] (appel.north west) -- ++(-0.5,0.3) node[label, anchor=south east] {5. Appel = Politesse finale};
\draw[arrow] (pol.south west) -- ++(-0.5,-0.3) node[label, anchor=north east] {7. Politesse adaptée};
\draw[arrow] (sign.south east) -- ++(0.5,-0.3) node[label, anchor=north west] {8. Signature + Nom};
\draw[arrow] (pj.south west) -- ++(-0.5,-0.3) node[label, anchor=north east] {6. Pièces jointes listées};
\end{tikzpicture}
\legende{Schéma annoté de la structure externe d'une lettre de motivation : les 8 zones obligatoires et leur emplacement normé.}
\end{popfigure}

\courssub{Le lieu, la date, l'objet et les références : le repérage spatio-temporel et thématique}

Sous le bloc destinataire (ou parfois aligné à gauche sous l'expéditeur selon les normes d'entreprise, mais \emph{à droite} est la norme administrative camerounaise classique), on inscrit le lieu de rédaction et la date.

\begin{propriete}[Règle du lieu et de la date]
Format : \textbf{Ville, le [jour en chiffres] [mois en toutes lettres] [année en 4 chiffres]}.
Exemple : \emph{Douala, le 15 mars 2026}.
\emph{Attention :} Pas de « le » devant le jour si on écrit « Fait à Douala, le 15 mars 2026 ». Le mois s'écrit en minuscule. N'utilisez jamais le format numérique seul (15/03/2026) dans une lettre formelle.
\end{propriete}

L'\textbf{Objet} est la ligne la plus lue après le nom. Il doit être \textbf{court, précis, sans verbe}, et visuellement distinct (gras, souligné, ou majuscules).

\begin{methode}[Rédiger un objet efficace — 3 étapes]
\begin{enumerate}
    \item \textbf{Identifiez la nature de l'acte :} Candidature, Demande de stage, Réponse à annonce, Relance.
    \item \textbf{Précisez le poste visé :} Intitulé exact de l'offre (ex. : « Technicien de maintenance industrielle », « Assistant comptable »).
    \item \textbf{Ajoutez la référence de l'offre (si existante) :} « Réf. : Annonce n°45 du 10/03/2026 » ou « Réf. : Votre courrier n°123/DRH/2026 ».
\end{enumerate}
\end{methode}

\begin{attention}[Piège classique : l'objet « fourre-tout »]
\textbf{Interdit :} « Objet : Lettre de motivation », « Objet : Candidature », « Objet : Demande d'emploi ».
\textbf{Correct :} « Objet : Candidature au poste de Technicien de maintenance (Réf. Annonce n°45) ».
\emph{Pourquoi ?} Le recruteur gère souvent 5 postes différents simultanément. Un objet vague oblige à ouvrir la lettre pour savoir de quel poste il s'agit ; un objet précis permet le routage immédiat vers le bon service.
\end{attention}

\courssub{La formule d'appel, la mention des pièces jointes et la signature : la politesse codifiée}

La \textbf{formule d'appel} (ou salutation) ouvre le dialogue épistolaire. Elle doit être \emph{strictement symétrique} à la formule de politesse finale (la « formule de clôture »).

\begin{definition}[Formule d'appel]
C'est la salutation initiale adressée au destinataire. Elle reprend \emph{exactement} le titre et la fonction indiqués dans les coordonnées du destinataire, précédés de « Monsieur » ou « Madame ».
\begin{itemize}
    \item Coordonnées : \emph{Monsieur le Directeur des Ressources Humaines} $\rightarrow$ Appel : \textbf{Monsieur le Directeur des Ressources Humaines,}
    \item Coordonnées : \emph{Madame la Directrice Générale} $\rightarrow$ Appel : \textbf{Madame la Directrice Générale,}
    \item Si le nom est connu : \emph{Monsieur DUPONT, Directeur} $\rightarrow$ Appel : \textbf{Monsieur le Directeur,} (on évite « Monsieur Dupont » seul dans une lettre formelle).
\end{itemize}
\end{definition}

\begin{aretenir}[Règle d'or : Appel = Politesse]
La formule d'appel et la formule de politesse finale forment un couple indissociable. Si vous appelez « Monsieur le Directeur », vous terminez par « Je vous prie d'agréer, \textbf{Monsieur le Directeur}, l'expression de... ». Toute asymétrie (ex. Appel : « Monsieur », Politesse : « Monsieur le Directeur ») trahit un manque de rigueur.
\end{aretenir}

La \textbf{mention des pièces jointes} (PJ) se place en bas à gauche, après la formule de politesse (ou parfois à droite sous la signature). Elle énumère \emph{précisément} les documents joints : « CV détaillé », « Copie certifiée conforme du BTS Maintenance », « Attestation de stage chez ENEO », etc. Évitez « PJ : 3 pièces » sans détail.

La \textbf{signature} (manuscrite pour un envoi postal, scannée ou typographiée pour un email) se place en bas à droite. Elle est précédée de votre nom tapé en majuscules pour lisibilité.

\begin{savaistu}[Les 30 secondes fatidiques]
Des études en psychologie cognitive appliquée au recrutement (eye-tracking) montrent qu'un recruteur passe en moyenne \textbf{6 à 30 secondes} sur le premier tri visuel d'une lettre. Il ne « lit » pas : il \emph{scanne} la structure externe (Objet, Nom, Mise en page, Longueur). Une lettre mal alignée, tachée, sans objet clair, ou avec une formule d'appel générique (« Madame, Monsieur » au lieu de « Madame la Directrice ») est souvent classée « non pertinent » avant même que le premier paragraphe ne soit lu. \textbf{La forme est le premier message de votre compétence professionnelle.}
\end{savaistu}

\courssub{Mise en page et présentation : aération, paragraphes et propreté}

La structure externe ne se limite pas aux champs obligatoires ; elle englobe la \emph{géométrie de la page}. Une lettre lisible respire.

\begin{propriete}[Normes de mise en page professionnelle]
\begin{itemize}
    \item \textbf{Marges :} 2,5 cm minimum de chaque côté (gauche, droite, haut, bas).
    \item \textbf{Police :} Classique, serif (Times New Roman, Garamond) ou sans serif sobre (Arial, Calibri, Helvetica), corps 11 ou 12 pt. \emph{Une seule police dans toute la lettre.}
    \item \textbf{Interligne :} 1,15 ou 1,5 pour l'aération.
    \item \textbf{Alignement :} Justifié (bloc rectangle net) ou aligné à gauche (plus moderne, moins de « rivières » blanches). Jamais centré pour le corps.
    \item \textbf{Paragraphes :} Séparés par une ligne blanche (saut de ligne double), \textbf{sans retrait (alinéa)} pour le style bloc (style anglo-saxon/administratif moderne) ou \textbf{avec retrait} pour le style français classique. Choisissez une convention et tenez-la.
    \item \textbf{Longueur :} \textbf{Une page maximum} (recto seul). Une lettre de motivation n'est pas un roman.
    \item \textbf{Propreté :} Zéro rature, zéro tache, impression laser/jet d'encre qualité haute sur papier blanc 80-90g (pour envoi postal). Pour email : PDF unique (Lettre + CV) nommé \texttt{Nom\_Prenom\_Poste.pdf}.
\end{itemize}
\end{propriete}

\begin{popfigure}
\begin{tikzpicture}[
    page/.style={draw=popdark, thick, rectangle, minimum width=7cm, minimum height=9.5cm, inner sep=5pt},
    bad/.style={fill=poppink!20, draw=poppink, thick},
    good/.style={fill=popgreen!20, draw=popgreen, thick},
    txt/.style={font=\ttfamily\footnotesize, align=left},
    title/.style={font=\bfseries\small, anchor=south}
]
% MAUVAISE MISE EN PAGE
\node[page, bad] (badpage) at (0,0) {};
\node[title, text=poppink] at (badpage.north) {MAUVAISE MISE EN PAGE (Échec visuel)};
\node[txt, anchor=north west, text width=6.5cm, text=popdark] at ([yshift=-5pt]badpage.north west) {
MARGES TROP ETROITES (1cm)
POLICE FANTAISISTE (Comic Sans)
INTERLIGNE SERRÉ (1.0)
OBJET: lettre de motivation
Monsieur,
Je vous ecris pour postuler au poste...
[Texte collé, pas de paragraphes distincts, phrases longues sans respiration]
Je suis tres motive et dynamique.
Cordialement,
Jean Mbarga
PJ: pieces
};
% BONNE MISE EN PAGE
\node[page, good] (goodpage) at (8.5,0) {};
\node[title, text=popgreen] at (goodpage.north) {BONNE MISE EN PAGE (Réussite visuelle)};
\node[txt, anchor=north west, text width=6.5cm, text=popdark] at ([yshift=-5pt]goodpage.north west) {
MARGES NORMALES (2.5cm)
POLICE SOBRE (Arial 11)
INTERLIGNE 1.15
\textbf{OBJET : Candidature Technicien Maintenance (Réf. 45)}
Monsieur le Directeur Général,
[Paragraphes aérés, blocs distincts]
\textbf{Introduction :} Accroche + Poste visé.
\textbf{Développement :} Compétences techniques (BTS) + Expérience stage ENEO + Savoir-être.
\textbf{Conclusion :} Disponibilité entretien + Formule politesse complète.
Je vous prie d'agréer, Monsieur le Directeur Général, l'expression de ma considération distinguée.
Limbé, le 15 mars 2026
\textsc{MBARGA} Jean
\textit{(Signature)}
\textbf{Pièces jointes :}
- CV détaillé
- Copie BTS Maintenance
- Attestation stage ENEO
};
\end{tikzpicture}
\legende{Comparaison visuelle : à gauche, une lettre « tueuse » (marges étroites, police fantaisiste, objet vague, pas de paragraphes) ; à droite, une lettre « professionnelle » (aérée, structurée, objet précis, politesse adaptée).}
\end{popfigure}

\courssub{Exercice résolu : Audit de structure externe}

\begin{exoresolu}[Analyse d'une en-tête défaillante]
\textbf{Énoncé :} Voici l'en-tête d'une lettre envoyée par un candidat, Pierre \textsc{NGONO}, pour un poste de « Commercial terrain » à la brasserie \textsc{UCB} (Douala). Identifiez 5 erreurs ou manquements à la structure externe.

\begin{flushleft}
\texttt{\small
Pierre Ngono\\
Akwa, Douala\\
06 99 88 77 66
\\
Monsieur le Directeur\\
UCB\\
Douala
\\
Objet : Candidature
\\
Monsieur,...
}
\end{flushleft}

\textbf{Solution détaillée :}
\begin{enumerate}
    \item \textbf{Nom incomplet :} « Pierre Ngono » au lieu de \textsc{NGONO} Pierre (Nom en majuscules pour repérage rapide).
    \item \textbf{Adresse incomplète :} Pas de BP, pas de quartier précis, pas de pays. Téléphone sans indicatif pays (+237).
    \item \textbf{Coordonnées destinataire insuffisantes :} « Monsieur le Directeur » (fonction floue : Directeur Général ? Commercial ? RH ?), pas de nom d'entreprise complet (Union Camerounaise de Brasseries), pas d'adresse (BP, ville).
    \item \textbf{Objet trop vague :} « Candidature » ne dit pas \emph{à quel poste}. Doit être : « Candidature au poste de Commercial Terrain (Réf. Annonce n°...) ».
    \item \textbf{Formule d'appel asymétrique :} L'en-tête dit « Monsieur le Directeur », l'appel dit « Monsieur, ». Si l'en-tête vise le Directeur Général, l'appel doit être « Monsieur le Directeur Général, ». Si c'est le DRH, « Monsieur le Directeur des Ressources Humaines, ».
    \item \textbf{(Bonus) :} Absence de Lieu/Date. Absence de mention PJ. Police non spécifiée mais présentation « bloc-notes » peu professionnelle.
\end{enumerate}
\end{exoresolu}

\begin{methode}[Grille de vérification « Pré-vol » — 8 points à cocher avant d'envoyer]
Imprimez cette grille (mentalement ou sur papier) et cochez \textbf{chaque case} avant de cliquer sur « Envoyer » ou de poster l'enveloppe.
\begin{center}
\begin{tabularx}{\linewidth}{|c|X|c|}\hline
\textbf{N°} & \textbf{Élément à vérifier} & \textbf{OK ?} \\ \hline
1 & \textbf{Expéditeur :} Nom (MAJ), Prénom, Adresse complète, Tel (+237), Email pro. & $\square$ \\ \hline
2 & \textbf{Destinataire :} Civilité + Fonction exacte, Entreprise, Adresse complète. & $\square$ \\ \hline
3 & \textbf{Lieu \& Date :} « Ville, le [jour] [mois] [année] » (à droite). & $\square$ \\ \hline
4 & \textbf{Objet :} Précis (Poste exact), Gras/Souligné, Référence annonce si existe. & $\square$ \\ \hline
5 & \textbf{Appel :} Reprend \emph{exactement} la civilité/fonction du destinataire (ex: « Monsieur le Directeur Général, »). & $\square$ \\ \hline
6 & \textbf{Corps :} Aéré (paragraphes séparés), Justifié/Gauche, Police 11/12, 1 page max. & $\square$ \\ \hline
7 & \textbf{Politesse :} Symétrique à l'appel, complète, adaptée (pas « Cordialement » seul). & $\square$ \\ \hline
8 & \textbf{Signature + PJ :} Nom tapé + Signature (manuscrite/scan), Liste détaillée des PJ. & $\square$ \\ \hline
\end{tabularx}
\end{center}
\emph{Si une case n'est pas cochée : corrigez avant envoi.}
\end{methode}

\courssub{Transition vers la structure interne}

La structure externe est désormais maîtrisée : votre lettre a une « tenue » impeccable, elle respecte les codes, elle invite à la lecture. Mais l'enveloppe, si belle soit-elle, ne suffit pas. Il faut maintenant remplir l'intérieur avec une argumentation qui tient la route. La prochaine section, \emph{« La structure interne : introduction et conclusion »}, nous apprendra à construire le « squelette » argumentatif : comment accrocher dès la première phrase (introduction) et comment partir sur une note professionnelle engageante (conclusion), avant d'aborder le cœur du réacteur : le corps de la lettre (développement).

\coursec{La structure interne : introduction et conclusion}

Dans la section précédente, nous avons étudié la \cle{structure externe} de la lettre de motivation : l'en-tête (expéditeur, destinataire, lieu, date), l'objet, la formule d'appel. Ces éléments forment le « cadre » formel, indispensable mais insuffisant. C'est le \cle{texte} lui-même — la \cle{structure interne} — qui persuade le recruteur. Une lettre de motivation n'est pas un simple résumé de parcours : c'est un discours persuasif organisé en trois temps, que les spécialistes résument par la formule \emph{Vous – Moi – Nous}. Nous allons maîtriser cette organisation en nous concentrant sur ses deux extrémités : l'\cle{introduction} et la \cle{conclusion}.

\courssub{Les trois parties du texte : le schéma « Vous – Moi – Nous »}

\begin{definition}[Structure interne de la lettre de motivation]
La \cle{structure interne} organise le corps du texte en trois parties fonctionnelles :
\begin{itemize}
\item l'\textbf{introduction} (le \emph{Vous}) : on parle de l'entreprise, du poste, de l'offre ;
\item le \textbf{corps de la lettre} (le \emph{Moi}) : on présente ses diplômes, compétences, qualités, expériences en lien avec le poste ;
\item la \textbf{conclusion} (le \emph{Nous}) : on évoque la collaboration future et on sollicite un entretien.
\end{itemize}
\end{definition}

Pourquoi cet ordre ? Lisez ces deux ouvertures :

\begin{itemize}
\item Lettre A : « Je m'appelle Mbarga Jean, j'ai 22 ans, je cherche du travail... »
\item Lettre B : « Votre entreprise, leader camerounais de la distribution d'électricité, recrute des agents commerciaux dynamiques : cette annonce a immédiatement retenu mon attention... »
\end{itemize}

Laquelle donne envie de continuer ? La lettre B : elle montre que le candidat s'intéresse d'abord à l'\emph{entreprise}. Le recruteur, flatté et intrigué, poursuit sa lecture. Le schéma « Vous – Moi – Nous » repose sur cette psychologie : capter l'attention en parlant du destinataire, démontrer sa valeur, projeter une relation de travail.

\begin{methode}[Le schéma « Vous – Moi – Nous » en trois étapes]
Pour organiser le texte de votre lettre :
\begin{enumerate}
\item \textbf{Le « Vous » (introduction)} : parlez de l'entreprise. Montrez que vous la connaissez (activité, réputation, offre) et annoncez votre candidature.
\item \textbf{Le « Moi » (corps)} : présentez votre profil. Reliez chaque diplôme, compétence, qualité aux besoins du poste. Ne répétez pas votre CV : commentez-le.
\item \textbf{Le « Nous » (conclusion)} : projetez la collaboration. Proposez un entretien, affirmez votre disponibilité, remerciez, concluez par la formule de politesse.
\end{enumerate}
\end{methode}

\begin{popfigure}
\begin{center}
\begin{tikzpicture}[node distance=0.4cm]
\node[draw=popblue, fill=popblueL, rounded corners, minimum width=3.6cm, minimum height=2.4cm, align=center] (vous) {\textbf{VOUS}\\[2pt] \emph{Introduction}\\[3pt] \footnotesize Parler de l'entreprise,\\ \footnotesize de l'offre, du poste.\\ \footnotesize Accrocher le lecteur.};
\node[draw=poporange, fill=poporangeL, rounded corners, minimum width=3.6cm, minimum height=2.4cm, align=center, right=0.9cm of vous] (moi) {\textbf{MOI}\\[2pt] \emph{Corps de la lettre}\\[3pt] \footnotesize Diplômes, compétences,\\ \footnotesize qualités, expériences\\ \footnotesize au service du poste.};
\node[draw=popgreen, fill=popgreenL, rounded corners, minimum width=3.6cm, minimum height=2.4cm, align=center, right=0.9cm of moi] (nous) {\textbf{NOUS}\\[2pt] \emph{Conclusion}\\[3pt] \footnotesize Entretien souhaité,\\ \footnotesize disponibilité,\\ \footnotesize remerciement, politesse.};
\draw[-{Stealth}, very thick, popblue] (vous) -- (moi);
\draw[-{Stealth}, very thick, poporange] (moi) -- (nous);
\end{tikzpicture}
\end{center}
\legende{Le schéma « Vous – Moi – Nous » : les trois parties du texte et la fonction de chacune.}
\end{popfigure}

\courssub{La rédaction de l'introduction : capter l'attention}

L'introduction est la partie la plus délicate : le recruteur lit souvent des dizaines de lettres ; les trois premières lignes décident du sort de la vôtre. Une bonne introduction remplit \textbf{trois fonctions} :

\begin{enumerate}
\item \textbf{L'accroche} : première phrase qui attire l'attention et donne envie de lire la suite ;
\item \textbf{La référence à l'offre ou à l'entreprise} : on précise l'annonce (date, source, intitulé du poste) ou l'entreprise visée (candidature spontanée) ;
\item \textbf{L'annonce de la candidature} : on déclare clairement qu'on postule, avec une première raison.
\end{enumerate}

\begin{definition}[L'accroche]
L'\cle{accroche} est la première phrase. Elle doit être brève, précise, personnelle. Trois techniques principales :
\begin{itemize}
\item la \textbf{question} : « Cherchez-vous un agent commercial capable de fidéliser une clientèle exigeante ? » ;
\item le \textbf{fait marquant de l'entreprise} : « L'extension récente de votre réseau dans la région du Littoral témoigne du dynamisme de votre entreprise. » ;
\item la \textbf{motivation directe} : « Technicien en froid et climatisation, je souhaite mettre mon savoir-faire au service de votre société. »
\end{itemize}
\end{definition}

\begin{attention}[Formules banales à bannir]
Évitez les accroches vides, lues cent fois : « Je me permets de vous écrire... », « Suite à votre annonce parue récemment... », « Je viens très respectueusement solliciter un emploi... ». Elles ne disent rien de vous ni de l'entreprise. Préférez une phrase \emph{concrète} citant un fait, un chiffre, une qualité précise.
\end{attention}

\begin{exemplebox}[Introduction réussie, commentée]
Candidature à un poste d'agent commercial à ENEO (candidat : Baccalauréat technique G2) :

\emph{« Entreprise de référence dans la distribution de l'énergie électrique au Cameroun, ENEO recherche des agents commerciaux pour renforcer ses équipes de la région du Centre. Titulaire d'un Baccalauréat technique G2 (Gestion – Comptabilité) et fort d'un stage de six mois dans une agence bancaire de Yaoundé, je vous propose ma candidature à ce poste. »}

\begin{itemize}
\item \textbf{Phrase 1 (accroche + référence)} : le candidat parle de l'entreprise (« Entreprise de référence... ») et cite l'offre avec précision (poste, région). Le recruteur voit que la lettre n'est pas un modèle-type.
\item \textbf{Phrase 2 (annonce de la candidature)} : deux atouts majeurs (diplôme, expérience) puis annonce claire : « je vous propose ma candidature ». En deux phrases : qui écrit, pour quel poste, avec quels arguments.
\end{itemize}
\end{exemplebox}

\begin{exemplebox}[Introduction pour candidature spontanée]
\emph{« Votre récente implantation à Douala-Bassa, annoncée dans la presse économique, confirme l'ambition de la SABC de renforcer sa présence logistique dans le Littoral. Technicien de maintenance titulaire d'un Baccalauréat technique F3 (Électromécanique), je souhaite anticiper vos besoins en proposant ma candidature pour tout poste correspondant à mon profil. »}

\textit{Analyse} : accroche par fait marquant (implantation), référence à l'actualité de l'entreprise, annonce de candidature avec diplôme et posture proactive (« anticiper vos besoins »).
\end{exemplebox}

\courssub{La rédaction de la conclusion : projeter la collaboration}

La conclusion est le « Nous » : elle doit laisser une impression positive et provoquer une réponse. Elle comprend \textbf{quatre éléments}, dans cet ordre :

\begin{enumerate}
\item \textbf{La demande d'entretien} : « Je serais heureux de vous exposer ma motivation lors d'un entretien que vous voudrez bien m'accorder. » ;
\item \textbf{La disponibilité} : « Je reste à votre entière disposition pour tout renseignement complémentaire. » ;
\item \textbf{Le remerciement} : « Vous remerciant de l'attention portée à ma candidature... » ;
\item \textbf{La formule de politesse} : phrase de déférence qui clôt officiellement la lettre.
\end{enumerate}

\begin{attention}[Erreur fréquente : fin familière]
Dans une lettre officielle camerounaise, il est \textbf{interdit} de terminer par « Merci d'avance », « Cordialement », « Bien à vous », « Amitiés ». Ces formules appartiennent au courrier familier ou au courriel. La lettre de motivation exige une \emph{formule de déférence complète}, développée, reprenant le titre exact du destinataire.
\end{attention}

\courssub{Les formules de politesse officielles}

\begin{propriete}[Cohérence appel / formule de politesse]
La formule de politesse reprend \textbf{toujours le titre exact} de la formule d'appel : même titre, même personne, même fonction.
\begin{itemize}
\item Appel : « Monsieur le Directeur » $\rightarrow$ Politesse : « Je vous prie d'agréer, \textbf{Monsieur le Directeur}, l'expression de ma haute considération. »
\item Appel : « Madame la Directrice des Ressources Humaines » $\rightarrow$ Politesse : « Je vous prie d'agréer, \textbf{Madame la Directrice des Ressources Humaines}, l'expression de ma haute considération. »
\end{itemize}
Écrire « Monsieur le Directeur » en appel et « Madame » dans la politesse est une faute grave d'incohérence.
\end{propriete}

Les formules varient selon le degré de déférence :

\begin{itemize}
\item \textbf{Supérieur hiérarchique ou haute autorité} (ministre, directeur général, recteur) : « Je vous prie d'agréer, Monsieur le Ministre, l'expression de ma \emph{très haute considération}. » ;
\item \textbf{Destinataire de rang équivalent ou responsable d'entreprise} : « Je vous prie d'agréer, Monsieur le Directeur, l'expression de ma \emph{haute considération} » ou « ... de mes \emph{salutations distinguées} » ;
\item \textbf{Administration en général} : « Je vous prie de croire, Monsieur le Préfet, en l'assurance de ma \emph{considération distinguée}. »
\end{itemize}

\begin{popfigure}
\begin{center}
\begin{tikzpicture}[grow=right, level distance=4.2cm,
level 1/.style={sibling distance=2.6cm},
edge from parent/.style={draw=popmuted, thick, -{Stealth}},
every node/.style={align=center}]
\node[draw=popdark, fill=popcream, rounded corners, font=\bfseries, align=center]{Formule\\ de politesse}
child { node[draw=popblue, fill=popblueL, rounded corners, font=\footnotesize, align=center]{Administration\\ « ... l'assurance de ma\\ considération distinguée. »} }
child { node[draw=poporange, fill=poporangeL, rounded corners, font=\footnotesize, align=center]{Responsable / égal\\ « ... l'expression de ma\\ haute considération »\\ ou « mes salutations distinguées »} }
child { node[draw=popgreen, fill=popgreenL, rounded corners, font=\footnotesize, align=center]{Supérieur / haute autorité\\ « ... l'expression de ma\\ très haute considération. »} };
\end{tikzpicture}
\end{center}
\legende{Arbre de choix des formules de politesse selon le destinataire. Dans tous les cas, le titre de l'appel est repris mot pour mot.}
\end{popfigure}

\begin{savaistu}[La formule de politesse au Probatoire / Baccalauréat technique]
À l'examen, la présentation de la lettre officielle fait l'objet d'une notation précise. L'\emph{absence} de formule, une formule \emph{familière} (« Cordialement ») ou une \emph{incohérence} appel/politesse coûtent des points de présentation, même si le contenu est excellent. Relisez toujours ces deux lignes avant de rendre votre copie !
\end{savaistu}

\begin{exoresolu}[Améliorer une introduction plate et une conclusion familière]
\textbf{Énoncé.} Début et fin de la lettre d'Etoundi, candidat technicien de maintenance à la SABC Douala :

\emph{Introduction : « Je me permets de vous écrire cette lettre pour vous demander du travail dans votre société. J'ai vu votre annonce quelque part. »}

\emph{Conclusion : « Voilà, j'espère que vous allez me prendre. Merci d'avance. Cordialement. »}

Réécrivez selon les règles.

\tcblower
\textbf{Solution détaillée.}

\textbf{Introduction corrigée} (accroche par fait marquant + référence + annonce) :

« Leader camerounais de l'industrie des boissons, votre société recherche, par annonce parue dans \emph{Cameroon Tribune} du 12 mars, un technicien de maintenance pour son site de Douala-Bassa. Titulaire d'un Baccalauréat technique F3 (Électromécanique) et formé à la maintenance des lignes de production lors d'un stage à la Chocolaterie Confiserie du Cameroun, je vous propose ma candidature à ce poste. »

\textit{Commentaire} : accroche banale remplacée par fait précis ; offre citée avec source et date ; candidature annoncée avec deux arguments concrets.

\textbf{Conclusion corrigée} (entretien + disponibilité + remerciement + déférence) :

« Je serais honoré de vous exposer ma motivation lors d'un entretien que vous voudrez bien m'accorder, et je reste à votre entière disposition pour tout renseignement complémentaire. Vous remerciant de l'attention portée à ma candidature, je vous prie d'agréer, Monsieur le Directeur des Ressources Humaines, l'expression de ma haute considération. »

\textit{Commentaire} : « Merci d'avance » et « Cordialement » supprimés ; politesse reprend exactement le titre de l'appel.
\end{exoresolu}

\begin{exercice}[À vous de rédiger]
Vous répondez à une offre de la CAMTEL pour un poste d'agent d'accueil à Bafoussam. Rédigez :
\begin{enumerate}
\item une introduction complète (accroche par question, référence à l'offre, annonce de la candidature) ;
\item une conclusion complète avec formule de politesse adaptée, la lettre étant adressée à « Madame la Cheffe d'Agence ».
\end{enumerate}
\end{exercice}

\begin{corrige}[Exercice : introduction et conclusion]
\textbf{Introduction possible} : « Souhaitez-vous offrir à vos abonnés de Bafoussam un accueil professionnel et chaleureux ? Votre annonce parue le 5 avril, concernant le recrutement d'un agent d'accueil à l'agence CAMTEL de Bafoussam, correspond exactement à mon profil. Titulaire d'un Baccalauréat technique G1 et parlant couramment le français et l'anglais, je vous propose ma candidature à ce poste. »

\textbf{Conclusion possible} : « Je serais heureuse de vous présenter ma motivation lors d'un entretien, et je demeure à votre disposition pour tout complément d'information. Vous remerciant de l'attention accordée à ma candidature, je vous prie d'agréer, Madame la Cheffe d'Agence, l'expression de ma haute considération. »

\textit{Vérification} : politesse reprend exactement le titre de l'appel ; aucune formule familière.
\end{corrige}

\begin{aretenir}[L'essentiel de la structure interne]
\begin{itemize}
\item Le texte suit le schéma \textbf{Vous} (introduction : l'entreprise) – \textbf{Moi} (corps : le candidat) – \textbf{Nous} (conclusion : la collaboration).
\item L'introduction = accroche originale + référence précise à l'offre + annonce de la candidature.
\item La conclusion = demande d'entretien + disponibilité + remerciement + formule de politesse.
\item La formule de politesse reprend \emph{toujours} le titre exact de la formule d'appel ; jamais de « Cordialement » ni de « Merci d'avance » dans une lettre officielle.
\end{itemize}
\end{aretenir}

\coursec{Le corps de la lettre : convaincre avec style}

Dans la section précédente, nous avons appris à construire l'introduction (qui accroche le lecteur en annonçant le poste visé) et la conclusion (qui sollicite un entretien avec une formule de politesse adaptée). Mais entre ces deux portes, il reste le cœur du bâtiment : le \cle{corps de la lettre}. C'est là que tout se joue. Le recruteur y cherche une réponse à trois questions : \emph{que sait faire ce candidat ? qu'a-t-il déjà fait ? pourquoi veut-il travailler chez nous, précisément ?} Cette section vous apprend à répondre à ces trois questions en trois paragraphes, avec des connecteurs logiques, des figures de style au service de la persuasion, et un registre soutenu qui inspire confiance.

\courssub{L'architecture du corps : trois paragraphes, trois fonctions}

Une lettre de motivation efficace n'est pas un inventaire : c'est une \cle{démonstration}. Comme dans un raisonnement, chaque paragraphe a une fonction précise et prépare le suivant.

\begin{definition}[Le corps de la lettre de motivation]
Le corps de la lettre de motivation est la partie centrale du texte, située entre l'introduction et la conclusion. Il comporte généralement \textbf{trois paragraphes} :
\begin{enumerate}
\item la \cle{formation} : diplômes et études en lien direct avec le poste ;
\item l'\cle{expérience} : stages, jobs, activités associatives, projets scolaires ;
\item les \cle{qualités personnelles} et la \cle{motivation} pour cette entreprise précise.
\end{enumerate}
Chaque paragraphe répond à une attente du recruteur : \emph{est-il formé ? est-il expérimenté ? est-il fait pour nous ?}
\end{definition}

\begin{popfigure}
\begin{center}
\begin{tikzpicture}[
  bloc/.style={draw=popblue, fill=popblueL, rounded corners, thick, text width=4.6cm, minimum height=1.7cm, align=center, font=\small},
  conn/.style={draw=poporange, fill=poporangeL, rounded corners, font=\small\itshape, text=popdark},
  fleche/.style={-{Stealth[length=3mm]}, very thick, poporange}
]
\node[bloc, align=center] (p1) at (0,0){\textbf{Paragraphe 1}\\ La formation et les diplômes en lien avec le poste};
\node[bloc, align=center] (p2) at (0,-2.8){\textbf{Paragraphe 2}\\ Les expériences : stages, jobs, projets, associations};
\node[bloc, align=center] (p3) at (0,-5.6){\textbf{Paragraphe 3}\\ Les qualités prouvées + la motivation pour l'entreprise};
\node[conn] (c1) at (5.3,-1.4) {ensuite, en outre};
\node[conn] (c2) at (5.3,-4.2) {c'est pourquoi, ainsi};
\draw[fleche] (p1) -- (p2);
\draw[fleche] (p2) -- (p3);
\end{tikzpicture}
\end{center}
\legende{Schéma du corps de la lettre : trois paragraphes enchaînés (formation $\rightarrow$ expérience $\rightarrow$ qualités et motivation), reliés par des connecteurs logiques.}
\end{popfigure}

\courssub{Paragraphe 1 : la formation et les diplômes en lien avec le poste}

L'erreur du débutant consiste à énumérer tous ses diplômes depuis le BEPC. Or le recruteur ne veut pas votre biographie scolaire : il veut savoir si votre formation \emph{correspond au poste}. Il faut donc \textbf{sélectionner} le diplôme le plus pertinent et \textbf{commenter} ce qu'il vous a appris d'utile pour l'emploi visé.

\begin{exemplebox}[Formation mise en lien avec le poste]
Pour un poste de technicien en maintenance à Douala :
\begin{quote}
« Titulaire d'un baccalauréat technique en électrotechnique obtenu au lycée technique de Bonabéri, j'ai acquis une solide maîtrise de la lecture de schémas électriques et du dépannage des installations industrielles, compétences directement utiles au poste de technicien de maintenance que vous proposez. »
\end{quote}
Remarquez la structure : \textbf{diplôme} $\rightarrow$ \textbf{compétences précises} $\rightarrow$ \textbf{lien explicite avec le poste}.
\end{exemplebox}

\begin{attention}[Le piège du catalogue]
Ne listez jamais tous vos diplômes et toutes vos matières. Une phrase comme « j'ai fait la sixième, puis la cinquième, puis la quatrième... » endort le lecteur et ne prouve rien. Choisissez \textbf{un ou deux diplômes}, ceux qui parlent au poste, et dites ce qu'ils vous ont réellement apporté.
\end{attention}

\courssub{Paragraphe 2 : les expériences qui prouvent}

Le deuxième paragraphe répond à la question : \emph{ce candidat a-t-il déjà pratiqué ?} Pour un élève de Première technique, l'expérience ne signifie pas forcément un emploi salarié : un \textbf{stage} en entreprise, un \textbf{job} de vacances, la gestion de la caisse d'un club, un \textbf{projet scolaire} (confection d'un circuit, organisation d'une journée portes ouvertes) sont des expériences valables, à condition de les présenter avec des faits.

\begin{methode}[Présenter une expérience en trois temps]
\begin{enumerate}
\item \textbf{Situer} : où, quand, dans quel cadre ? (« Lors de mon stage à la menuiserie Ets Ngo Bell, pendant les vacances de 2025... »)
\item \textbf{Décrire la tâche} : qu'avez-vous fait concrètement ? (« ...j'ai participé à la fabrication et à la pose de douze portes en bois massif... »)
\item \textbf{Dégager l'acquis} : quelle compétence cela prouve-t-il ? (« ...ce qui m'a appris la précision des mesures et le respect des délais. »)
\end{enumerate}
\end{methode}

\courssub{Paragraphe 3 : les qualités prouvées et la motivation}

C'est le paragraphe décisif : celui où le candidat cesse de parler de son passé pour parler de sa \textbf{valeur} et de son \textbf{désir}. Deux règles d'or : toute qualité affirmée doit être prouvée par un fait, et la motivation doit viser \emph{cette} entreprise, pas une entreprise quelconque.

\begin{methode}[Une qualité = une preuve]
Pour chaque qualité affirmée, apportez une preuve concrète et vérifiable :
\begin{itemize}
\item « rigoureux » $\rightarrow$ « j'ai tenu la caisse du club informatique sans aucune erreur pendant une année scolaire » ;
\item « sens du travail en équipe » $\rightarrow$ « j'ai coordonné un groupe de cinq élèves pour le projet d'éclairage de la salle des fêtes » ;
\item « ponctuel et fiable » $\rightarrow$ « je n'ai manqué aucun rendez-vous de stage en deux mois ».
\end{itemize}
Une qualité sans preuve est une prétention ; une qualité prouvée est un argument.
\end{methode}

\begin{savaistu}[La personnalisation paie]
Les recruteurs privilégient les lettres personnalisées : mentionner le \textbf{nom exact de l'entreprise} et l'une de ses \textbf{activités réelles} augmente nettement les chances d'obtenir un entretien. Une lettre qui commence par « votre entreprise, leader régional de la transformation du cacao » montre que le candidat s'est renseigné ; une lettre interchangeable envoyée à vingt adresses se reconnaît immédiatement... et finit à la corbeille.
\end{savaistu}

\courssub{L'articulation logique : les connecteurs}

Un corps de lettre n'est pas une pile de phrases : c'est un chemin. Les \cle{connecteurs logiques} sont les panneaux qui guident le lecteur d'une idée à l'autre.

\begin{aretenir}[Les connecteurs utiles au corps de la lettre]
\begin{itemize}
\item \textbf{Ordre et enchaînement} : \emph{d'abord, ensuite, puis, enfin} ;
\item \textbf{Addition d'arguments} : \emph{en outre, de plus, également, par ailleurs} ;
\item \textbf{Conséquence} : \emph{c'est pourquoi, ainsi, de ce fait, par conséquent} ;
\item \textbf{Opposition valorisante} : \emph{cependant, néanmoins, bien que}.
\end{itemize}
Exemple d'enchaînement : « \emph{D'abord} formé à la comptabilité, j'ai \emph{ensuite} pratiqué la saisie des factures en stage ; \emph{en outre}, je maîtrise le tableur ; \emph{c'est pourquoi} je serai immédiatement opérationnel dans votre service. »
\end{aretenir}

\courssub{La stylistique au service de la persuasion}

Convaincre, ce n'est pas seulement aligner des faits : c'est les \emph{mettre en valeur}. Trois figures de style, au programme de Première technique, sont particulièrement efficaces dans une lettre de motivation.

\begin{definition}[Gradation, antithèse, métonymie]
\begin{itemize}
\item La \cle{gradation} est une énumération dont les termes vont en \textbf{s'intensifiant} (gradation ascendante) ou en s'affaiblissant (gradation descendante). Exemple : « technicien appliqué, rigoureux, passionné, indispensable ».
\item L'\cle{antithèse} oppose deux idées pour en \textbf{valoriser une} par contraste. Exemple : « jeune, mais déjà expérimenté ; novice en apparence, mais formé aux gestes du métier ».
\item La \cle{métonymie} désigne une réalité par un \textbf{terme voisin} (la partie pour le tout, le contenant pour le contenu). Exemple : « rejoindre vos équipes » (les équipes pour l'entreprise), « mettre mes mains au service de votre atelier » (les mains pour le travail).
\end{itemize}
\end{definition}

\begin{popfigure}
\begin{center}
\begin{tikzpicture}[
  marche/.style={draw=popgreen, thick, fill=popgreenL, rounded corners, align=center, font=\small, minimum height=0.9cm}
]
\draw[-{Stealth}, very thick, popdark] (0,0) -- (10.5,0) node[below left, font=\small]{intensité croissante};
\node[marche, text width=2cm] at (1.3,0.7) {appliqué};
\node[marche, text width=2cm] at (3.9,1.5) {rigoureux};
\node[marche, text width=2cm] at (6.5,2.3) {passionné};
\node[marche, text width=2.2cm, fill=popgoldL, draw=popgold] at (9.1,3.1) {\textbf{indispensable}};
\draw[popgreen, thick, dashed] (0.4,0.2) -- (9.9,3.5);
\end{tikzpicture}
\end{center}
\legende{Gradation ascendante : chaque terme monte d'un cran en intensité, jusqu'au mot-clé final « indispensable », qui laisse une impression forte au recruteur.}
\end{popfigure}

\begin{exemplebox}[Un paragraphe de corps rédigé et annoté]
Candidature à un poste de \textbf{secrétaire bilingue} dans une entreprise de Buea :
\begin{quote}
« Secrétaire stagiaire pendant trois mois à la mairie de Buea, j'y ai accueilli les usagers en français et en anglais, classé plus de deux cents dossiers et rédigé des comptes rendus de réunion. Peu expérimentée hier, \textbf{mais} déjà opérationnelle aujourd'hui, je suis devenue une collaboratrice \textbf{ponctuelle, organisée, fiable, irremplaçable}. Rejoindre vos équipes, c'est pour moi la possibilité de mettre cette rigueur au service d'une entreprise reconnue dans toute la région du Sud-Ouest pour la qualité de son accueil client. »
\end{quote}
Analyse : l'\textbf{antithèse} « peu expérimentée hier, mais déjà opérationnelle aujourd'hui » transforme une faiblesse (la jeunesse) en atout ; la \textbf{gradation} « ponctuelle, organisée, fiable, irremplaçable » monte en puissance ; la \textbf{métonymie} « rejoindre vos équipes » dit l'entreprise autrement ; et la dernière phrase nomme l'entreprise et sa réputation : la lettre est personnalisée.
\end{exemplebox}

\courssub{Le texte descriptif : rendre le parcours concret et crédible}

Pour prouver une expérience ou une qualité, il faut parfois \textbf{décrire} brièvement. Le \cle{texte descriptif} s'appuie sur des détails concrets : nombres, lieux, objets, gestes. Ces détails créent la crédibilité : on croit ce qu'on peut voir.

\begin{exemplebox}[Décrire une réalisation en quelques lignes]
Version vague (non crédible) : « j'ai participé à un projet d'électricité à l'école. »\\
Version descriptive (crédible) : « Avec quatre camarades, j'ai câblé le tableau électrique de la salle informatique du lycée : quinze prises, deux disjoncteurs et un interrupteur différentiel, installés en trois semaines sous la supervision de notre professeur. Depuis, la salle fonctionne sans panne. »
\end{exemplebox}

La bonne description professionnelle est \textbf{brève} (deux ou trois phrases), \textbf{précise} (chiffres, lieux, durées) et \textbf{orientée} vers la compétence qu'elle doit prouver. Elle ne raconte pas une histoire pour le plaisir : elle fournit une preuve.

\courssub{La communication par l'image : ce qu'une affiche nous apprend de la persuasion}

Avant de rédiger, observons comment un message visuel persuade. Prenons une \textbf{affiche d'indication sécuritaire} que l'on voit dans les ateliers et les chantiers du Cameroun : un casque jaune sur fond bleu, un pictogramme d'ouvrier, et le slogan « Port du casque obligatoire ». Pourquoi cette affiche fonctionne-t-elle ?

\begin{experience}[Analyser une affiche sécuritaire comme un message persuasif]
\begin{enumerate}
\item \textbf{Observer} : l'affiche combine une image (le casque, l'ouvrier), un texte bref (la consigne) et souvent une couleur d'alerte (jaune, rouge).
\item \textbf{Dégager le procédé} : l'image montre immédiatement \emph{de quoi} il s'agit ; le slogan dit \emph{ce qu'il faut faire} ; la couleur crée l'\emph{émotion} (danger, urgence).
\item \textbf{Conclure} : un message persuasif efficace associe \textbf{preuve visible + consigne claire + impact émotionnel}.
\item \textbf{Transposer à l'écrit} : dans votre lettre, la « preuve visible » est le fait concret (chiffres, dates), la « consigne claire » est votre demande (le poste, l'entretien), et « l'impact émotionnel » est créé par les figures de style (gradation, antithèse).
\end{enumerate}
\end{experience}

\begin{popfigure}
\begin{center}
\begin{tikzpicture}[
  zone/.style={draw=popdark, thick, rounded corners, align=center, font=\small},
  effet/.style={draw=popgold, fill=popgoldL, rounded corners, align=center, font=\small\itshape, text width=3.2cm}
]
\node[zone, fill=popblueL, text width=3.2cm, minimum height=2.2cm] (img) at (0,0) {\textbf{IMAGE}\\[4pt] casque jaune\\ + ouvrier};
\fill[popgold] (-0.55,0.75) arc (180:0:0.55) -- cycle;
\draw[popdark, thick] (-0.55,0.75) arc (180:0:0.55) -- cycle;
\draw[popdark, thick] (-0.75,0.75) -- (0.75,0.75);
\node[zone, fill=popgreenL, text width=3.2cm, minimum height=2.2cm] (txt) at (4.6,0) {\textbf{TEXTE / SLOGAN}\\[4pt] « Port du casque\\ obligatoire »};
\node[zone, fill=poppinkL, text width=3.2cm, minimum height=2.2cm] (emo) at (9.2,0) {\textbf{COULEUR}\\[4pt] jaune et rouge :\\ alerte, urgence};
\node[effet] (eff) at (4.6,-2.6) {\textbf{EFFET} : le lecteur comprend vite, croit et obéit};
\draw[-{Stealth}, thick, poporange] (img.south) -- (eff.north west);
\draw[-{Stealth}, thick, poporange] (txt.south) -- (eff.north);
\draw[-{Stealth}, thick, poporange] (emo.south) -- (eff.north east);
\end{tikzpicture}
\end{center}
\legende{Analyse d'une affiche sécuritaire : image + slogan + couleur convergent vers un seul effet persuasif. Dans la lettre de motivation, les faits concrets, la demande claire et les figures de style jouent les mêmes rôles.}
\end{popfigure}

\courssub{Le registre soutenu : la crédibilité du candidat}

Une lettre de motivation se rédige en \cle{registre soutenu} : vocabulaire précis (« collaborer » plutôt que « bosser »), phrases claires et bien construites, formules de politesse, et surtout \textbf{absence totale de fautes}. Une faute d'orthographe dans une lettre de motivation est comme une tache sur une chemise le jour de l'entretien : elle détruit l'image de rigueur que tout le texte cherche à construire.

\begin{attention}[Trois fautes qui coûtent un entretien]
\begin{enumerate}
\item \textbf{Répéter le CV phrase par phrase} : la lettre \emph{sélectionne} et \emph{commente}, elle ne duplique pas. Le CV dit « stage à la SABC, 2025 » ; la lettre explique ce que ce stage vous a appris.
\item \textbf{Les formules vagues} : « je suis dynamique et motivé » sans aucune preuve ; « votre honorable entreprise » sans nom ni activité précise.
\item \textbf{Les fautes de langue} : confusion « a / à », « ce / se », accords négligés. Relisez toujours votre lettre à voix haute, puis faites-la relire.
\end{enumerate}
\end{attention}

\begin{exoresolu}[Rédiger le corps d'une lettre de motivation]
Énoncé : À partir de la situation suivante, rédigez les trois paragraphes du corps d'une lettre de motivation. Situation : Aminatou, élève en Première technique (spécialité secrétariat-bureautique) au lycée technique de Bafoussam, candidate à un poste de secrétaire bilingue chez « Highland Agro Services », entreprise de Buea spécialisée dans l'exportation du café. Éléments du parcours : stage de deux mois à la mairie de Bafoussam (accueil bilingue, classement de 150 dossiers) ; trésorière du club d'anglais ; très bon niveau en dactylographie (45 mots par minute).
\tcblower
Solution détaillée :
\begin{quote}
« Actuellement en classe de Première technique, spécialité secrétariat-bureautique, au lycée technique de Bafoussam, je maîtrise la dactylographie à raison de quarante-cinq mots par minute, ainsi que les outils de bureautique indispensables au poste de secrétaire bilingue que vous proposez.

Ensuite, mon expérience confirme cette formation : pendant deux mois, lors d'un stage à la mairie de Bafoussam, j'ai accueilli les usagers en français et en anglais et classé cent cinquante dossiers sans aucune erreur. En outre, trésorière du club d'anglais de mon lycée, j'ai géré le budget du club pendant un an avec une rigueur reconnue par tous.

Jeune candidate, mais déjà expérimentée, je suis devenue une collaboratrice ponctuelle, organisée, fiable, irremplaçable. C'est pourquoi rejoindre vos équipes chez Highland Agro Services, entreprise reconnue dans toute la région du Sud-Ouest pour la qualité de son café d'exportation, serait pour moi l'occasion de mettre ces qualités au service d'une structure exigeante. »
\end{quote}
Analyse : paragraphe 1 = formation liée au poste ; paragraphe 2 = deux expériences situées, décrites (chiffres : deux mois, cent cinquante dossiers) et valorisées ; paragraphe 3 = antithèse (« jeune, mais déjà expérimentée »), gradation (« ponctuelle, organisée, fiable, irremplaçable »), métonymie (« rejoindre vos équipes »), connecteurs (« ensuite », « en outre », « c'est pourquoi ») et personnalisation (nom et activité de l'entreprise).
\end{exoresolu}

\begin{exercice}[À vous de convaincre]
Situation : M. Etoundi, élève de Première technique en génie mécanique à Ebolowa, postule à un stage ouvrier dans un garage de Yaoundé spécialisé dans les poids lourds. Il a réparé le moteur du car de son lycée avec deux camarades et il est toujours ponctuel à l'atelier. Rédigez les trois paragraphes du corps de sa lettre en utilisant : au moins trois connecteurs logiques, une gradation ascendante, une antithèse et une métonymie. Soulignez chaque procédé utilisé.
\end{exercice}

\begin{corrige}[Exercice n]
Éléments de correction attendus :
\begin{itemize}
\item \textbf{Paragraphe 1 (formation)} : mention de la classe et de la spécialité en génie mécanique, compétences liées au poste (diagnostic, démontage, remontage moteur).
\item \textbf{Paragraphe 2 (expérience)} : la réparation du moteur du car présentée en trois temps (situation, tâche, acquis), avec détails concrets ; connecteur « ensuite » ou « en outre ».
\item \textbf{Paragraphe 3 (qualités et motivation)} : ponctualité prouvée (« aucun retard à l'atelier depuis la rentrée ») ; antithèse du type « élève encore, mais déjà mécanicien dans les gestes » ; gradation du type « appliqué, rigoureux, passionné, indispensable » ; métonymie du type « mettre mes mains au service de votre atelier » ; nom et spécialité du garage mentionnés ; connecteur de conclusion « c'est pourquoi ».
\end{itemize}
Barème indicatif : respect des trois paragraphes (4 points), connecteurs (4 points), figures de style (6 points), preuves concrètes (4 points), langue correcte (2 points).
\end{corrige}

\begin{aretenir}[L'essentiel du corps de la lettre]
\begin{itemize}
\item Trois paragraphes : \textbf{formation} liée au poste, \textbf{expériences} prouvées par des faits, \textbf{qualités} démontrées et \textbf{motivation} personnalisée.
\item Des \textbf{connecteurs} pour guider le lecteur : d'abord, ensuite, en outre, c'est pourquoi.
\item Trois figures de style pour persuader : \textbf{gradation} (montée d'intensité), \textbf{antithèse} (opposition valorisante), \textbf{métonymie} (dire autrement).
\item Des \textbf{détails concrets} (chiffres, lieux, durées) qui rendent le parcours crédible, comme l'image d'une affiche rend une consigne crédible.
\item Un \textbf{registre soutenu}, sans fautes : la forme est la première preuve de votre sérieux.
\end{itemize}
\end{aretenir}

\coursec{Production guidée : rédiger une lettre de motivation complète}

Dans les sections précédentes, nous avons appris à construire le corps de la lettre de motivation : trois paragraphes organisés selon le plan \emph{Vous -- Moi -- Nous}, enrichis de figures de style (gradation, antithèse, métonymie) et appuyés sur des preuves concrètes. Mais une belle lettre ne naît jamais du premier jet. Comme un technicien qui ne monte pas un système sans avoir d'abord vérifié chaque composant, le candidat sérieux suit une \cle{démarche de production} en étapes. C'est cette démarche complète — de la préparation à la version finale — que nous allons maintenant pratiquer, car c'est exactement celle que l'examinateur du Probatoire technique attend de vous le jour de l'épreuve de production d'écrits.

\begin{definition}[Lettre de motivation]
La \cle{lettre de motivation} est une lettre officielle, adressée à un employeur ou un responsable de formation, dans laquelle le candidat expose les raisons de sa candidature, met en valeur ses compétences et son adéquation avec le poste ou la formation visés, et sollicite un entretien. Elle respecte une \cle{structure externe} (mise en page codifiée) et une \cle{structure interne} (organisation des idées en paragraphes).
\end{definition}

\begin{propriete}[Les qualités d'une bonne lettre de motivation]
Une bonne lettre de motivation tient en \cle{une page}, comporte \cle{trois à quatre paragraphes} (introduction, un ou deux paragraphes de corps, conclusion) et répond à une seule question, celle que se pose le recruteur : \emph{« pourquoi vous et pas un autre ? »}. Toute phrase qui n'aide pas à répondre à cette question doit être supprimée.
\end{propriete}

\courssub{Vue d'ensemble : les cinq étapes de la production}

Rédiger une lettre de motivation, c'est suivre un parcours en cinq étapes : \cle{préparer}, \cle{rédiger un brouillon}, \cle{réviser}, \cle{améliorer} grâce au regard des autres, puis \cle{finaliser}. Aucune étape ne peut être sautée : celui qui écrit directement « au propre » produit presque toujours une lettre faible, car il n'a ni analysé l'offre ni relu son texte.

\begin{popfigure}
\begin{center}
\begin{tikzpicture}[
  node distance=0.55cm,
  etape/.style={rectangle, rounded corners=3pt, draw=popblue, fill=popblueL, thick, text width=3.1cm, minimum height=1.15cm, align=center, font=\small},
  desc/.style={rectangle, draw=popline, fill=popcream, text width=3.1cm, align=center, font=\scriptsize, below=0.12cm of \tikzlastnode},
  fleche/.style={-{Stealth[length=3mm]}, very thick, poporange}
]
\node[etape, align=center] (e1){\textbf{1. Préparer}\\ analyser l'offre};
\node[etape, right=0.7cm of e1, align=center] (e2){\textbf{2. Brouillon}\\ rédiger librement};
\node[etape, right=0.7cm of e2, align=center] (e3){\textbf{3. Réviser}\\ s'auto-évaluer};
\node[etape, below=1.1cm of e3, align=center] (e4){\textbf{4. Améliorer}\\ échange entre pairs};
\node[etape, left=0.7cm of e4, align=center] (e5){\textbf{5. Finaliser}\\ copie propre};
\draw[fleche] (e1) -- (e2);
\draw[fleche] (e2) -- (e3);
\draw[fleche] (e3) -- (e4);
\draw[fleche] (e4) -- (e5);
\draw[fleche, dashed, poppurple] (e5.west) .. controls +(-1.4,0.9) and +(-1.4,-0.9) .. (e1.west) node[midway, left, font=\scriptsize\itshape, text=poppurple, align=center]{si\\ nécessaire};
\end{tikzpicture}
\end{center}
\legende{Organigramme des cinq étapes de production d'une lettre de motivation. La flèche en pointillés rappelle qu'on peut revenir à la préparation si la révision révèle un problème de fond.}
\end{popfigure}

\courssub{Étape 1 — Préparation : analyser avant d'écrire}

\textbf{L'intuition.} Imaginez un menuisier qui fabriquerait une porte sans avoir mesuré l'ouverture : la porte sera belle, peut-être, mais elle ne fermera pas. De même, une lettre écrite sans analyse préalable de l'offre sera bien rédigée mais \emph{à côté du sujet}. La préparation est l'étape qui garantit la \cle{pertinence} de votre candidature.

\textbf{Le travail à mener} comporte trois opérations :

\begin{enumerate}
\item \textbf{Analyser l'offre d'emploi ou le sujet d'examen.} On relève : l'intitulé exact du \cle{poste}, l'entreprise ou l'administration, les \cle{qualités requises} (ponctualité, sens du travail en équipe, maîtrise d'un logiciel, permis de conduire...), les diplômes exigés et le \cle{destinataire} (Monsieur le Directeur des Ressources Humaines, Madame la Cheffe d'agence...).
\item \textbf{Faire le bilan de son propre profil} (brainstorming individuel) : diplômes, stages, expériences, qualités personnelles, et surtout les \cle{preuves} qui les illustrent.
\item \textbf{Construire le plan « Vous -- Moi -- Nous »} : ce que je sais de l'entreprise (Vous), ce que je peux lui apporter (Moi), ce que nous pouvons accomplir ensemble (Nous).
\end{enumerate}

\begin{methode}[La fiche de préparation en quatre questions]
Avant d'écrire la moindre phrase, répondez par écrit à quatre questions :
\begin{enumerate}
\item \textbf{À qui j'écris ?} — identité et fonction exactes du destinataire, pour choisir la formule d'appel et la formule de politesse adaptées.
\item \textbf{Quel poste ?} — intitulé précis, référence de l'annonce éventuelle, pour rédiger l'objet de la lettre.
\item \textbf{Quelles qualités sont demandées ?} — relevez-les dans l'offre : ce sont elles que votre lettre devra démontrer.
\item \textbf{Quelles preuves ai-je ?} — pour chaque qualité demandée, notez un fait précis (stage, résultat, responsabilité exercée) qui prouve que vous la possédez.
\end{enumerate}
Si vous ne pouvez pas répondre à l'une de ces questions, votre préparation n'est pas terminée : retournez à l'offre.
\end{methode}

\begin{exemplebox}[Fiche de préparation remplie]
Offre : \emph{« Entreprise SOCATRAV de Douala recrute un technicien en génie électrique. Profil : probatoire technique F3 minimum, rigueur, capacité à travailler en équipe, disponibilité immédiate. Écrire au Directeur des Ressources Humaines, BP 4521 Douala. »}

\begin{center}
\begin{tabularx}{\linewidth}{|l|X|}\hline
\rowcolor{popblueL}\textbf{Question} & \textbf{Réponse du candidat} \\ \hline
À qui j'écris ? & Monsieur le Directeur des Ressources Humaines de SOCATRAV, BP 4521 Douala. \\ \hline
Quel poste ? & Technicien en génie électrique ; objet : « Candidature au poste de technicien en génie électrique ». \\ \hline
Qualités demandées ? & Rigueur, travail en équipe, disponibilité, niveau probatoire technique F3. \\ \hline
Quelles preuves ? & Probatoire F3 obtenu en 2025 ; stage de 3 mois à ENEO Bonabéri (maintenance des installations) ; chef d'équipe du projet d'éclairage du lycée ; disponible dès le 1\textsuperscript{er} septembre. \\ \hline
\end{tabularx}
\end{center}
\end{exemplebox}

\courssub{Étape 2 — Le brouillon : écrire sans se censurer}

\textbf{L'intuition.} Au brouillon, l'objectif n'est pas la perfection mais la \emph{matière} : comme le sculpteur dégage d'abord un bloc avant de polir, vous couchez d'abord toutes vos idées sur le papier. Les corrections viendront après.

\textbf{La démarche} suit l'ordre de la lettre elle-même :

\begin{enumerate}
\item \textbf{La structure externe complète} : en-tête (coordonnées de l'expéditeur en haut à gauche, celles du destinataire à droite), lieu et date (« Douala, le 15 octobre 2026 »), objet de la lettre, formule d'appel (« Monsieur le Directeur, »).
\item \textbf{L'introduction} : deux à trois phrases qui annoncent la candidature, précisent le poste et donnent envie de lire la suite (l'accroche).
\item \textbf{Le corps en trois paragraphes} : \emph{Vous} (montrer qu'on connaît l'entreprise), \emph{Moi} (présenter diplômes, expériences et qualités avec leurs preuves), \emph{Nous} (projeter la collaboration, proposer un entretien).
\item \textbf{La conclusion} : formule de politesse adaptée au destinataire, signature.
\end{enumerate}

\begin{attention}[Le piège du brouillon « définitif »]
Beaucoup d'élèves rédigent directement au propre, par peur de manquer de temps. C'est une erreur stratégique : sans brouillon, aucune relecture véritable n'est possible, et \textbf{une seule faute d'accord peut discréditer toute la candidature}. Un recruteur qui lit « je suis \emph{motivée} » sous la plume d'un candidat masculin, ou « vos \emph{entreprise} », conclut immédiatement à la négligence — défaut rédhibitoire pour un poste exigeant de la rigueur. Au Probatoire, la correction de la langue est un critère d'évaluation à part entière : prévoyez toujours dix minutes de relecture.
\end{attention}

\courssub{Étape 3 — La révision : s'auto-évaluer avec une grille}

\textbf{L'intuition.} Relire sa lettre, ce n'est pas la parcourir des yeux : c'est la soumettre à un \emph{contrôle qualité}, comme on vérifie un appareil avant livraison. Pour être efficace, ce contrôle suit une grille précise, point par point.

\begin{methode}[Grille d'auto-évaluation en cinq points]
Relisez votre brouillon en cochant chaque critère :
\begin{enumerate}
\item \textbf{Structure externe} : en-tête complet, date, objet, appel, politesse, signature — rien ne manque ?
\item \textbf{Pertinence} : chaque qualité demandée dans l'offre est-elle traitée ? Chaque affirmation est-elle prouvée par un fait ?
\item \textbf{Cohérence} : le plan Vous -- Moi -- Nous est-il respecté ? Les paragraphes sont-ils liés par des connecteurs (d'abord, en outre, enfin) ?
\item \textbf{Langue et style} : au moins une figure de style étudiée (gradation, antithèse, métonymie) ? Phrases de longueur variée ? Registre soutenu ?
\item \textbf{Orthographe et grammaire} : accords sujet-verbe, accords des participes passés, accords en genre et en nombre, ponctuation.
\end{enumerate}
Tout critère non satisfait donne lieu à une correction immédiate sur le brouillon.
\end{methode}

\courssub{Étape 4 — L'amélioration : le regard des autres}

\textbf{L'intuition.} On est toujours mauvais juge de son propre texte : on relit ce qu'on a \emph{voulu} écrire, pas ce qu'on a écrit. Le regard d'un camarade, lui, est neuf.

\textbf{La démarche en classe} se déroule en deux temps :

\begin{enumerate}
\item \textbf{L'échange entre pairs} : par binômes, les élèves échangent leurs brouillons. Chaque relecteur applique la grille d'auto-évaluation à la lettre de son camarade et formule \emph{deux points forts} et \emph{deux axes d'amélioration}, par écrit, avec courtoisie (« ton introduction est accrocheuse, mais le paragraphe "Moi" manque de preuves »).
\item \textbf{La correction collective} : l'enseignant projette ou recopie au tableau une lettre \emph{anonyme} (volontairement imparfaite). La classe la corrige phrase par phrase : on barre, on remplace, on justifie chaque modification en citant la règle appliquée.
\end{enumerate}

\begin{exoresolu}[Corriger un extrait de lettre anonyme]
Énoncé. Voici le début du corps d'une lettre anonyme écrite au tableau. Relevez les défauts et proposez une version améliorée : \emph{« Je suis très bon en électricité. J'ai fait un stage. Votre entreprise est grande et connue. Je veux le poste parce que j'ai besoin d'argent. »}
\tcblower
Solution détaillée. Défauts relevés : (1) « très bon » est une affirmation sans preuve ; (2) « j'ai fait un stage » est vague : où, quand, quelles tâches ? (3) « grande et connue » montre qu'on ne connaît pas réellement l'entreprise ; (4) « j'ai besoin d'argent » répond à l'envers : le recruteur demande « pourquoi vous ? », pas « pourquoi avez-vous besoin de nous ? » ; (5) phrases courtes et sèches, aucune figure de style, registre familier (« très bon »).

Version améliorée : \emph{« Entreprise de référence dans le BTP camerounais, SOCATRAV mène des chantiers dont l'exigence technique correspond exactement à ma formation. Titulaire du probatoire technique F3, j'ai consolidé mes compétences lors d'un stage de trois mois à ENEO Bonabéri, où j'ai participé à la maintenance des installations basse tension : rigueur dans les diagnostics, rapidité dans les interventions, esprit d'équipe dans chaque mission. Intégrer vos équipes, c'est pour moi l'occasion de mettre cette expérience au service de vos chantiers et de grandir à vos côtés. »} On remarque la gradation (« rigueur, rapidité, esprit d'équipe »), le respect du plan Vous -- Moi -- Nous et des preuves précises (durée, lieu, tâches).
\end{exoresolu}

\courssub{Étape 5 — La version finale : la copie propre}

La version finale est la lettre telle que le destinataire la recevra. Elle doit être :

\begin{itemize}
\item \textbf{Propre et lisible} : manuscrite (encre bleue ou noire, écriture soignée, aucune rature) ou tapée à l'ordinateur (police classique, interligne aéré, marges régulières) ;
\item \textbf{Limitée à une page} : au-delà, le recruteur décroche ;
\item \textbf{Complète} : la mention des \cle{pièces jointes} figure en bas à gauche, avant ou après la signature : « P. J. : curriculum vitae, photocopie du diplôme, attestation de stage » ;
\item \textbf{Signée} : la signature manuscrite engage le candidat, même sur une lettre tapée.
\end{itemize}

\begin{exemplebox}[Du brouillon à la version finale : annotations]
Extrait du brouillon de Marie, élève de Première F3, avec les annotations de la correction collective et la version finale correspondante.

\begin{center}
\begin{tabularx}{\linewidth}{|X|X|}\hline
\rowcolor{popblueL}\textbf{Brouillon annoté} & \textbf{Version finale} \\ \hline
« Je suis \emph{motivée} » $\rightarrow$ \emph{accord fautif : Marie signe "M. Ngo Bell", vérifier ; ici on reformule pour éviter le piège} & « Ma motivation est entière. » \\ \hline
« J'ai travaillé durant mes vacances » $\rightarrow$ \emph{trop vague : où ? quoi ?} & « Pendant les grandes vacances 2025, j'ai assisté l'équipe de câblage de l'entreprise ELEC+ à Akwa. » \\ \hline
« Je serais content de venir » $\rightarrow$ \emph{registre familier} & « Je me tiens à votre disposition pour un entretien. » \\ \hline
\emph{Pas de pièces jointes mentionnées} & « P. J. : CV, photocopie du probatoire, attestation de stage. » \\ \hline
\end{tabularx}
\end{center}
\end{exemplebox}

\courssub{Évaluation de la production : la grille de notation}

La lettre finale est évaluée sur vingt points, selon une grille qui reprend exactement les critères du Probatoire technique.

\begin{popfigure}
\begin{center}
\begin{tikzpicture}[x=1cm, y=1cm]
\draw[fill=popblue, draw=popblue] (0,0) rectangle (10.5,0.9);
\node[text=white, font=\small\bfseries, anchor=west] at (0.25,0.45) {Critère d'évaluation};
\node[text=white, font=\small\bfseries, anchor=east] at (10.25,0.45) {Barème};
\draw[fill=popcream, draw=popline] (0,-0.9) rectangle (10.5,0);
\node[font=\small, anchor=west] at (0.25,-0.45) {\textbf{Structure externe et interne} : en-tête, objet, appel, plan, politesse};
\node[font=\small\bfseries, text=popblue, anchor=east] at (10.25,-0.45) {/3};
\draw[fill=white, draw=popline] (0,-1.8) rectangle (10.5,-0.9);
\node[font=\small, anchor=west] at (0.25,-1.35) {\textbf{Contenu} : pertinence, preuves, réponse à « pourquoi vous ? »};
\node[font=\small\bfseries, text=popblue, anchor=east] at (10.25,-1.35) {/6};
\draw[fill=popcream, draw=popline] (0,-2.7) rectangle (10.5,-1.8);
\node[font=\small, anchor=west] at (0.25,-2.25) {\textbf{Langue} : orthographe, syntaxe, registre, figures de style};
\node[font=\small\bfseries, text=popblue, anchor=east] at (10.25,-2.25) {/6};
\draw[fill=white, draw=popline] (0,-3.6) rectangle (10.5,-2.7);
\node[font=\small, anchor=west] at (0.25,-3.15) {\textbf{Présentation} : propreté, lisibilité, une page, pièces jointes};
\node[font=\small\bfseries, text=popblue, anchor=east] at (10.25,-3.15) {/5};
\draw[fill=poporangeL, draw=poporange, thick] (0,-4.5) rectangle (10.5,-3.6);
\node[font=\small\bfseries, anchor=west] at (0.25,-4.05) {TOTAL};
\node[font=\small\bfseries, text=poporange, anchor=east] at (10.25,-4.05) {/20};
\end{tikzpicture}
\end{center}
\legende{Grille d'évaluation de la lettre de motivation : structure (3 points), contenu (6 points), langue (6 points), présentation (5 points), pour un total de 20 points.}
\end{popfigure}

\begin{savaistu}[Les critères du Probatoire technique]
Au Probatoire technique, la production écrite est évaluée selon quatre grandes familles de critères : la \emph{pertinence} (le texte répond-il à la situation de communication ?), la \emph{cohérence} (les idées sont-elles organisées et liées ?), la \emph{correction de la langue} (orthographe, grammaire, vocabulaire) et la \emph{présentation} (lisibilité, respect de la forme du type de texte demandé). La grille ci-dessus est donc un entraînement direct à l'examen : chaque lettre rédigée en classe est une répétition de l'épreuve réelle.
\end{savaistu}

\courssub{Justifier sa démarche : l'oral de métalangage}

Dernière étape, souvent négligée et pourtant décisive pour apprendre : \textbf{expliquer ses choix}. Chaque élève présente oralement à la classe \emph{deux choix de rédaction} effectués dans sa lettre et les justifie. Trois types de choix peuvent être défendus :

\begin{itemize}
\item \textbf{Le choix de l'accroche} : « J'ai ouvert par une phrase sur le chantier récent de SOCATRAV pour montrer que je connais l'entreprise, car le plan Vous -- Moi -- Nous commence par s'intéresser au destinataire. »
\item \textbf{Le choix d'une figure de style} : « J'ai utilisé une gradation — rigueur, rapidité, esprit d'équipe — pour énumérer mes qualités dans un ordre croissant d'importance et frapper la mémoire du lecteur. »
\item \textbf{Le choix d'une preuve} : « J'ai cité mon stage à ENEO plutôt que mes notes, parce qu'une expérience réelle prouve mieux qu'un résultat scolaire que je maîtrise le métier. »
\end{itemize}

Cet exercice de justification développe le savoir-faire essentiel du programme : \emph{rédiger et justifier une démarche, une réponse ou une conclusion}. Celui qui sait dire \emph{pourquoi} il a écrit ce qu'il a écrit ne rédige plus au hasard : il écrit en professionnel.

\begin{exercice}[Production complète en situation]
Situation d'intégration : \emph{L'entreprise camerounaise AGROPLUS, spécialisée dans la transformation du manioc à Bafoussam, recrute un assistant de production. Profil exigé : probatoire technique (série F ou G), sens de l'organisation, goût du travail en équipe. Candidature à adresser à Monsieur le Directeur Général, BP 210 Bafoussam.}

Travail à faire, en suivant les cinq étapes de la leçon : (1) remplissez la fiche de préparation en quatre questions ; (2) rédigez le brouillon complet de votre lettre (structure externe, introduction, corps en trois paragraphes, conclusion) ; (3) auto-évaluez-le avec la grille en cinq points ; (4) échangez avec un camarade et améliorez ; (5) copiez la version finale et préparez la justification orale de deux de vos choix de rédaction.
\end{exercice}

\begin{corrige}[Exercice — éléments de réussite]
Une copie réussie présente : une structure externe complète et sans erreur (coordonnées, date, objet précis, « Monsieur le Directeur Général ») ; une introduction qui annonce la candidature avec une accroche liée à AGROPLUS ; un corps Vous -- Moi -- Nous où chaque qualité exigée (organisation, travail en équipe) est prouvée par un fait concret (stage, responsabilité associative, projet scolaire) ; au moins une figure de style pertinente ; une formule de politesse adaptée (« Veuillez agréer, Monsieur le Directeur Général, l'expression de ma haute considération ») ; la mention des pièces jointes ; une page maximum, sans faute d'accord. L'oral de justification doit citer la règle ou la stratégie qui motive chaque choix.
\end{corrige}

\begin{aretenir}[L'essentiel de la production guidée]
Produire une lettre de motivation, c'est suivre cinq étapes indissociables : \textbf{préparer} (fiche en quatre questions : à qui ? quel poste ? quelles qualités ? quelles preuves ?), \textbf{rédiger un brouillon} (structure externe, introduction, corps Vous -- Moi -- Nous, conclusion), \textbf{réviser} (grille d'auto-évaluation en cinq points), \textbf{améliorer} (échange entre pairs et correction collective), \textbf{finaliser} (copie propre d'une page, pièces jointes mentionnées, signature). La lettre est évaluée sur 20 points : structure /3, contenu /6, langue /6, présentation /5. Enfin, savoir justifier oralement ses choix de rédaction transforme un exercice scolaire en compétence professionnelle durable.
\end{aretenir}

\coursec{Compte-rendu, évaluation et remédiation}

Après avoir rédigé, à la section précédente, une lettre de motivation complète en production guidée, chaque élève dispose désormais d'un texte personnel. Mais une compétence n'est réellement acquise que si l'on sait dire ce que l'on a appris, mesurer ce qui reste fragile et corriger ses propres erreurs. C'est précisément l'objet de cette dernière étape de la séquence : faire le \cle{compte-rendu} des acquis, vérifier la maîtrise de la compétence par une \cle{évaluation} individuelle, puis organiser une \cle{remédiation} ciblée pour consolider ce qui n'est pas encore solide. Enfin, nous réinvestirons la compétence dans d'autres situations de la vie courante, car une lettre officielle ne sert pas seulement à chercher un emploi.

\courssub{Compte-rendu de la séquence : synthèse collective des acquis}

Le compte-rendu est un retour collectif et organisé sur tout le travail de la séquence. Il se construit en classe entière, à partir des productions des élèves et des notes prises depuis la découverte des textes iconiques (lettres de motivation authentiques et affiches d'indications sécuritaires). L'objectif est de transformer des connaissances dispersées en un savoir stable, formulé clairement et mémorisable.

\begin{definition}[Compte-rendu de séquence]
Le \cle{compte-rendu} est une synthèse écrite et orale qui récapitule, à la fin d'une séquence, les notions étudiées, les méthodes acquises et les difficultés rencontrées. Il est élaboré collectivement, validé par l'enseignant, puis consigné par chaque élève dans son cahier de cours.
\end{definition}

La démarche se déroule en trois temps. D'abord, chaque élève relit sa lettre produite en section 5 et note sur une feuille : ce qu'il a réussi, ce qui lui a paru difficile, ce qu'il retient de la séquence (structure, style, figures de style, communication par l'image). Ensuite, par groupes de quatre, les élèves confrontent leurs bilans et formulent ensemble les acquis essentiels. Enfin, la classe entière met en commun et le professeur guide la formulation d'une synthèse unique, qui donnera naissance à la fiche-bilan affichée au mur.

La synthèse collective fait ressortir quatre piliers :

\begin{itemize}
\item \textbf{La forme} : la lettre de motivation possède une \cle{structure externe} codifiée, composée de huit éléments obligatoires (coordonnées de l'expéditeur, coordonnées du destinataire, lieu et date, objet, formule d'appel, corps de la lettre, formule de politesse, signature).
\item \textbf{Le fond} : la \cle{structure interne} suit le schéma argumentatif « Vous – Moi – Nous » : parler de l'entreprise ou de l'institution, parler de soi et de ses compétences, puis projeter la collaboration.
\item \textbf{Le style} : le candidat convainc grâce à un vocabulaire précis, des connecteurs logiques et des figures de style maîtrisées (gradation, antithèse, métonymie).
\item \textbf{La méthode} : on ne rédige jamais directement au propre ; on suit les étapes de production : analyse de l'annonce, plan, brouillon, rédaction, relecture, amélioration.
\end{itemize}

\begin{popfigure}
\begin{tikzpicture}[mindmap, grow cyclic, every node/.style={concept, text=white}, root concept/.append style={concept color=popblue, font=\large\bfseries}, level 1/.append style={level distance=3.6cm, sibling angle=90}]
\node[root concept]{La lettre de motivation}
  child[concept color=poporange]{ node {La forme} 
    child[concept color=poporange!70, level distance=2.6cm]{ node[concept, font=\small] {8 éléments externes} }
    child[concept color=poporange!70, level distance=2.6cm]{ node[concept, font=\small] {Objet précis} }
  }
  child[concept color=popgreen]{ node {Le fond}
    child[concept color=popgreen!70, level distance=2.6cm]{ node[concept, font=\small] {Vous} }
    child[concept color=popgreen!70, level distance=2.6cm]{ node[concept, font=\small] {Moi} }
    child[concept color=popgreen!70, level distance=2.6cm]{ node[concept, font=\small] {Nous} }
  }
  child[concept color=poppurple]{ node {Le style}
    child[concept color=poppurple!70, level distance=2.6cm]{ node[concept, font=\small] {Gradation} }
    child[concept color=poppurple!70, level distance=2.6cm]{ node[concept, font=\small] {Antithèse} }
    child[concept color=poppurple!70, level distance=2.6cm]{ node[concept, font=\small] {Métonymie} }
  }
  child[concept color=popgold]{ node {La méthode}
    child[concept color=popgold!70, level distance=2.6cm]{ node[concept, font=\small] {Analyser l'annonce} }
    child[concept color=popgold!70, level distance=2.6cm]{ node[concept, font=\small] {Brouillon \& relecture} }
  };
\end{tikzpicture}
\legende{Carte mentale de synthèse de la séquence : la lettre de motivation repose sur quatre piliers complémentaires — la forme (structure externe), le fond (schéma « Vous – Moi – Nous »), le style (figures et connecteurs) et la méthode (étapes de production).}
\end{popfigure}

\begin{aretenir}[Les acquis essentiels de la séquence]
\begin{enumerate}
\item Une lettre de motivation comporte \textbf{8 éléments externes} placés à des positions fixes.
\item Son plan interne suit le schéma \textbf{Vous – Moi – Nous}.
\item Elle convainc par des \textbf{qualités prouvées} (exemples, résultats) et un style soigné.
\item Elle se construit en \textbf{étapes} : annonce analysée, plan, brouillon, relecture.
\item La réception de textes iconiques (lettres, affiches sécuritaires) développe la compétence \emph{communication par l'image}.
\end{enumerate}
\end{aretenir}

\courssub{La fiche-bilan affichée en classe}

Pour rendre la synthèse visible et durable, la classe élabore ensemble une \cle{fiche-bilan} qui sera affichée au mur de la salle. Cette fiche récapitule les trois outils majeurs de la séquence.

\textbf{1. Les huit éléments de la structure externe.}

\begin{center}
\begin{tabularx}{\linewidth}{|c|l|X|}
\hline
\rowcolor{popblueL} \textbf{N°} & \textbf{Élément} & \textbf{Position et règle} \\
\hline
1 & Coordonnées de l'expéditeur & En haut à gauche : nom, prénom, adresse, téléphone, courriel. \\
\hline
2 & Coordonnées du destinataire & En haut à droite : fonction, entreprise/organisation, ville. \\
\hline
3 & Lieu et date & À droite, sous le destinataire : « Douala, le 15 mars 2026 ». \\
\hline
4 & Objet & Précis et court : « Objet : candidature au poste de technicien comptable ». \\
\hline
5 & Formule d'appel & « Monsieur le Directeur, », « Madame, » — jamais de nom propre. \\
\hline
6 & Corps de la lettre & Introduction, développement (Vous – Moi – Nous), conclusion. \\
\hline
7 & Formule de politesse & Soutenue et complète : « Veuillez agréer, Monsieur le Directeur, l'expression de mes salutations distinguées. » \\
\hline
8 & Signature & Nom, prénom et paraphe, en bas à droite. \\
\hline
\end{tabularx}
\end{center}

\textbf{2. Le schéma « Vous – Moi – Nous ».} Ce schéma organise le fond de la lettre : le paragraphe « Vous » montre que le candidat connaît l'entreprise et le poste ; le paragraphe « Moi » présente la formation, les compétences et les qualités \emph{appuyées par des preuves} ; le paragraphe « Nous » projette la collaboration et sollicite un entretien.

\textbf{3. Les trois figures de style au service de la persuasion.}

\begin{center}
\begin{tabularx}{\linewidth}{|l|X|X|}
\hline
\rowcolor{popblueL} \textbf{Figure} & \textbf{Définition} & \textbf{Exemple dans une lettre} \\
\hline
Gradation & Énumération de termes d'intensité croissante. & « Rigoureux, méthodique et parfaitement autonome, je m'adapterai vite. » \\
\hline
Antithèse & Opposition de deux termes ou idées. & « Peu expérimenté, mais extrêmement motivé. » \\
\hline
Métonymie & On désigne une chose par le nom d'une chose liée. & « J'ai fait mes classes au lycée technique de Bafoussam » (l'établissement pour la formation). \\
\hline
\end{tabularx}
\end{center}

\courssub{Évaluation formative : produire une lettre à partir d'une annonce réelle}

L'évaluation formative vérifie, en situation individuelle, que chaque élève sait mobiliser seul l'ensemble de la compétence. Elle porte sur une annonce authentique publiée au Cameroun.

\begin{exercice}[Évaluation formative — durée : 1 heure]
Voici une annonce parue dans un quotidien camerounais : « \emph{La SODECOTON de Garoua recrute des techniciens de maintenance industrielle titulaires d'un Baccalauréat technique ou d'un BTS. Envoyer lettre de motivation et CV au Directeur des Ressources Humaines, BP 302, Garoua.} »

Rédigez la lettre de motivation complète d'un candidat de votre profil (coordonnées fictives). Votre lettre doit respecter la structure externe (8 éléments), le schéma « Vous – Moi – Nous », employer au moins deux connecteurs logiques et une figure de style, et comporter une formule de politesse soutenue.
\end{exercice}

Le barème de correction, communiqué aux élèves avant l'épreuve, répartit les points ainsi : structure externe (4 points), structure interne et pertinence du fond (6 points), qualités prouvées et argumentation (4 points), style et figures (3 points), orthographe et accords (3 points), soit 20 points.

\begin{exemplebox}[Barème commenté d'une copie type Baccalauréat technique]
Copie de l'élève A., notée 13/20.
\begin{itemize}
\item \textbf{Points gagnés} : les 8 éléments externes sont présents et bien placés (4/4) ; le schéma « Vous – Moi – Nous » est respecté et le paragraphe « Vous » cite précisément l'entreprise (5/6) ; la gradation « ponctuel, rigoureux et infatigable » est pertinente (2/3).
\item \textbf{Points perdus} : l'objet est trop vague (« Objet : emploi » au lieu de « candidature au poste de technicien de maintenance ») : $-1$ point ; la qualité « je suis dynamique » n'est appuyée par aucune preuve : $-2$ points ; deux fautes d'accord (« les compétences que j'ai \emph{acquis} » au lieu d'« acquises ») : $-2$ points ; formule de politesse écourtée (« Merci d'avance ») : $-1$ point.
\end{itemize}
\textbf{Leçon à tirer} : la copie perd 7 points non pas sur les grandes connaissances, mais sur des détails de précision et de relecture. Une relecture méthodique de dix minutes aurait permis de gagner au moins 4 points.
\end{exemplebox}

\courssub{Analyse des erreurs récurrentes de la classe}

Après correction, le professeur dépouille les copies et classe les erreurs par catégories. Ce travail statistique, fait avec la classe, permet de visualiser les difficultés collectives et d'orienter la remédiation. Dans notre classe de Première technique (fictive, 45 élèves), le dépouillement donne les résultats suivants.

\begin{popfigure}
\begin{tikzpicture}
\begin{axis}[
  ybar, bar width=0.7cm,
  width=0.95\linewidth, height=7cm,
  ymin=0, ymax=30,
  ylabel={Nombre d'élèves concernés},
  symbolic x coords={ObjetVague, PolitesseFamiliere, QualitesSansPreuves, FautesAccord, ConnecteursAbsents},
  xtick=data,
  xticklabels={Objet vague, Politesse familière, Qualités sans preuves, Fautes d'accord, Connecteurs absents},
  x tick label style={rotate=20, anchor=east, font=\small},
  nodes near coords, nodes near coords style={font=\small},
  axis line style={popline}, tick style={popline}
]
\addplot[fill=poporange, draw=popdark] coordinates {(ObjetVague,18) (PolitesseFamiliere,24) (QualitesSansPreuves,27) (FautesAccord,21) (ConnecteursAbsents,12)};
\end{axis}
\end{tikzpicture}
\legende{Histogramme des erreurs relevées dans les copies de la classe (effectif fictif de 45 élèves). Trois difficultés dominent : les qualités affirmées sans preuves, les formules de politesse familières et les fautes d'accord. Ce diagnostic justifie les ateliers de remédiation.}
\end{popfigure}

Quatre erreurs reviennent le plus souvent :

\begin{enumerate}
\item \textbf{L'objet vague} : « Objet : demande d'emploi » ne renseigne pas le destinataire. L'objet doit nommer le poste précis : « Objet : candidature au poste de technicien de maintenance industrielle ».
\item \textbf{La politesse familière} : « Salut », « Merci d'avance », « Amicalement » sont des formules de lettre amicale, interdites dans une lettre officielle.
\item \textbf{Les qualités sans preuves} : écrire « je suis sérieux » ne convainc personne ; il faut une preuve : « sérieux, comme en témoigne mon assiduité récompensée par le prix d'excellence de ma classe ».
\item \textbf{Les fautes d'accord} : accord du participe passé avec « avoir » quand le COD est placé avant (« les compétences que j'ai acquises »), accord sujet-verbe (« mes qualités \emph{sont} »).
\end{enumerate}

\courssub{Activités de remédiation différenciées}

\begin{attention}[La remédiation n'est pas une punition]
La \cle{remédiation} est un \textbf{second passage ciblé} sur le point non acquis. Elle ne concerne que la notion fragile, sous une forme nouvelle et plus courte. Être orienté vers un atelier de remédiation n'est pas une sanction : c'est une chance de transformer une erreur en acquis définitif. L'élève qui a réussi peut servir de tuteur dans un atelier, ce qui consolide aussi sa propre maîtrise.
\end{attention}

Chaque élève est orienté, selon ses erreurs dominantes, vers l'un des ateliers suivants, animés en petits groupes.

\textbf{Atelier 1 — Les formules de politesse.} Les élèves trient des formules en deux colonnes (familières / soutenues), puis réécrivent les conclusions de trois lettres fautives. Ils mémorisent un modèle complet : « Dans l'attente d'une réponse que j'espère favorable, veuillez agréer, Monsieur le Directeur, l'expression de ma considération distinguée. »

\textbf{Atelier 2 — Les connecteurs logiques.} À partir d'un paragraphe « Moi » décousu, les élèves insèrent les connecteurs adaptés (d'abord, ensuite, en effet, de plus, c'est pourquoi, enfin) et justifient chaque choix. L'objectif est de faire sentir que le connecteur rend visible la logique du raisonnement.

\textbf{Atelier 3 — L'orthographe et les accords.} Dictée ciblée de phrases extraites des copies (anonymisées), suivie d'une correction collective argumentée : chaque accord est expliqué par sa règle, puis réinvesti dans deux phrases produites par l'élève.

\textbf{Réécriture ciblée.} Chaque élève reprend sa propre copie et réécrit \emph{uniquement} les passages signalés : objet, conclusion, ou paragraphe des qualités. La nouvelle version est comparée à l'ancienne pour mesurer le progrès.

\begin{methode}[La fiche de remédiation personnelle]
Chaque élève remplit sa fiche en quatre étapes :
\begin{enumerate}
\item \textbf{Identifier} : je relis ma copie corrigée et je note mes \textbf{deux erreurs principales} (par exemple : formule de politesse familière ; qualités sans preuves).
\item \textbf{Comprendre} : pour chaque erreur, j'écris la règle exacte que j'ai enfreinte (par exemple : une lettre officielle exige une formule de politesse complète et soutenue).
\item \textbf{Corriger} : je réécris le passage fautif de ma lettre en appliquant la règle.
\item \textbf{Vérifier} : je rédige une phrase ou un paragraphe \emph{nouveau} qui applique la règle, et je le fais valider par le tuteur ou le professeur.
\end{enumerate}
La fiche est conservée dans le cahier : elle servira de pense-bête avant chaque production d'écrit.
\end{methode}

\begin{exoresolu}[Corriger une conclusion fautive]
Énoncé. Voici la conclusion écrite par un élève : « Voilà, j'espère que vous allez me prendre. Merci d'avance et à bientôt. » Relevez les erreurs et réécrivez cette conclusion dans un style soutenu, adapté à une lettre de motivation adressée au Directeur des Ressources Humaines de la SODECOTON.
\tcblower
Solution détaillée. Trois erreurs : « Voilà » est une transition familière ; « j'espère que vous allez me prendre » est une formulation directe et familière qui ne projette pas la collaboration (pas de « Nous ») ; « Merci d'avance et à bientôt » est une formule de lettre amicale.

Réécriture attendue : « Convaincu que mon dynamisme et ma formation technique sauront répondre aux exigences de votre entreprise, je me tiens à votre disposition pour un entretien. Dans l'attente d'une réponse que j'espère favorable, veuillez agréer, Monsieur le Directeur, l'expression de ma considération distinguée. » Cette version projette la collaboration (« Nous »), sollicite un entretien et se termine par une formule de politesse complète.
\end{exoresolu}

\courssub{Réinvestissement : appliquer la compétence à d'autres situations}

La compétence acquise est transférable : toute lettre officielle repose sur la même structure externe, le même respect du destinataire et la même clarté argumentative. Seuls l'objet et le contenu du corps changent.

\begin{exercice}[Réinvestissement — au choix]
\begin{enumerate}
\item \textbf{Demande de stage} : rédigez une lettre au directeur d'une entreprise de votre ville (garage, hôpital, banque) pour solliciter un stage de découverte de deux semaines.
\item \textbf{Candidature à un concours} : rédigez une lettre au recteur d'université ou au directeur d'une école de formation pour demander l'autorisation de composer à un concours d'entrée.
\item \textbf{Lettre au maire} : rédigez une lettre au maire de votre commune pour demander l'éclairage public d'un quartier ou la réparation d'une piste, en argumentant avec des faits précis.
\end{enumerate}
Pour chaque lettre, vérifiez : les 8 éléments externes, un objet précis, un plan clair, une formule de politesse soutenue.
\end{exercice}

\begin{corrige}[Exercice de réinvestissement — indications]
\begin{enumerate}
\item \textbf{Demande de stage} : objet « demande de stage de découverte » ; le « Vous » montre la connaissance de l'entreprise ; le « Moi » présente la filière technique suivie ; le « Nous » insiste sur l'intérêt mutuel (main-d'œuvre motivée, découverte du métier).
\item \textbf{Concours} : ton administratif ; joindre la mention des pièces exigées ; le corps expose le diplôme obtenu et le projet de formation.
\item \textbf{Lettre au maire} : registre de la demande citoyenne ; faits précis (rue, nombre de foyers concernés) ; arguments d'intérêt général (sécurité, vie du quartier) ; formule de politesse adaptée à une autorité (« Veuillez agréer, Monsieur le Maire, l'expression de ma haute considération »).
\end{enumerate}
\end{corrige}

\begin{savaistu}[Une compétence pour toute la vie]
Savoir rédiger une lettre officielle sert toute la vie : obtenir un \textbf{emploi}, décrocher un \textbf{stage}, s'inscrire à un \textbf{concours}, adresser une \textbf{réclamation} à une entreprise, formuler une \textbf{demande de bourse} ou écrire à une administration. Au Cameroun comme ailleurs, beaucoup de candidatures sont écartées non pas par manque de compétences, mais à cause d'une lettre mal présentée ou mal écrite. La lettre est votre première image : elle parle de vous avant même que vous n'entriez dans la pièce. La maîtrise de la \emph{communication par l'image} (affiches sécuritaires) complète cette compétence en développant la lecture critique des supports visuels officiels.
\end{savaistu}

\begin{aretenir}[Ce qu'il faut retenir de cette séquence]
\begin{itemize}
\item Le compte-rendu collectif transforme les acquis dispersés en savoir stable, consigné sur une fiche-bilan affichée.
\item L'évaluation formative vérifie la compétence en situation individuelle, avec un barème connu à l'avance.
\item Le dépouillement des erreurs (objet vague, politesse familière, qualités sans preuves, fautes d'accord) oriente des ateliers de remédiation différenciés.
\item La fiche de remédiation personnelle (identifier, comprendre, corriger, vérifier) rend chaque élève acteur de son progrès.
\item La compétence se réinvestit dans toutes les lettres officielles : stage, concours, bourse, réclamation, lettre au maire.
\end{itemize}
\end{aretenir}

\newpage

\coursec{Activité d'intégration}

\coursec{Produire une lettre officielle : la lettre de motivation}

\courssub{Objectifs de l'activité}

À l'issue de cette activité d'intégration, l'élève sera capable de :
\begin{itemize}
\item analyser une annonce d'offre de stage ou d'emploi pour en dégager les qualités et compétences attendues par le destinataire ;
\item rédiger une lettre de motivation complète respectant la \cle{structure externe} (en-tête, objet, formule d'appel, formule de politesse, signature) et la \cle{structure interne} (introduction en « vous », corps en « moi », conclusion en « nous ») ;
\item employer à bon escient les figures de style étudiées dans l'unité : la \cle{gradation}, l'\cle{antithèse} et la \cle{métonymie}, pour rendre son propos plus convaincant ;
\item analyser une affiche d'indications sécuritaires (communication par l'image) et en dégager le message ;
\item évaluer la production d'un camarade à l'aide d'une grille critériée, puis améliorer sa propre lettre et justifier oralement ses choix de rédaction.
\end{itemize}

\courssub{Situation}

\textbf{Contexte.} À l'approche des grandes vacances, le club d'animation pédagogique du Lycée technique de Douala-Bassa organise un concours intitulé « Je décroche mon stage ». Trois offres réelles, adaptées au niveau Première technique, sont affichées au tableau d'information. Chaque élève de la classe doit choisir \textbf{une seule} offre et produire une lettre de motivation complète. Les trois meilleures lettres seront lues en classe, affichées au tableau d'honneur et effectivement envoyées.

\textbf{Offre n\textsuperscript{o} 1 --- Stage de vacances à la SODECOTON de Garoua.} La Société de Développement du Coton du Cameroun (SODECOTON), sise à Garoua, recrute dix stagiaires de niveau Première technique pour un stage d'observation et de découverte des métiers industriels, du 15 juillet au 31 août. Profil recherché : élève sérieux, ponctuel, intéressé par la mécanique ou l'électrotechnique, sachant travailler en équipe. Dossier à adresser au Chef du personnel, B.P. 302, Garoua, avant le 20 juin.

\textbf{Offre n\textsuperscript{o} 2 --- Poste de caissier vacataire dans une microfinance de Bafoussam.} La COOPEC « Épargne et Avenir » de Bafoussam cherche deux caissiers vacataires pour le mois d'août. Profil recherché : élève de Première ou Terminale technique (spécialité comptabilité ou gestion appréciée), honnête, rigoureux, à l'aise avec les chiffres et avec la clientèle. Lettre de motivation à déposer au siège, quartier Marché A, Bafoussam, ou à envoyer au Directeur de l'agence.

\textbf{Offre n\textsuperscript{o} 3 --- Formation de technicien au Lycée technique de Douala-Bassa.} Dans le cadre d'un partenariat avec une entreprise de BTP, le Lycée technique de Douala-Bassa sélectionne cinq élèves de Première technique pour une formation pratique de quatre semaines en maintenance des engins de chantier. Profil recherché : élève motivé, assidu, ayant de bons résultats en technologie, respectueux des consignes de sécurité. Candidature à remettre au Proviseur avant le 10 juin.

\textbf{Préambule obligatoire.} Avant de rédiger, la classe observe deux affiches d'indications sécuritaires apportées par le professeur : une affiche « Port du casque obligatoire » (chantier de la SODECOTON) et une affiche « Défense de fumer --- zone de stockage » (atelier du lycée). Il s'agit de vérifier que chacun sait lire un message transmis par l'image : un pictogramme, une couleur (rouge pour l'interdiction, bleu pour l'obligation), une consigne courte.

\courssub{Consignes}

\begin{enumerate}
\item \textbf{Lecture de l'image.} Pour chaque affiche, décris ce que tu vois (pictogramme, couleurs, texte) et dégage le message transmis. À quelle offre chaque affiche peut-elle être rattachée ? Pourquoi le respect des consignes de sécurité peut-il devenir un argument dans ta lettre ?
\item \textbf{Analyse de l'annonce.} Choisis une offre. Relève dans un tableau à deux colonnes : (a) ce que le destinataire attend (qualités, compétences) ; (b) ce que tu possèdes qui y répond (études, qualités personnelles, expériences même modestes).
\item \textbf{Préparation.} Rédige le plan de ta lettre : introduction en « vous » (qui montre que tu connais l'entreprise et le poste), corps en « moi » (ta formation, tes qualités, ta motivation), conclusion en « nous » (la collaboration future, la demande d'entretien).
\item \textbf{Rédaction du brouillon.} Rédige individuellement ta lettre complète, en respectant la structure externe : en-tête (coordonnées de l'expéditeur, lieu et date, coordonnées du destinataire), objet, formule d'appel, corps en trois paragraphes, formule de politesse, signature. Ton corps de lettre doit contenir \textbf{au moins une gradation et une antithèse}.
\item \textbf{Évaluation croisée.} Par groupes de quatre, échangez vos brouillons et évaluez la lettre d'un camarade à l'aide de la grille critériée ci-dessous. Note chaque critère sur 2 et propose une amélioration concrète.
\item \textbf{Version finale.} Produis ta version définitive en tenant compte des remarques reçues. Prépare-toi à justifier oralement \textbf{deux choix de rédaction} (par exemple : pourquoi cette gradation, pourquoi cette formule de politesse, pourquoi cet ordre d'arguments).
\end{enumerate}

\begin{center}
\begin{tabularx}{\linewidth}{|l|c|X|}
\hline
\rowcolor{popblueL} \textbf{Critère} & \textbf{Note /2} & \textbf{Remarque du correcteur} \\
\hline
Structure externe complète et bien placée & & \\
\hline
Introduction en « vous » adaptée au destinataire & & \\
\hline
Corps en « moi » : qualités en lien avec l'annonce & & \\
\hline
Présence d'une gradation et d'une antithèse correctes & & \\
\hline
Conclusion en « nous » + formule de politesse & & \\
\hline
Orthographe, grammaire, présentation soignée & & \\
\hline
\end{tabularx}
\end{center}

\courssub{Ressources à mobiliser}

\begin{aretenir}[Ce qu'il faut se rappeler]
\begin{itemize}
\item \textbf{Structure externe} : coordonnées de l'expéditeur (en haut à gauche), lieu et date (à droite), coordonnées du destinataire (à droite, sous la date), objet de la lettre, formule d'appel (\emph{Monsieur le Directeur, Madame, Monsieur}), corps de la lettre, formule de politesse, signature.
\item \textbf{Structure interne} : le schéma \cle{Vous -- Moi -- Nous}. L'introduction parle du destinataire et de l'offre ; le corps présente le candidat et ses atouts ; la conclusion évoque la future collaboration et sollicite une réponse.
\item \textbf{Gradation} : énumération de termes de force croissante (ou décroissante) : \emph{« sérieux, appliqué, passionné »}.
\item \textbf{Antithèse} : opposition de deux termes ou idées pour faire ressortir une pensée : \emph{« jeune encore, mais déjà déterminé »}.
\item \textbf{Métonymie} : désigner une chose par le nom d'une chose liée : \emph{« mettre la main à la pâte »} pour dire travailler concrètement.
\item \textbf{Communication par l'image} : une affiche sécuritaire associe un pictogramme, une couleur codée (rouge = interdiction, bleu = obligation, jaune = danger) et une consigne brève.
\end{itemize}
\end{aretenir}

\courssub{Méthode : Structurer sa lettre selon le schéma « Vous -- Moi -- Nous »}

\begin{methode}[Les trois temps de la lettre de motivation]
\begin{enumerate}
\item \textbf{Introduction (VOUS) :} Accrocher le lecteur en montrant que l'on connaît l'entreprise, l'offre et ses enjeux. Citer l'origine de l'information (annonce, site web, recommandation). Exprimer l'intérêt pour la structure.
\item \textbf{Corps (MOI) :} Présenter son parcours (formation, spécialité), ses qualités (en lien direct avec l'annonce) et sa motivation. C'est ici que l'on insère la \cle{gradation} (pour monter en puissance) et l'\cle{antithèse} (pour transformer un point faible en force). Utiliser une \cle{métonymie} pour montrer son engagement concret.
\item \textbf{Conclusion (NOUS) :} Projeter la collaboration future (« Ce stage nous permettrait de... »), demander un entretien, remercier. Terminer par une formule de politesse administrative complète et adaptée au destinataire.
\end{enumerate}
\end{methode}

\courssub{Étude de documents iconiques : les affiches de sécurité}

\begin{experience}[Analyse d'affiches de consignes de sécurité --- Communication par l'image]
\textbf{Matériel :} Deux affiches projetées ou affichées au tableau.
\begin{itemize}
\item \textbf{Affiche A :} Fond bleu, pictogramme blanc (tête casquée), texte « PORT DU CASQUE OBLIGATOIRE ».
\item \textbf{Affiche B :} Fond rouge, pictogramme blanc (cigarette barrée), texte « DÉFENSE DE FUMER --- ZONE DE STOCKAGE ».
\end{itemize}

\textbf{Protocole d'analyse (à reproduire dans le cahier) :}
\begin{enumerate}
\item \textbf{Description objective :} Forme, couleurs dominantes, pictogramme, texte, police.
\item \textbf{Décodage du sens :} Type de consigne (obligation / interdiction / danger), zone concernée, public visé.
\item \textbf{Lien avec l'offre :} Affiche A $\rightarrow$ Offre n\textsuperscript{o} 1 (SODECOTON, chantier/industrie). Affiche B $\rightarrow$ Offre n\textsuperscript{o} 3 (BTP, atelier) ou Offre n\textsuperscript{o} 1 (stockage coton).
\item \textbf{Argumentaire :} « Le respect strict des consignes de sécurité (port du casque, interdiction de fumer) fait partie de mes habitudes d'atelier au lycée. »
\end{enumerate}

\begin{popfigure}
\begin{tikzpicture}[scale=0.8, every node/.style={transform shape}]
  % Affiche A : Obligation (Bleu)
  \begin{scope}[xshift=-5cm]
    \fill[popblue] (-2.5,-3) rectangle (2.5,3);
    \node[white, font=\Large\bfseries] at (0,1.5) {\textsf{PORT DU CASQUE}};
    \node[white, font=\Large\bfseries] at (0,0.5) {\textsf{OBLIGATOIRE}};
    % Pictogramme casque simplifié
    \draw[white, line width=3pt] (0,-0.5) ellipse (1.2 and 0.6);
    \draw[white, line width=3pt] (-1.2,-0.5) -- (-1.2,-1.5) -- (1.2,-1.5) -- (1.2,-0.5);
    \draw[white, line width=2pt] (-0.5,-0.5) .. controls (-0.5,-1.2) and (0.5,-1.2) .. (0.5,-0.5);
    \node[white, font=\small] at (0,-2.5) {ISO 7010 - M002};
  \end{scope}
  % Affiche B : Interdiction (Rouge)
  \begin{scope}[xshift=5cm]
    \fill[poporange!80!popdark] (-2.5,-3) rectangle (2.5,3); % Rouge vif
    \node[white, font=\Large\bfseries] at (0,1.5) {\textsf{DÉFENSE DE FUMER}};
    \node[white, font=\Large\bfseries] at (0,0.5) {\textsf{ZONE DE STOCKAGE}};
    % Pictogramme cigarette barrée
    \draw[white, line width=3pt] (-1,1) -- (1,-1); % Barre
    \draw[white, line width=2pt] (-1.2,-0.5) -- (-0.2,-0.5); % Filtre
    \draw[white, line width=2pt] (-1.2,-0.5) -- (-1.2,0.5); % Côté filtre
    \draw[white, line width=2pt] (-0.2,-0.5) -- (-0.2,0.5); % Côté allumette
    \draw[white, line width=2pt] (-1.2,0.5) -- (-0.2,0.5); % Haut
    \draw[white, line width=2pt, decorate, decoration={snake, amplitude=1pt, segment length=4pt}] (-0.7,0.5) -- (-0.7,1.2); % Fumée
    \node[white, font=\small] at (0,-2.5) {ISO 7010 - P002};
  \end{scope}
\end{tikzpicture}
\legende{Modèles simplifiés d'affiches de sécurité normalisées (ISO 7010). À gauche : obligation (fond bleu) ; à droite : interdiction (fond rouge).}
\end{popfigure}
\end{experience}

\courssub{Démarche et corrigé commenté}

\begin{exoresolu}[Modèle de lettre réussie (offre n\textsuperscript{o} 1 : SODECOTON)]
Énoncé : un élève de Première technique F (électrotechnique) répond à l'offre de stage de la SODECOTON. Produire le brouillon puis la version finale de sa lettre.

\tcblower

\textbf{Étape 1 --- Analyse de l'annonce.} Le destinataire attend : sérieux, ponctualité, intérêt pour la mécanique ou l'électrotechnique, esprit d'équipe. Le candidat possède : une formation en électrotechnique, de bons résultats en technologie, l'expérience des travaux pratiques en atelier, le respect des consignes de sécurité (vérifié dans le préambule grâce aux affiches).

\textbf{Étape 2 --- Plan.} Introduction « vous » : la SODECOTON, moteur de l'économie du Nord, propose un stage. Corps « moi » : formation, qualités (avec gradation et antithèse), motivation. Conclusion « nous » : collaboration utile, demande d'entretien, formule de politesse.

\textbf{Étape 3 --- Version finale commentée.}

\begin{exemplebox}[Lettre de motivation --- Candidature stage SODECOTON Garoua]
\emph{Etoudi Marcel, élève en Première technique F4\\
Lycée technique de Garoua, B.P. 45 Garoua\\
Tél. : 6 90 00 00 00}

\emph{Garoua, le 5 juin 2025}

\emph{À Monsieur le Chef du personnel de la SODECOTON\\
B.P. 302, Garoua}

\textbf{Objet :} Candidature au stage de vacances du 15 juillet au 31 août 2025.

Monsieur le Chef du personnel,

Votre entreprise, qui fait vivre des milliers de familles du Grand Nord et place le coton camerounais sur les marchés du monde entier, propose aux élèves de Première technique un stage de découverte des métiers industriels. Cette opportunité correspond exactement au projet professionnel que je construis depuis mon entrée au lycée technique.

Élève en Première technique F4, spécialité électrotechnique, je me forme chaque semaine aux gestes du métier dans les ateliers de mon établissement, où le respect des consignes de sécurité --- port du casque, port des gants, ordre et propreté --- est devenu pour moi une seconde nature. Mes professeurs me reconnaissent un élève \textbf{attentif, sérieux, passionné} ; et si je suis \textbf{jeune encore, je suis déjà déterminé} à mettre la main à la pâte. Travailler en équipe ne m'effraie pas : au contraire, c'est auprès des autres que j'apprends le plus vite.

Ce stage serait pour moi l'occasion de découvrir le monde de l'entreprise, et pour vous celui de compter sur un stagiaire disponible et volontaire. Je reste à votre disposition pour tout entretien que vous jugerez utile.

Dans l'attente d'une réponse favorable, je vous prie d'agréer, Monsieur le Chef du personnel, l'expression de ma respectueuse considération.

\emph{Signature}
\end{exemplebox}

\textbf{Étape 4 --- Justification orale de deux choix.}
\begin{itemize}
\item \emph{Choix 1 (Gradation) :} La gradation ascendante « attentif, sérieux, passionné » place le terme le plus fort à la fin ; elle donne un élan à la phrase et marque une motivation qui grandit.
\item \emph{Choix 2 (Antithèse) :} L'antithèse « jeune encore, mais déjà déterminé » transforme un point faible (le jeune âge, le manque d'expérience) en argument : la détermination compense l'inexpérience.
\item \emph{Choix 3 (Métonymie) :} « Mettre la main à la pâte » évoque le travail concret, manuel, sans le nommer ; cela prouve que le candidat n'a pas peur de l'effort physique en atelier.
\end{itemize}

\textbf{Étape 5 --- Vérification avec la grille.}
\begin{center}
\begin{tabularx}{\linewidth}{|l|c|X|}
\hline
\rowcolor{popblueL} \textbf{Critère} & \textbf{Note /2} & \textbf{Remarque du correcteur} \\
\hline
Structure externe complète et bien placée & 2 & En-tête, date, destinataire, objet, appel, politesse, signature : tout y est. \\
\hline
Introduction en « vous » adaptée au destinataire & 2 & Citation de l'entreprise, de son rôle, lien avec le projet pro. \\
\hline
Corps en « moi » : qualités en lien avec l'annonce & 2 & Électrotechnique, sécurité, équipe, figures de style. \\
\hline
Présence d'une gradation et d'une antithèse correctes & 2 & Gradation (3 termes), Antithèse (jeune/déterminé). \\
\hline
Conclusion en « nous » + formule de politesse & 2 & Projection collaboration + demande entretien + formule complète. \\
\hline
Orthographe, grammaire, présentation soignée & 2 & Aucune faute, paragraphes aérés, mise en page standard. \\
\hline
\end{tabularx}
\end{center}
\textbf{Total : 12/12.} Lettre prête pour le tableau d'honneur et l'envoi effectif.
\end{exoresolu}

\begin{attention}[Erreurs fréquentes à éviter]
\begin{itemize}
\item Commencer la lettre par « moi » (« Je m'appelle..., je veux... ») : l'introduction doit d'abord parler du destinataire et de l'offre.
\item Oublier l'objet de la lettre ou le placer après la formule d'appel.
\item Employer une formule de politesse familière (« amicalement », « merci d'avance », « cordialement » seul) dans une lettre officielle. \textbf{Exiger la formule administrative complète.}
\item Forcer les figures de style : une gradation maladroite (« bon, très bon, excellent, génial ») affaiblit la lettre au lieu de la renforcer. Les termes doivent appartenir au registre formel.
\item Recopier l'annonce sans montrer ce que l'on peut apporter à l'entreprise (pas de copier-coller).
\item Négliger la présentation : ratures, marges non respectées, en-tête incomplet.
\end{itemize}
\end{attention}

\courssub{Bilan de l'activité}

Cette activité d'intégration a permis de mettre en évidence trois acquis essentiels. D'abord, une lettre de motivation efficace naît d'une \textbf{lecture précise de l'annonce} : on ne convainc un destinataire qu'en répondant point par point à ce qu'il attend. Ensuite, la réussite repose sur le respect conjoint de la \textbf{structure externe} (la forme visible, qui donne la première impression) et de la \textbf{structure interne} en « vous -- moi -- nous », qui organise l'argumentation. Enfin, les \textbf{figures de style} --- gradation, antithèse, métonymie --- ne sont pas des ornements inutiles : bien employées, elles donnent de la force et du rythme à la candidature. L'évaluation croisée et la justification orale des choix de rédaction ont en outre montré que rédiger, c'est réécrire : la version finale, améliorée grâce au regard des pairs, est toujours supérieure au premier brouillon. Ces compétences, directement mobilisables au Probatoire et au Baccalauréat technique, serviront aussi dans la vie réelle : stage, emploi, formation.

\newpage

\coursec{Méthodes et exercices résolus}

\coursec{Produire une lettre officielle : la lettre de motivation}

\courssub{Découverte : analyse de lettres de motivation authentiques (Textes iconiques)}

\begin{methode}[Analyser des modèles pour en extraire les codes professionnels]
\textbf{Objectif :} Identifier, par l'observation de lettres réelles (textes iconiques 1 à 5), les constantes structurelles, le ton professionnel et les stratégies d'argumentation efficaces.
\begin{enumerate}
\item \textbf{Collecte :} Constituer un corpus de 3 à 5 lettres de motivation réussies (stages, emplois techniques, BTS/DUT) issues d'anciens élèves ou de banques de modèles validées par l'enseignant.
\item \textbf{Grille d'analyse (en groupes) :} Pour chaque lettre, relever :
\begin{itemize}
\item La structure externe (alignements, mentions obligatoires).
\item L'accroche (Introduction) : présence du « Vous » (entreprise), du « Moi » (poste), du « Lien » (source).
\item Le corps : nombre de paragraphes, idée directrice par paragraphe, verbes d'action, données chiffrées, vocabulaire technique.
\item La conclusion : appel à l'action, disponibilité, formule de politesse.
\item Le style : figures de style (gradation, métonymie), connecteurs, absence de familier.
\end{itemize}
\item \textbf{Mise en commun :} Établir au tableau la « check-list des indispensables » qui servira de référence pour toute la séquence.
\end{enumerate}
\cle{Compétence visée :} Passer du statut d'élève (qui écrit pour une note) à celui de professionnel (qui écrit pour obtenir un entretien).
\end{methode}

\begin{exemplebox}[Exemple d'analyse guidée (Texte iconique 1 -- Extrait)]
\textbf{Extrait :} \textit{« Votre projet d'extension de l'unité de traitement des eaux usées de Douala-Bonabéri (phase 2) répond à un enjeu majeur de santé publique. Fort de mon stage de 6 mois à la CAMWATER où j'ai piloté la réhabilitation de 3 postes de refoulement, je vous propose ma candidature au poste de Technicien Hydraulique (Ref: ONEP-2024/012). »}
\textbf{Analyse :}
\begin{itemize}
\item \textbf{Vous :} Projet concret, localisation précise, enjeu identifié (santé publique).
\item \textbf{Moi :} Poste exact, référence de l'annonce.
\item \textbf{Preuve :} Stage CAMWATER, action « piloté », résultat « 3 postes », durée « 6 mois ».
\item \textbf{Ton :} Professionnel, direct, pas de « Par la présente... ».
\end{itemize}
\end{exemplebox}

\courssub{Analyser et respecter la structure externe de la lettre de motivation}

\begin{methode}[Identifier et ordonner les éléments de la structure externe]
\textbf{Objectif :} Maîtriser la mise en page normée (normes administratives camerounaises) pour une présentation professionnelle immédiate.
\begin{enumerate}
\item \textbf{En-tête expéditeur (haut gauche) :} Nom, Prénom, Adresse complète (quartier, ville), Téléphone, Email. \emph{Astuce :} Aligner à gauche, police lisible (Times New Roman 12 ou Arial 11).
\item \textbf{En-tête destinataire (haut droite, sous l'expéditeur) :} Nom du responsable (si connu) ou Titre du poste (ex: Directeur des Ressources Humaines), Nom de l'entreprise/structure, Adresse de l'entreprise.
\item \textbf{Lieu et Date (à droite, sous le destinataire) :} Ex: \textit{Yaoundé, le 15 octobre 2023}. Mois en toutes lettres, pas de chiffre romain.
\item \textbf{Objet (aligné à gauche, gras, souligné) :} \textbf{Objet : Candidature au poste de [Intitulé exact du poste] / Demande de stage [Spécialité]}. Mentionner la référence de l'annonce si existante (Ref: N°...).
\item \textbf{Formule d'appel (salutation) :} \textit{Monsieur le Directeur,} ou \textit{Madame la Directrice,} (suivi d'une virgule, pas de point). Éviter « Monsieur, » seul si le titre est connu.
\item \textbf{Corps de la lettre :} Introduction, Développement (2-3 paragraphes), Conclusion. Paragraphes séparés par une ligne blanche.
\item \textbf{Formule de politesse de clôture (bas gauche) :} Longue, adaptée au destinataire (voir fiche mémo). Ex: \textit{Je vous prie d'agréer, Monsieur le Directeur, l'expression de ma considération distinguée.}
\item \textbf{Signature (bas droite) :} Manuscritte (sur papier) ou image insérée (numérique), précédée du Nom Prénom tapuscrit.
\item \textbf{Mention des pièces jointes (bas gauche, sous la formule de politesse) :} \textit{Pièces jointes : CV, Copies des diplômes, Attestations de travail, Extrait d'acte de naissance.}
\end{enumerate}
\cle{Règle d'or :} La lettre tient sur \textbf{une seule page A4}. Mares : 2,5 cm (haut/bas), 3 cm (gauche/droite). Interligne 1,15 ou 1,5.
\end{methode}

\begin{exoresolu}[Reconstituer la structure externe d'une lettre officielle]
\textbf{Énoncé :} Vous trouvez ci-dessous les éléments d'une lettre de motivation mélangés. Replacez-les dans l'ordre exact de la structure externe normée en indiquant l'alignement (Gauche / Droite) pour chaque bloc.

\begin{itemize}
\item[A.] \textit{Je vous prie d'agréer, Monsieur le Directeur, l'expression de ma considération distinguée.}
\item[B.] \textit{Objet : Candidature au poste de Technicien de Maintenance Industrielle (Ref: ANPE-2023/045)}
\item[C.] \textit{Monsieur le Directeur,}
\item[D.] \textit{Ngonso, le 20 novembre 2023}
\item[E.] \textit{M. Mbarga Jean, Directeur des Ressources Humaines, Société CAMINDUS, B.P. 123 Douala}
\item[F.] \textit{Nom : FOPA Eric \quad B.P. 456 Yaoundé \quad Tel: 6XX XX XX XX \quad Email: eric.fopa@email.cm}
\item[G.] \textit{Pièces jointes : CV, Copie du BTS Maintenance, Attestation de stage}
\item[H.] \textit{FOPA Eric} (signature)
\item[I.] \textit{[Corps de la lettre : Introduction, Développement, Conclusion]}
\end{itemize}

\tcblower
\textbf{Solution détaillée :}

Voici l'ordre correct (1 à 9) avec alignement et justifications :

\begin{center}
\begin{tabularx}{\linewidth}{|c|X|c|X|}\hline
\textbf{Ordre} & \textbf{Contenu} & \textbf{Alignement} & \textbf{Justification / Règle} \\ \hline
1 & F. \textit{FOPA Eric, B.P. 456 Yaoundé...} & Gauche & En-tête expéditeur (identité complète + contacts). \\ \hline
2 & E. \textit{M. Mbarga Jean, Directeur RH...} & Droite & En-tête destinataire (sous l'expéditeur, à droite). \\ \hline
3 & D. \textit{Ngonso, le 20 novembre 2023} & Droite & Lieu (commune/ville) + Date du jour (mois en lettres). \\ \hline
4 & B. \textit{Objet : Candidature au poste...} & Gauche & \textbf{Gras + Souligné}. Précis, cite la référence. \\ \hline
5 & C. \textit{Monsieur le Directeur,} & Gauche & Formule d'appel (virgule finale, pas de point). \\ \hline
6 & I. \textit{[Corps de la lettre]} & Justifié & Introduction, Corps (2-3 §), Conclusion. Ligne blanche entre §. \\ \hline
7 & A. \textit{Je vous prie d'agréer...} & Gauche & Formule de politesse \textbf{longue}, adaptée à l'appel (Monsieur le Directeur). \\ \hline
8 & H. \textit{FOPA Eric} & Droite & Signature (manuscrite idéalement) + Nom Prénom tapuscrit au-dessus. \\ \hline
9 & G. \textit{Pièces jointes : CV...} & Gauche & Liste énumérée, en dessous de la formule de politesse (ou à gauche de la signature). \\ \hline
\end{tabularx}
\end{center}

\textbf{Conclusion :} Le respect scrupuleux de cet ordre et de ces alignements conditionne la première impression du recruteur (lisibilité, rigueur professionnelle). Toute inversion (ex: Date avant Destinataire) ou oubli (Objet, Pièces jointes) est pénalisant.
\end{exoresolu}

\courssub{Rédiger une introduction accrocheuse et une conclusion engageante}

\begin{methode}[Structurer l'introduction (le « Vous ») et la conclusion (l'appel à l'action)]
\textbf{Objectif :} Capter l'attention dès la première ligne (Introduction) et provoquer la réaction souhaitée : l'entretien (Conclusion).
\medskip
\textbf{A. L'Introduction (3 à 4 lignes max) : La technique du « Vous $\rightarrow$ Moi $\rightarrow$ Lien »}
\begin{enumerate}
\item \textbf{Le « Vous » (Accroche) :} Montrez que vous connaissez l'entreprise. Citez un fait marquant, une valeur, un projet récent. \textit{Ex: « Votre récente certification ISO 9001 témoigne de l'excellence que vise CAMINDUS... »}
\item \textbf{Le « Moi » (Déclaration d'intention) :} Annoncez clairement la candidature. \textit{Ex: « C'est donc avec un vif intérêt que je vous adresse ma candidature au poste de Technicien de Maintenance... »}
\item \textbf{Le « Lien » (Source) :} Précisez l'origine de l'info. \textit{Ex: « ...suite à votre annonce parue sur le site de l'ANPE le 10 novembre 2023 (Ref: ANPE-2023/045). »}
\end{enumerate}
\cle{Interdit :} « Par la présente, je vous adresse ma candidature... » (Trop passif, banal). « Je suis titulaire d'un BTS... » (Place au corps de lettre).

\medskip
\textbf{B. La Conclusion (2 à 3 lignes) : La technique de l'« Appel à l'action »}
\begin{enumerate}
\item \textbf{Disponibilité :} \textit{« Disponible immédiatement pour un entretien à votre convenance... »}
\item \textbf{Rappel de la valeur ajoutée :} \textit{« ...afin de vous exposer comment mes compétences en maintenance préventive pourront contribuer à la performance de vos lignes de production. »}
\item \textbf{Formule de politesse longue :} Obligatoire, adaptée au titre de l'appel (voir méthode 1).
\end{enumerate}
\cle{Astuce Bac :} Évitez « Dans l'attente de votre réponse, je vous salue. » (Trop faible). Préférez l'engagement actif.
\end{methode}

\begin{exoresolu}[Rédiger l'introduction et la conclusion pour une offre précise]
\textbf{Énoncé :} Offre d'emploi : \textit{« Entreprise AGROCAM (Douala) recherche un Technicien Agroalimentaire (BTS/DUT). Missions : Contrôle qualité HACCP, suivi production, hygiène. Expérience stage exigée. »}
Rédigez l'\textbf{Introduction} et la \textbf{Conclusion} d'une lettre de motivation pour le candidat \textbf{KAMGA Paul}, BTS Agroalimentaire (spécialité Contrôle Qualité), ayant effectué un stage de 3 mois chez CHOCOCAM (ligne boissons).

\tcblower
\textbf{Solution détaillée :}

\textbf{Introduction proposée :}
\begin{quote}
\textit{Votre engagement reconnu dans la promotion des produits locaux « Made in Cameroon » et la rigueur de votre démarche qualité HACCP, certifiée depuis 2021, font d'AGROCAM une référence incontournable du secteur agroalimentaire national. C'est donc avec une motivation déterminée que je vous présente ma candidature au poste de Technicien Agroalimentaire, suite à votre offre publiée sur le portail Emploi.cm en date du 5 janvier 2024.}
\end{quote}
\textbf{Analyse :}
\begin{itemize}
\item \textbf{Vous :} « Votre engagement... démarche qualité HACCP certifiée » (Connaissance entreprise + vocabulaire technique).
\item \textbf{Moi :} « Je vous présente ma candidature au poste de Technicien Agroalimentaire » (Poste exact).
\item \textbf{Lien :} « Suite à votre offre publiée sur Emploi.cm... » (Source précise).
\end{itemize}

\medskip
\textbf{Conclusion proposée :}
\begin{quote}
\textit{Fort de mon stage chez CHOCOCAM où j'ai participé à la validation des plans HACCP sur la ligne boissons, je suis pleinement opérationnel pour intégrer vos équipes de contrôle qualité. Disponible dès à présent pour un entretien, je souhaite vous démontrer en quoi ma rigueur analytique et ma maîtrise des normes ISO 22000 seront un atout pour la satisfaction de vos clients. \newline
Je vous prie d'agréer, Monsieur le Directeur, l'expression de ma considération distinguée.}
\end{quote}
\textbf{Analyse :}
\begin{itemize}
\item \textbf{Preuve concrète :} Stage CHOCOCAM + HACCP (répond à l'exigence « expérience stage »).
\item \textbf{Disponibilité + Objectif :} « Disponible dès à présent pour un entretien... ».
\item \textbf{Projection (Nous) :} « Maîtrise normes ISO 22000... atout pour la satisfaction de vos clients ».
\item \textbf{Politesse :} Formule longue correcte (Monsieur le Directeur).
\end{itemize}
\end{exoresolu}

\courssub{Construire le corps de la lettre : la stratégie « Vous / Moi / Nous »}

\begin{methode}[Argumenter par paragraphes thématiques avec la structure en 3 temps]
\textbf{Objectif :} Démontrer l'adéquation Profil/Poste sans répéter le CV, en prouvant la valeur ajoutée.
\medskip
\textbf{Structure type en 2 ou 3 paragraphes distincts (1 idée = 1 paragraphe) :}

\begin{enumerate}
\item \textbf{Paragraphe 1 : Le « Vous » (L'entreprise et ses besoins) -- Optionnel si bien fait en intro, sinon indispensable.}
Montrez que vous comprenez les enjeux techniques de l'entreprise. Vocabulaire métier requis.
\item \textbf{Paragraphe 2 : Le « Moi » (Mes compétences techniques \textbf{prouvées}) -- \textbf{Le cœur de la lettre.}}
\begin{itemize}
\item Ne listez pas les diplômes (CV s'en charge).
\item Citez \textbf{une réalisation concrète} (Stage, Projet tutoré, Emploi précédent) : \textit{Contexte $\rightarrow$ Action (verbes d'action) $\rightarrow$ Résultat chiffré/qualifié}.
\item Mots-clés de l'offre : \cle{HACCP}, \cle{Maintenance préventive}, \cle{Automatisme Siemens}, \cle{Gestion de stock}, \cle{Soudure TIG/MIG}...
\end{itemize}
\item \textbf{Paragraphe 3 : Le « Nous » (La collaboration future / Savoir-être / Soft skills).}
Lien entre vos soft skills (rigueur, esprit d'équipe, adaptabilité, sécurité) et la culture de l'entreprise. \textit{Ex: « Mon sens de la sécurité, forgé sur le terrain, s'aligne parfaitement sur votre politique "Zéro Accident". »}
\end{enumerate}

\cle{Verbes d'action à privilégier (passé composé ou présent) :} \textbf{Conçu, Réalisé, Programmé, Diagnostiqué, Optimisé, Géré, Coordonné, Validé, Amélioré, Réduit (coûts/délais), Assuré.}
\cle{Verbes à bannir :} \textit{Aider, Participer, Apprendre, Voir, Faire, Être} (trop passifs).
\end{methode}

\begin{exoresolu}[Transformer des lignes de CV en paragraphes argumentés « Moi »]
\textbf{Énoncé :} Extrait du CV de \textbf{NGONO Alice}, candidate au poste d'\textbf{Assistant Comptable} (Offre : « Maîtrise Sage Saari, TVA, Rapprochement bancaire, Rigueur »).
\textit{Lignes CV :}
\begin{itemize}
\item Stage 6 mois (2023) : Cabinet EXPERTIS (Douala). Saisie pièces comptables (200/mois). Déclaration TVA mensuelle. Rapprochement bancaire.
\item Projet fin d'études : Mise en place procédure archivage numérique (Gain temps 30\%).
\item BTS Comptabilité-Gestion (Mention Bien).
\end{itemize}
Rédigez le paragraphe « Moi » (compétences techniques prouvées) de sa lettre de motivation.

\tcblower
\textbf{Solution détaillée :}

\textbf{Paragraphe « Moi » rédigé :}
\begin{quote}
\textit{Mon stage de six mois au cabinet d'expertise comptable EXPERTIS à Douala m'a permis de maîtriser concrètement le cycle complet de la tenue comptable sur \textbf{Sage Saari Comptabilité 100}. J'ai assuré en autonomie la \textbf{saisie et le lettrage de 200 pièces mensuelles}, la \textbf{préparation et la télédéclaration des liasses de TVA} dans les délais légaux, ainsi que le \textbf{rapprochement bancaire mensuel} de trois comptes courants. Par ailleurs, mon projet de fin d'études, consistant en la \textbf{conception d'une procédure d'archivage numérique}, a réduit de \textbf{30\% le temps de recherche des justificatifs}, témoignant de ma rigueur organisationnelle et de ma capacité à optimiser les processus.}
\end{quote}

\textbf{Analyse technique (Pourquoi ça marche au Bac) :}
\begin{center}
\begin{tabularx}{\linewidth}{|l|X|}\hline
\textbf{Critère} & \textbf{Application dans le texte} \\ \hline
\textbf{Verbes d'action} & \textbf{Assuré, Maîtrisé, Préparé, Conçu, Réduit} (Pas de « J'ai aidé à... »). \\ \hline
\textbf{Mots-clés offre} & \textbf{Sage Saari, TVA, Rapprochement bancaire, Rigueur} (Repris explicitement). \\ \hline
\textbf{Preuve chiffrée} & \textit{200 pièces/mois, 3 comptes, 30\% gain temps, 6 mois, 3 comptes.} \\ \hline
\textbf{Contexte $\rightarrow$ Action $\rightarrow$ Résultat} & Stage EXPERTIS (Contexte) $\rightarrow$ Saisie/TVA/Rappro (Actions) $\rightarrow$ Autonomie/Délais (Résultats). Projet FIN (Contexte) $\rightarrow$ Conception procédure (Action) $\rightarrow$ Gain 30\% (Résultat). \\ \hline
\textbf{Soft skill prouvée} & « Rigueur organisationnelle », « Capacité à optimiser » (Déduites des faits, pas juste affirmées). \\ \hline
\end{tabularx}
\end{center}
\textbf{Conclusion :} Ce paragraphe remplace avantageusement la liste du CV par une narration professionnelle orientée « Résultat ».
\end{exoresolu}

\courssub{Mobiliser les figures de style pour renforcer la persuasion (Gradation, Antithèse, Métonymie)}

\begin{methode}[Utiliser la Gradation, l'Antithèse et la Métonymie au service de l'argumentation]
\textbf{Objectif :} Élever le style (critère « Qualité de la langue / Expression » au barème) pour marquer l'esprit du correcteur/recruteur.
\medskip
\begin{enumerate}
\item \textbf{La Gradation (ou Climax) :} \textit{Ordre croissant de l'intensité.} Pour montrer une progression des compétences ou l'ampleur d'un résultat.
\begin{itemize}
\item \textit{Structure :} « J'ai \textbf{assisté}, puis \textbf{réalisé}, enfin \textbf{piloté}... »
\item \textit{Effet :} Montre l'évolution, la montée en charge, l'autonomie acquise.
\end{itemize}
\item \textbf{L'Antithèse :} \textit{Opposition de deux idées dans une même phrase (souvent « Mais / Cependant / Alors que »)}. Pour transformer une faiblesse apparente en force, ou contraster besoin/offre.
\begin{itemize}
\item \textit{Ex:} « \textbf{Alors que} le poste exige une réactivité immédiate, \textbf{mon expérience en maintenance curative} m'a forgé une capacité d'analyse sous pression. »
\item \textit{Ex (Jeune diplômé) :} « \textbf{Si} je suis jeune diplômé, \textbf{ma maîtrise des normes ISO 45001} est déjà opérationnelle. »
\end{itemize}
\item \textbf{La Métonymie :} \textit{Remplacer un concept par un terme lié (Contenu/Contenant, Lieu/Personne, Marque/Objet, Auteur/Œuvre).} Pour faire « pro », concis, technique.
\begin{itemize}
\item \textit{Contenant/Contenu :} « La \textbf{direction} a validé » (au lieu de « Le directeur et son équipe »).
\item \textit{Lieu/Personne :} « \textbf{Le terrain} m'a appris... » (Expérience pratique).
\item \textit{Marque/Objet (Très technique) :} « Je pilote des \textbf{Siemens S7-1200} » (Automates), « Je gère la \textbf{paie sous Sage} » (Logiciel).
\item \textit{Auteur/Œuvre :} « J'applique les préceptes de \textbf{Deming} » (Qualité).
\end{itemize}
\end{enumerate}
\cle{Attention :} Une figure de style par paragraphe max. Elle doit servir le sens, pas faire « littéraire ». Au Bac Tech, la \textbf{Gradation} (parcours) et la \textbf{Métonymie} (vocabulaire pro) sont les plus rentables.
\end{methode}

\begin{exoresolu}[Améliorer des phrases plates par l'intégration de figures de style]
\textbf{Énoncé :} Voici trois phrases « brutes » extraites d'une lettre de motivation pour un poste de \textbf{Chargé de Clientèle Banque}. Réécrivez chacune en y intégrant \textbf{une} figure de style imposée (indiquée entre parenthèses) pour renforcer l'impact professionnel.

\begin{enumerate}
\item \textit{(Gradation) :} « J'ai fait de l'accueil, puis j'ai géré des comptes, et maintenant je conseille des clients. »
\item \textit{(Antithèse) :} « Je n'ai pas beaucoup d'expérience mais je suis très motivé pour apprendre. »
\item \textit{(Métonymie) :} « Je connais bien les logiciels de la banque comme celui pour les crédits et celui pour l'épargne. »
\end{enumerate}

\tcblower
\textbf{Solution détaillée :}

\begin{enumerate}
\item \textbf{Gradation (Montée en expertise) :}
\begin{quote}
\textit{« De l'\textbf{accueil physique} de la clientèle à la \textbf{gestion courante des comptes courants}, j'ai progressé jusqu'au \textbf{conseil personnalisé en produits d'épargne et de crédit}, développant ainsi une vision globale de la relation client bancaire. »}
\end{quote}
\textit{Analyse :} Trois paliers : Accueil (exécution) $\rightarrow$ Gestion (technique) $\rightarrow$ Conseil (expertise/relationnel). Verbes : Accueillir $\rightarrow$ Gérer $\rightarrow$ Conseiller.

\item \textbf{Antithèse (Faiblesse apparente vs Force réelle) :}
\begin{quote}
\textit{« \textbf{Si mon parcours junior} peut sembler léger en années d'ancienneté, \textbf{ma certification AMF (Autorité des Marchés Financiers) et ma pratique intensive du simulateur de crédit en formation} garantissent une opérationnalité immédiate sur les dossiers complexes. »}
\end{quote}
\textit{Analyse :} Opposition « Junior / Léger » vs « Certification AMF / Pratique intensive / Opérationnalité immédiate ». Transforme le défaut (jeunesse) en atout (formation certifiée, outils maîtrisés).

\item \textbf{Métonymie (Vocabulaire métier / Outils) :}
\begin{quote}
\textit{« Ma maîtrise des \textbf{core banking} (notamment le module \textbf{crédit} et l'outil \textbf{épargne}) me permet de traiter un dossier de A à Z, de l'instruction au déblocage, sans intermédiaire. »}
\end{quote}
\textit{Analyse :} « Core banking » (Métonymie : Lieu/Système central pour désigner l'ERP bancaire), « Module crédit / Outil épargne » (Contenant/Contenu ou Marque/Module pour désigner les processus métiers). Remplace « logiciels de la banque » par le jargon pro.
\end{enumerate}
\end{exoresolu}

\courssub{Produire une lettre de motivation complète en situation d'examen (Probatoire)}

\begin{methode}[Démarche chronométrée (45 min) pour la production intégrée]
\textbf{Objectif :} Livrer une lettre conforme, argumentée, sans fautes graves, dans le temps imparti.
\medskip
\textbf{Étape 1 : Analyse du sujet \& de l'offre (5 min) -- \textbf{Brouillon obligatoire}}
\begin{itemize}
\item Identifier : \textbf{Poste exact}, \textbf{Entreprise}, \textbf{3 compétences techniques clés} exigées, \textbf{2 soft skills} implicites (ex: « terrain » $\rightarrow$ mobilité/résistance ; « équipe » $\rightarrow$ communication).
\item Cochez sur votre CV : \textbf{Quelles expériences prouvent ces 3 compétences ?} (Notez mots-clés, chiffres, verbes d'action).
\item Définir le « Vous » : 1 fait marquant entreprise (site web, actualité, réputation locale).
\end{itemize}

\textbf{Étape 2 : Squelette \& Mise en page (5 min) -- \textbf{Sur la copie finale}}
\begin{itemize}
\item Tracez la structure externe (Expéditeur, Destinataire, Lieu/Date, Objet, Appel).
\item Écrivez la \textbf{Formule de politesse finale} et \textbf{Pièces jointes} \textbf{tout de suite} (évite l'oubli en fin de temps).
\item Laissez l'espace pour le corps (Intro, §1, §2, §3, Conclusion).
\end{itemize}

\textbf{Étape 3 : Rédaction du corps (25 min) -- \textbf{Une phrase par idée}}
\begin{itemize}
\item \textbf{Intro (3 lignes) :} Vous (Entreprise) $\rightarrow$ Moi (Poste) $\rightarrow$ Lien (Source).
\item \textbf{§1 Technique (Le « Moi » fort) :} Expérience phare $\rightarrow$ Verbes d'action $\rightarrow$ Chiffres $\rightarrow$ Mots-clés offre. \textit{(Intégrer 1 figure de style : Gradation ou Métonymie)}.
\item \textbf{§2 Transversal (Le « Moi » soft/autre compétence) :} Autre stage/projet $\rightarrow$ Soft skill prouvée (Rigueur, Sécurité, Autonomie).
\item \textbf{§3 Projection (Le « Nous ») :} Adéquation valeurs / Enjeux entreprise $\rightarrow$ Ce que j'apporte demain.
\item \textbf{Conclusion (2 lignes) :} Disponibilité entretien + Valeur ajoutée future + Formule longue.
\end{itemize}

\textbf{Étape 4 : Rellecture active « Chasse aux points » (10 min) -- \textbf{Grille de 5 points}}
\begin{enumerate}
\item \textbf{Structure :} Objet ? Appel ? Politesse ? Pièces jointes ? Signature ? Une page ?
\item \textbf{Orthographe lexicale :} Mots techniques (HACCP, Saari, Siemens, TVA, BTS), Nom entreprise, Nom destinataire.
\item \textbf{Grammaire/Conjugaison :} Accords (participes passés avec avoir $\rightarrow$ COD avant), Temps (Passé composé pour faits passés, Présent pour compétences actuelles, Futur/Conditionnel pour projection).
\item \textbf{Syntaxe/Ponctuation :} Phrases courtes (max 2 lignes). Pas de « et » en chaîne. Virgules pour isoler les incises.
\item \textbf{Ton/Style :} « Je » assumé mais pas arrogant. Vocabulaire précis (pas de « trucs, choses, boulot »). Figures de style vérifiées.
\end{enumerate}
\end{methode}

\begin{exoresolu}[Production intégrée : Sujet type Probatoire]
\textbf{Sujet :} \textit{Vous êtes MESSI Arnaud, titulaire d'un BTS Électrotechnique (Option Énergies Renouvelables), résidant à Bafoussam. Vous répondez à l'offre ci-dessous. Rédigez votre lettre de motivation.}

\begin{center}
\begin{tabularx}{\linewidth}{|X|}\hline
\rowcolor{popblueL} \textbf{OFFRE D'EMPLOI -- SOCIETE SOLARIS CAMEROUN (Douala)} \\
\textbf{Poste :} Technicien Installation Solaire (H/F) -- CDD 12 mois renouvelable. \\
\textbf{Missions :} Étude dimensionnement, Installation panneaux/onduleurs, Mise en service, Maintenance préventive/curative, Relation client. \\
\textbf{Profil :} BTS/DUT Électrotechnique ou Énergies Renouvelables. Maîtrise normes NFC 15-100. Expérience terrain (stage/emploi) impérative. Permis B. Sens du service, autonomie, rigueur sécurité. \\
\textbf{Contact :} Directeur Technique, B.P. 789 Douala. Ref: SOL-2024/011. \\ \hline
\end{tabularx}
\end{center}

\textit{Informations candidat (CV) :}
\begin{itemize}
\item Stage 4 mois (2023) : ENERGY+ (Yaoundé). Dimensionnement 15 installations résidentielles (AutoCAD, PV*SOL). Installation terrain 8 sites (toiture, au sol). Mise en service, formation clients.
\item Projet fin d'études : Système solaire hybride (Solaire + Groupe + Batteries) pour centre de santé rural. Dimensionnement complet, schéma unifilaire, note de calcul.
\item Permis B (2022). Habilitation électrique BR/BC (Valide).
\end{itemize}

\tcblower
\textbf{Solution détaillée (Lettre complète type « Copie élève excellente ») :}

\begin{center}
\begin{tabularx}{\linewidth}{|X|}\hline
\textbf{MESSI Arnaud} \\
B.P. 105 Bafoussam \quad Tel: 6XX XX XX XX \quad Email: arnaud.messi@email.cm \\ \hline
\textbf{M. Le Directeur Technique} \\
\textbf{SOLARIS CAMEROUN} \\
B.P. 789 Douala \\ \hline
\textbf{Bafoussam, le 15 mars 2024} \\ \hline
\textbf{Objet : Candidature au poste de Technicien Installation Solaire (Ref: SOL-2024/011)} \\ \hline
\textbf{Monsieur le Directeur Technique,} \\ \hline
\textit{Votre positionnement de leader dans l'électrification rurale au Cameroun, illustré par le récent projet « Solarisation 100 Centres de Santé » dans la région de l'Est, confirme l'expertise technique et l'impact social qui animent SOLARIS. C'est donc avec une motivation profonde que je vous propose ma candidature au poste de Technicien Installation Solaire, suite à votre avis de recrutement publié sur le site de l'ANPE (Ref: SOL-2024/011).} \\ \hline
\textit{\textbf{Mon stage de quatre mois chez ENERGY+ à Yaoundé} m'a confronté à la réalité du terrain photovoltaïque. J'ai \textbf{réalisé le dimensionnement complet de 15 installations résidentielles} sous AutoCAD et PV*SOL, respectant la norme NFC 15-100. Sur le terrain, j'ai \textbf{installé et mis en service 8 sites} (toiture et au sol), de la pose des structures au câblage onduleurs/batteries, en assurant la formation des utilisateurs aux gestes de maintenance de premier niveau. Cette expérience, couronnée par l'obtention de mon habilitation électrique BR/BC, m'a forgé l'autonomie et la rigueur sécurité indispensables à vos chantiers.} \\ \hline
\textit{\textbf{Mon projet de fin d'études}, la conception d'un \textbf{système hybride Solaire/Groupe/Batteries pour un centre de santé rural}, a mobilisé mes compétences en \textbf{étude de faisabilité, calcul de charge, choix des onduleurs hybrides et rédaction de notes de calcul techniques}. Ce projet démontre ma capacité à appréhender des solutions complexes, garantissant la continuité de service -- un enjeu majeur pour vos clients professionnels.} \\ \hline
\textit{Titulaire du \textbf{Permis B} et habitué aux déplacements interurbains, je suis prêt à rejoindre vos équipes itinérantes dès la signature du contrat. Mon sens du service client, affiné lors des formations utilisateurs, et ma connaissance des contraintes climatiques locales (poussière, chaleur, orage) seront, j'en suis convaincu, des atouts pour la satisfaction et la fidélisation de votre clientèle.} \\ \hline
\textit{Disponible immédiatement pour un entretien technique, je souhaite vous exposer concrètement ma méthode de dimensionnement et de mise en service. \newline
Je vous prie d'agréer, Monsieur le Directeur Technique, l'expression de ma considération distinguée.} \\ \hline
\textbf{Arnaud MESSI} \hfill \textit{(Signature)} \\
\textbf{Pièces jointes :} CV détaillé, Copie BTS Électrotechnique, Attestation stage ENERGY+, Habilitation BR/BC, Permis B, Extrait acte de naissance. \\ \hline
\end{tabularx}
\end{center}

\textbf{Commentaire correcteur (Points forts) :}
\begin{itemize}
\item \textbf{Structure externe parfaite} (Alignements, Objet précis avec Ref, Politesse adaptée « Directeur Technique », Pièces jointes cohérentes).
\item \textbf{Intro :} « Vous » (Projet concret Solarisation 100 centres) + « Moi » (Poste exact) + « Lien » (ANPE/Ref).
\item \textbf{§1 (Technique) :} Gradation implicite (Dimensionnement $\rightarrow$ Installation $\rightarrow$ Mise en service $\rightarrow$ Formation). Mots-clés : NFC 15-100, AutoCAD, PV*SOL, Habilitation BR/BC, Autonomie, Sécurité. Chiffres : 15, 8, 4 mois.
\item \textbf{§2 (Projet/Complexité) :} Hybride, Centre de santé, Continuité de service (Enjeu client).
\item \textbf{§3 (Nous/Soft) :} Permis B, Déplacements, Service client, Contraintes locales.
\item \textbf{Conclusion :} Disponibilité + Proposition entretien technique + Politesse longue.
\item \textbf{Langue :} Zéro faute, verbes d'action, vocabulaire précis, ton professionnel.
\end{itemize}
\end{exoresolu}

\courssub{Relire et corriger sa lettre : Grille d'auto-évaluation « Zéro Faute »}

\begin{methode}[Check-list de relecture en 5 passes ciblées (Méthode « Passer au peigne fin »)]
\textbf{Objectif :} Éliminer les erreurs « bêtes » qui font chuter la note de « Qualité de la langue » (souvent 4 à 6 pts sur 20).
\medskip
\textbf{Ne relisez PAS tout d'un coup.} Faites \textbf{5 passages courts}, chacun avec \textbf{UN SEUL FOCUS} :

\begin{enumerate}
\item \textbf{Pass 1 : La « Squelette » (Mise en page \& Obligatoires) -- 1 min.}
\begin{itemize}
\item [$\square$] En-têtes (Expéditeur G / Destinataire D) ? Lieu/Date (D) ? Objet (G, Gras, Souligné, Ref) ?
\item [$\square$] Appel (G, Virgule) ? Politesse finale (G, Longue, Accord avec Appel) ? Signature (D) ? Pièces jointes (G) ?
\item [$\square$] Une seule page ? Paragraphes séparés (ligne blanche) ? Mares respectées ?
\end{itemize}
\item \textbf{Pass 2 : L'Orthographe Lexicale (Mots « images ») -- 2 min.}
\begin{itemize}
\item [$\square$] Noms propres : Entreprise, Destinataire, Ville, Logiciels (Sage, AutoCAD, PV*SOL), Normes (NFC 15-100, ISO, HACCP), Diplômes (BTS, DUT, Baccalauréat).
\item [$\square$] Mots techniques du métier (Onduleur, Dimensionnement, Rapprochement, Télédéclaration, Métonymies utilisées).
\item [$\square$] Homophones classiques : \textit{a/à, et/est, ou/où, ce/se, son/sont, ce/ceux, leur/leurs}.
\end{itemize}
\item \textbf{Pass 3 : La Grammaire / Conjugaison (Accords \& Temps) -- 3 min.}
\begin{itemize}
\item [$\square$] \textbf{Participes passés (avoir) :} COD avant ? \textit{« Les installations \textbf{que} j'ai \textbf{réalisées} »} (féminin pluriel). \textit{« J'ai \textbf{rédigé} une note »} (pas de COD avant).
\item [$\square$] \textbf{Participes passés (être/pronominal) :} Accord au sujet. \textit{« Elle \textbf{est partie} » / « Elles \textbf{se sont formées} »}.
\item [$\square$] \textbf{Temps :} Passé composé (Expériences passées) / Présent (Compétences actuelles, Vérités générales) / Futur/Conditionnel (Projection « Nous »).
\item [$\square$] \textbf{Accords GN :} \textit{« Mes compétences techniques »} (pluriel), \textit{« Une expérience significative »} (singulier).
\end{itemize}
\item \textbf{Pass 4 : La Syntaxe / Ponctuation / Fluidité -- 2 min.}
\begin{itemize}
\item [$\square$] Phrases $\le$ 2 lignes. Sujet - Verbe - Complément proches.
\item [$\square$] Ponctuation : Virgules pour incises / énumérations. Point-virgule pour lier 2 propositions indépendantes liées. Deux-points pour annoncer explication/liste.
\item [$\square$] Connecteurs logiques variés : \textit{Par ailleurs, En outre, De plus, Cependant, Ainsi, En effet, C'est pourquoi.}
\item [$\square$] Pas de répétitions (mots-outils : \textit{aussi, donc, alors} $\rightarrow$ varier).
\end{itemize}
\item \textbf{Pass 5 : Le « Test à voix basse » (Oreille) -- 2 min.}
\begin{itemize}
\item [$\square$] Lisez en murmurant (bougez les lèvres). L'oreille entend les lourdeurs, les manques de liaison, les phrases qui « ne tombent pas juste ».
\item [$\square$] Vérifiez le \textbf{Ton} : Professionnel, positif, concret. Pas de familier (\textit{« boulot, truc, super »}), pas de condescendance (\textit{« Je vous apprends que... »}).
\end{itemize}
\end{enumerate}
\cle{Astuce gain de temps :} Sur brouillon, codez vos passages : \textcircled{1} Squelette, \textcircled{2} Mots, \textcircled{3} Verbes, \textcircled{4} Phrases, \textcircled{5} Oreille. Cochez au fur et à mesure.
\end{methode}

\begin{exoresolu}[Correction guidée : « Chasse aux erreurs » sur une lettre défaillante]
\textbf{Énoncé :} Le texte ci-dessous contient \textbf{10 erreurs majeures} (structure, orthographe, grammaire, syntaxe, style) pénalisantes au Bac. Identifiez-les, classez-les par type (Pass 1 à 5), et proposez la correction.

\textit{« Monsieur le Directeur,
Objet : Candidature poste Technicien Maintenance
Par la presente, je vous ecris pour postuler a l'offre que j'ai vu sur Facebook. Je suis titulaire d'un BTS Maintenance Industrielle obtenue en 2023. Durant mon stage chez CAMRAIL, j'ai participer a la maintenance des locomotives diesel. J'ai aussi fait de la soudure et du diagnostic de pannes. Je suis très motivé et dynamique. Je veux bien travailler chez vous pour apprendre encore plus.
Dans l'attente, je vous salut.
Messi Arnaud
Pièces jointes : CV »}

\tcblower
\textbf{Solution détaillée : Analyse et Corrections}

\begin{center}
\begin{tabularx}{\linewidth}{|c|l|X|X|}\hline
\textbf{N°} & \textbf{Type (Pass)} & \textbf{Erreur identifiée dans le texte} & \textbf{Correction / Justification} \\ \hline
1 & \textbf{Structure (P1)} & \textbf{Manque En-têtes Expéditeur/Destinataire, Lieu/Date.} Objet mal placé (avant appel), pas gras/souligné, incomplet (pas de Ref). & \textbf{Ajouter bloc complet en haut.} Objet : \textbf{Objet : Candidature au poste de Technicien Maintenance Industrielle (Ref: ...)} \\ \hline
2 & \textbf{Structure (P1)} & \textbf{Appel incorrect :} « Monsieur le Directeur, » placé \textbf{après} l'Objet (doit être après Objet, avant corps). Manque virgule finale dans l'énoncé (mais présent ici). & \textbf{Ordre :} Expéditeur $\rightarrow$ Destinataire $\rightarrow$ Lieu/Date $\rightarrow$ \textbf{Objet} $\rightarrow$ \textbf{Monsieur le Directeur,} $\rightarrow$ Corps. \\ \hline
3 & \textbf{Style/Intro (P4/P5)} & \textbf{« Par la présente, je vous écris pour postuler... »} (Formule administrative vide, passive, « vu sur Facebook » peu pro). & \textbf{Intro « Vous/Moi/Lien » :} « Votre politique de maintenance prédictive sur le réseau ferroviaire Nord m'interpelle. Je vous adresse ma candidature au poste de Technicien Maintenance Industrielle, suite à votre annonce parue sur le site de l'ANPE (Ref: ...). » \\ \hline
4 & \textbf{Orthographe (P2)} & \textbf{« presente » (accent), « ecris » (accent), « a » (à), « vu » (vue -- COD « que » avant verbe avoir).} & \textbf{présente, écris, à, vue.} (Règle : COD « que » = « l'offre » fém. sg. $\rightarrow$ accord). \\ \hline
5 & \textbf{Grammaire (P3)} & \textbf{« obtenue »} (accord avec « BTS » masc. sg. $\rightarrow$ \textbf{obtenu}). & \textit{BTS} masculin $\rightarrow$ obtenu. \\ \hline
6 & \textbf{Grammaire (P3)} & \textbf{« j'ai participer »} (Infinitif au lieu de Participe Passé). & \textbf{participé} (verbe du 1er groupe : é). \\ \hline
7 & \textbf{Vocab/Style (P4/P5)} & \textbf{« J'ai aussi fait de la soudure et du diagnostic »} (Verbe « faire » faible, énumération vague, pas de verbes d'action, pas de chiffre). & \textbf{« J'ai \textbf{réalisé} des opérations de \textbf{soudure TIG/MIG} et \textbf{diagnostiqué} des pannes sur variateurs de fréquence, sur 3 séries de locomotives. »} \\ \hline
8 & \textbf{Style/Conclusion (P4/P5)} & \textbf{« Je veux bien travailler chez vous pour apprendre encore plus. »} (Trop familier, posture « demandeur/élève » pas « professionnel/expert »). & \textbf{Projection « Nous » :} « Intégrer vos équipes me permettrait de valoriser ma rigueur diagnostique tout en développant mon expertise sur vos équipements de traction modernes. » \\ \hline
9 & \textbf{Structure/Politesse (P1)} & \textbf{« Dans l'attente, je vous salut. »} (Formule de politesse \textbf{courte, incorrecte, familière}. « salut » = verbe 1er groupe $\rightarrow$ salue ; mais surtout formule interdite en lettre officielle). & \textbf{Formule longue obligatoire :} « Je vous prie d'agréer, Monsieur le Directeur, l'expression de ma considération distinguée. » \\ \hline
10 & \textbf{Structure (P1)} & \textbf{Signature :} « Messi Arnaud » (Prénom Nom au lieu de Nom PRÉNOM). Manque pièces jointes détaillées (Diplômes, Attestations). & \textbf{MESSI Arnaud (Signature). Pièces jointes : CV, Copie BTS, Attestation stage CAMRAIL, Habilitations, Extrait acte naissance.} \\ \hline
\end{tabularx}
\end{center}

\textbf{Bilan :} Sur 20 pts (Qualité langue + Contenu + Structure), cette version brute obtiendrait \textbf{$\le$ 5/20}. La version corrigée vise \textbf{16-18/20}. La relecture par « Pass » permet de corriger systématiquement ces 10 types d'erreurs récurrentes.
\end{exoresolu}

\begin{attention}[Les 5 erreurs qui coûtent des points au Bac (Probatoire / Baccalauréat Technique)]
\begin{enumerate}
\item \textbf{L'oubli ou l'inversion des mentions obligatoires de la structure externe (Pass 1).} \textit{Conséquence :} Perte de 2 à 4 pts sur « Respect des consignes / Structure » + impression de négligence immédiate. \textbf{Règle :} Écrire l'en-tête, l'objet, la politesse, la signature \textbf{AVANT} le corps de lettre.
\item \textbf{L'introduction « Par la présente... » ou « Je m'appelle X, j'ai Y ans... » (Pass 3/5).} \textit{Conséquence :} Note « Expression/Style » plafonnée (0-1/4). Le correcteur attend l'accroche « Vous » (Entreprise) $\rightarrow$ « Moi » (Poste) $\rightarrow$ « Lien » (Source).
\item \textbf{Le corps de lettre qui répète le CV (Liste diplômes/stages) sans argumenter (Pass 4/5).} \textit{Conséquence :} Note « Contenu/Argumentation » faible. \textbf{Règle :} 1 Paragraphe = 1 Compétence clé prouvée (Contexte $\rightarrow$ Action \textbf{verbe fort} $\rightarrow$ Résultat \textbf{chiffré}) + Mots-clés de l'offre.
\item \textbf{Les fautes d'accord du Participe Passé (avec avoir) et les homophones (a/à, et/est, ou/où, ce/se, son/sont) (Pass 2/3).} \textit{Conséquence :} Pénalité « Qualité de la langue » sévère (souvent -0,5 pt par faute grave, plafond bas). \textbf{Réflexe :} Pass 3 dédié \textbf{uniquement} aux accords PP (COD avant ?) et aux homophones au stylo rouge.
\item \textbf{La formule de politesse finale courte, mal accordée, ou absente (Pass 1).} \textit{Conséquence :} Perte de 1 à 2 pts « Conventions épistolaires » + signal négatif fort (manque de savoir-vivre pro). \textbf{Règle :} Formule \textbf{longue}, \textbf{adaptée au titre de l'appel} (Monsieur le Directeur $\rightarrow$ Monsieur le Directeur), \textbf{manuscrite si papier}.
\end{enumerate}
\cle{Mémo :} Une lettre de motivation technique réussie = \textbf{Structure impeccable} + \textbf{Vocabulaire métier précis} + \textbf{Preuves chiffrées} + \textbf{Zéro faute d'accord PP} + \textbf{Politesse parfaite}.
\end{attention}

\newpage

\coursec{Exercices}

\coursec{Produire une lettre officielle : la lettre de motivation}

\courssub{Découverte : lire des lettres de motivation authentiques}

\begin{aretenir}[Les textes supports de la leçon]
Cette séquence s'appuie sur cinq textes iconiques (lettres de motivation authentiques) et une affiche de sécurité. Leur lecture progressive permet d'identifier les codes de la correspondance professionnelle, la structure type et les procédés stylistiques efficaces.
\end{aretenir}

\begin{exercice}[QCM : la structure externe]
Pour chaque question, choisis la bonne réponse.
\begin{enumerate}
\item Dans une lettre de motivation, la formule d'appel correcte est :
\begin{enumerate}
\item[a)] Salut Monsieur ;
\item[b)] \cle{Monsieur le Directeur,} ;
\item[c)] Cher ami.
\end{enumerate}
\item L'objet de la lettre se place :
\begin{enumerate}
\item[a)] après la signature ;
\item[b)] \cle{avant la formule d'appel, à gauche} ;
\item[c)] dans l'en-tête du destinataire.
\end{enumerate}
\item La mention de la ville et de la date s'écrit :
\begin{enumerate}
\item[a)] en bas à gauche ;
\item[b)] \cle{en haut à droite, sous les coordonnées de l'expéditeur} ;
\item[c)] au milieu de la page.
\end{enumerate}
\item Une formule de politesse correcte est :
\begin{enumerate}
\item[a)] \cle{« Je vous prie d'agréer, Monsieur le Directeur, l'expression de ma haute considération. »} ;
\item[b)] « À plus ! » ;
\item[c)] « Merci d'avance, ciao. »
\end{enumerate}
\end{enumerate}
\end{exercice}

\begin{exercice}[Vrai ou faux justifié]
Dis si chaque affirmation est vraie ou fausse, puis justifie ta réponse en une phrase.
\begin{enumerate}
\item La lettre de motivation accompagne généralement le curriculum vitae (CV).
\item On peut commencer une lettre de motivation par : « Je viens par cette lettre vous demander du travail. »
\item Le schéma « Vous -- Moi -- Nous » signifie : parler d'abord de l'entreprise, puis de soi, puis de la collaboration possible.
\item La signature de l'expéditeur se place en haut à gauche de la lettre.
\item Une lettre de motivation peut contenir des abréviations comme « tkt » ou « svp ».
\end{enumerate}
\end{exercice}

\begin{exercice}[Définitions]
Définis en deux ou trois phrases chacun des termes suivants, puis donne un exemple pour les trois derniers :
\begin{enumerate}
\item la \cle{lettre de motivation} ;
\item la \cle{formule d'appel} ;
\item la \cle{gradation} ;
\item l'\cle{antithèse} ;
\item la \cle{métonymie}.
\end{enumerate}
\end{exercice}

\begin{exercice}[Les éléments de la structure externe]
Cite, dans l'ordre où ils apparaissent sur la page, les huit éléments de la structure externe d'une lettre de motivation, en précisant pour chacun sa place (haut/bas, gauche/droite).
\end{exercice}

\courssub{La structure externe de la lettre de motivation}

\begin{aretenir}[Les 8 éléments de la structure externe (ordre d'apparition)]
\begin{enumerate}
\item \cle{Coordonnées de l'expéditeur} : nom, prénom, adresse, téléphone, email — \textbf{en haut à gauche}.
\item \cle{Coordonnées du destinataire} : nom/qualité, structure, boîte postale, ville — \textbf{en haut à gauche, sous l'expéditeur}.
\item \cle{Lieu et date d'écriture} : ville, jour mois année — \textbf{en haut à droite}.
\item \cle{Objet de la lettre} : mention « Objet : » + titre précis + référence annonce — \textbf{à gauche, avant la formule d'appel}.
\item \cle{Formule d'appel} : « Monsieur le Directeur, » / « Madame la Directrice, » — \textbf{à gauche, sous l'objet}.
\item \cle{Corps de la lettre} : introduction, développement (1 à 3 paragraphes), conclusion — \textbf{au centre de la page}.
\item \cle{Formule de politesse finale} : formule longue adaptée au destinataire — \textbf{en bas à gauche}.
\item \cle{Signature} : prénom NOM manuscrit — \textbf{en bas à droite, sous la formule de politesse}.
\end{enumerate}
\end{aretenir}

\begin{attention}[Pièges fréquents]
\begin{itemize}
\item Ne pas confondre l'ordre : l'expéditeur est \emph{au-dessus} du destinataire.
\item La date est \emph{à droite}, jamais à gauche.
\item L'objet n'est pas une phrase complète, pas de point final.
\item La formule d'appel ne contient pas le nom propre (pas « Monsieur Dupont, »).
\item La signature est \emph{manuscrite} (même sur traitement de texte, on laisse l'espace).
\end{itemize}
\end{attention}

\begin{exercice}[Remettre dans l'ordre]
Voici, mélangés, les huit éléments de la structure externe d'une lettre de motivation :
\begin{center}
\begin{tabularx}{\linewidth}{|c|X|}\hline
\textbf{Lettre} & \textbf{Élément} \\ \hline
a & Signature \\ \hline
b & Nom, prénom, adresse et téléphone de l'expéditeur \\ \hline
c & Formule de politesse finale \\ \hline
d & Objet de la lettre \\ \hline
e & Nom et adresse du destinataire \\ \hline
f & Lieu et date d'écriture \\ \hline
g & Formule d'appel \\ \hline
h & Corps de la lettre \\ \hline
\end{tabularx}
\end{center}
\begin{enumerate}
\item Remets ces éléments dans l'ordre correct d'apparition sur la page (de 1 à 8).
\item Indique pour chaque élément sa position (en haut à gauche, en haut à droite, en bas à droite, etc.).
\end{enumerate}
\end{exercice}

\courssub{La structure interne : introduction et conclusion}

\begin{methode}[Rédiger une introduction efficace (schéma « Vous -- Moi -- Nous »)]
\begin{enumerate}
\item \textbf{Vous} (1 phrase) : Montrez que vous connaissez l'entreprise, citez l'annonce, valorisez l'organisme. Ex. : « Votre réputation d'excellence dans le secteur brassicole et votre annonce n\textsuperscript{o} BR/26/112 pour un poste de technicien de maintenance ont retenu toute mon attention. »
\item \textbf{Moi} (1--2 phrases) : Présentez votre profil (diplôme, spécialité, stage clé) en lien avec le poste. Ex. : « Titulaire d'un Baccalauréat technique en Maintenance des Systèmes Automatisés, j'ai consolidé mes compétences lors d'un stage de trois mois aux Brasseries du Cameroun, service maintenance préventive. »
\item \textbf{Nous} (1 phrase) : Projetez la collaboration. Ex. : « Je souhaite mettre ma rigueur et mon sens de l'équipe au service de la performance de vos lignes de production. »
\end{enumerate}
\end{methode}

\begin{methode}[Rédiger une conclusion professionnelle]
\begin{enumerate}
\item \textbf{Disponibilité} : « Je me tiens à votre disposition pour un entretien à votre convenance. »
\item \textbf{Relance} : « Dans l'attente de votre réponse, je vous prie d'agréer... »
\item \textbf{Formule de politesse} : Adaptée au destinataire (voir l'encadré ci-dessous).
\end{enumerate}
\end{methode}

\begin{aretenir}[Formules d'appel et de politesse : correspondance obligatoire]
\begin{center}
\begin{tabularx}{\linewidth}{|l|X|X|}\hline
\textbf{Destinataire} & \textbf{Formule d'appel} & \textbf{Formule de politesse (exemple)} \\ \hline
Monsieur le Directeur / Directeur Général & Monsieur le Directeur, & Je vous prie d'agréer, Monsieur le Directeur, l'expression de ma haute considération. \\ \hline
Madame la Directrice & Madame la Directrice, & Je vous prie d'agréer, Madame la Directrice, l'expression de ma haute considération. \\ \hline
Monsieur le Chef de Service / Personnel & Monsieur le Chef de Service, & Je vous prie d'agréer, Monsieur le Chef de Service, l'expression de mes salutations distinguées. \\ \hline
Madame, Monsieur (inconnu) & Madame, Monsieur, & Je vous prie d'agréer, Madame, Monsieur, l'expression de mes salutations distinguées. \\ \hline
\end{tabularx}
\end{center}
\textbf{Règle d'or} : Le titre dans la formule de politesse \textbf{reprend exactement} le titre de la formule d'appel.
\end{aretenir}

\begin{exercice}[Réécrire une introduction]
Voici l'introduction banale rédigée par un élève :
\begin{center}
\emph{« Je viens par cette lettre vous demander du travail. J'ai fini mes études et je suis disponible. »}
\end{center}
\begin{enumerate}
\item Releve deux faiblesses de cette introduction.
\item Réécris-la en une introduction accrocheuse de quatre à six lignes, en suivant le schéma « Vous -- Moi -- Nous », pour un poste de technicien de maintenance dans une brasserie.
\end{enumerate}
\end{exercice}

\courssub{Le corps de la lettre : convaincre avec style}

\begin{aretenir}[Figures de style au service de la persuasion]
\begin{definition}[Gradation]
Enchaînement de termes de force croissante (ou décroissante) pour amplifier une idée. \emph{Ex. : « Sérieux, rigoureux, infatigable... »}
\end{definition}
\begin{definition}[Antithèse]
Rapprochement de deux contraires pour créer un contraste saisissant. \emph{Ex. : « Le travail bien fait n'est pas un fardeau, mais une fierté. »}
\end{definition}
\begin{definition}[Métonymie]
Remplacement d'un mot par un autre qui lui est lié (contenant/contenu, lieu/habitants, institution/personnes). \emph{Ex. : « C'est toute la Brasserie qui m'a formé » (l'institution pour les formateurs/collectif).}
\end{definition}
\end{aretenir}

\begin{attention}[Registre de langue]
La lettre de motivation exige le \cle{registre soutenu} : pas de familier (« salut », « boulot », « tkt »), pas de tutoiement, pas d'abréviations SMS, orthographe et accords irréprochables. Les qualités s'\textbf{prouvent} par des faits (stage, projet, résultat), ne s'\textbf{affirment} pas (« je suis très doué »).
\end{attention}

\begin{exercice}[Figures de style]
Lis l'extrait suivant, tiré d'une lettre de motivation :
\begin{center}
\emph{« Sérieux, rigoureux, infatigable, j'ai appris dans vos ateliers que le travail bien fait n'est pas un fardeau, mais une fierté. C'est toute la Brasserie qui m'a formé, et c'est à elle que je veux aujourd'hui rendre service. »}
\end{center}
\begin{enumerate}
\item Releve dans cet extrait une gradation, une antithèse et une métonymie, en citant les mots exacts.
\item Explique l'effet produit par chacune de ces figures sur le lecteur.
\item Rédige trois phrases de candidature employant chacune l'une de ces trois figures de style.
\end{enumerate}
\end{exercice}

\begin{exercice}[Corriger une lettre fautive]
Voici un extrait de lettre rédigé par un élève :
\begin{center}
\emph{« Salut le chef, je t'écris pour le boulot que vous avez annoncé. Moi je suis très doué et tout le monde le dit. J'ai fais mon stage quelque part à Douala. Réponds-moi vite. »}
\end{center}
\begin{enumerate}
\item Releve huit fautes ou maladresses (registre familier, objet vague, fautes d'orthographe et d'accord, qualités affirmées sans preuve, formule finale incorrecte, tutoiement/vouvoiement instable, ponctuation, syntaxe).
\item Corrige chacune d'elles et justifie ta correction.
\item Réécris l'extrait complet sous une forme correcte et convaincante.
\end{enumerate}
\end{exercice}

\courssub{Production guidée : rédiger une lettre de motivation complète}

\begin{exercice}[Type Baccalauréat technique : lettre complète]
\textbf{Situation.} Tu es titulaire d'un Baccalauréat technique (spécialité de ton choix) et tu as effectué un stage de trois mois dans un atelier de maintenance. Tu lis l'annonce suivante :
\begin{center}
\emph{« Les Brasseries du Cameroun recrutent des techniciens de maintenance (réf. BR/26/112). Profil : Bac technique, sens de l'organisation, esprit d'équipe. Adresser lettre de motivation et CV à Monsieur le Directeur des Ressources Humaines, B.P. 4036 Douala. »}
\end{center}
\begin{enumerate}
\item Rédige la lettre de motivation complète que tu adresses à Monsieur le Directeur des Ressources Humaines, en respectant la structure externe (coordonnées, lieu et date, objet, formule d'appel, formule de politesse, signature) et la structure interne (introduction selon le schéma « Vous -- Moi -- Nous », corps en deux ou trois paragraphes, conclusion).
\item Dans le corps de ta lettre, présente au moins deux qualités, chacune justifiée par une preuve concrète (formation, stage, expérience).
\item Emploie au moins une \cle{gradation} et une \cle{antithèse} dans ta lettre.
\end{enumerate}
\emph{Barème indicatif : présentation et structure externe /5 ; contenu et argumentation /8 ; correction de la langue /7.}
\end{exercice}

\begin{exercice}[Type Probatoire : introduction et conclusion]
\textbf{Situation.} Tu réponds à l'annonce de stage suivante, parue dans \emph{Cameroon Tribune} :
\begin{center}
\emph{« La SODECOTON de Garoua offre des stages de perfectionnement aux élèves de classes de Première et Terminale techniques, filière cotonnière, pour la période de juillet à septembre. Écrire à Monsieur le Chef du Personnel, B.P. 302 Garoua. »}
\end{center}
\begin{enumerate}
\item Rédige l'en-tête complet de ta lettre (coordonnées de l'expéditeur, coordonnées du destinataire, lieu et date, objet précis).
\item Rédige l'introduction de ta lettre (quatre à six lignes) en suivant le schéma « Vous -- Moi -- Nous » et avec la formule d'appel appropriée.
\item Rédige la conclusion de ta lettre (trois à cinq lignes), comportant une demande d'entretien ou de réponse, suivie de la formule de politesse adaptée au destinataire.
\item Justifie en deux phrases le choix de ta formule d'appel et de ta formule de politesse.
\end{enumerate}
\emph{Barème indicatif : en-tête /4 ; introduction /6 ; conclusion /6 ; langue /4.}
\end{exercice}

\courssub{Compte-rendu, évaluation et remédiation}

\begin{exercice}[Communication par l'image et texte descriptif]
\textbf{Situation.} Dans l'atelier de ton lycée technique est affichée une affiche de sécurité : on y voit un ouvrier portant casque, gants et lunettes, avec le slogan « Un geste de prudence vaut mille regrets » et la consigne « Port des équipements obligatoire ».
\begin{enumerate}
\item Analyse l'affiche en identifiant : l'image, le slogan et la consigne ; puis explique le rôle de chacun dans la communication.
\item Explique en trois phrases comment cette affiche persuade le lecteur (appel aux émotions, à la raison, à la peur ou au sens des responsabilités).
\item Le slogan contient une figure de style : laquelle ? Justifie ta réponse.
\item Rédige un court paragraphe descriptif (huit à douze lignes) décrivant l'affiche de façon objective (éléments visuels, couleurs, disposition), puis subjectif (impression produite sur le lecteur).
\end{enumerate}
\emph{Barème indicatif : analyse /6 ; figure de style /4 ; paragraphe descriptif /10.}
\end{exercice}

\begin{aretenir}[Bilan : checklist avant de rendre sa lettre]
\begin{itemize}
\item[$\square$] Structure externe complète et bien placée (8 éléments).
\item[$\square$] Objet précis avec référence annonce.
\item[$\square$] Formule d'appel et formule de politesse cohérentes (même titre).
\item[$\square$] Introduction « Vous -- Moi -- Nous » (4--6 lignes).
\item[$\square$] Corps : 2--3 paragraphes, qualités \textbf{prouvées} par exemples.
\item[$\square$] Au moins une figure de style (gradation, antithèse, métonymie).
\item[$\square$] Conclusion : disponibilité + relance + formule de politesse.
\item[$\square$] Signature manuscrite (espace prévu).
\item[$\square$] Relecture : orthographe, grammaire, ponctuation, registre soutenu.
\end{itemize}
\end{aretenir}

\newpage

\coursec{Corrigés des exercices}

\coursec{Exercices corrigés et remédiation — La lettre de motivation}

\courssub{Étape 2 : La structure externe de la lettre de motivation}

\begin{corrige}[Exercice 1 — QCM : la structure externe]
\begin{enumerate}
\item \textbf{Réponse b)} : « \cle{Monsieur le Directeur,} ». La formule d'appel d'une lettre officielle mentionne la \emph{qualité} du destinataire, jamais son nom propre, et exclut toute familiarité (« Salut », « Cher ami » sont réservés à la correspondance privée).
\item \textbf{Réponse b)} : l'objet se place \cle{avant la formule d'appel, à gauche}. Il annonce le sujet de la lettre avant même que celle-ci ne commence ; il ne peut donc figurer ni après la signature ni dans l'en-tête du destinataire.
\item \textbf{Réponse b)} : la ville et la date s'écrivent \cle{en haut à droite}, sous les coordonnées de l'expéditeur. C'est une convention fixe de la correspondance administrative et professionnelle.
\item \textbf{Réponse a)} : « Je vous prie d'agréer, Monsieur le Directeur, l'expression de ma haute considération. » C'est la formule de politesse longue, soutenue, qui reprend le titre de la formule d'appel. Les formules b) et c) relèvent du registre familier, interdit dans une lettre officielle.
\end{enumerate}
\textbf{Point de vigilance} : retenir la règle d'or — la formule de politesse reprend \emph{exactement} le titre employé dans la formule d'appel.
\end{corrige}

\begin{corrige}[Exercice 4 — Les éléments de la structure externe]
Dans l'ordre d'apparition sur la page :
\begin{enumerate}
\item \textbf{Coordonnées de l'expéditeur} (nom, prénom, adresse, téléphone, email) — \emph{en haut à gauche}.
\item \textbf{Coordonnées du destinataire} (qualité, structure, boîte postale, ville) — \emph{en haut à gauche, sous l'expéditeur}.
\item \textbf{Lieu et date d'écriture} (« Douala, le 15 mars 2026 ») — \emph{en haut à droite}.
\item \textbf{Objet de la lettre} (« Objet : candidature au poste de... », sans point final) — \emph{à gauche, avant la formule d'appel}.
\item \textbf{Formule d'appel} (« Monsieur le Directeur, ») — \emph{à gauche, sous l'objet}.
\item \textbf{Corps de la lettre} (introduction, développement, conclusion) — \emph{au centre de la page}.
\item \textbf{Formule de politesse finale} — \emph{en bas à gauche}.
\item \textbf{Signature manuscrite} (prénom NOM) — \emph{en bas à droite}.
\end{enumerate}
\end{corrige}

\begin{corrige}[Exercice 5 — Remettre dans l'ordre]
\textbf{1. Ordre correct d'apparition (de 1 à 8) :}
\begin{center}
\begin{tabularx}{\linewidth}{|c|X|l|}\hline
\textbf{Ordre} & \textbf{Élément} & \textbf{Position} \\ \hline
1 & b — Nom, prénom, adresse et téléphone de l'expéditeur & en haut à gauche \\ \hline
2 & e — Nom et adresse du destinataire & en haut à gauche, sous l'expéditeur \\ \hline
3 & f — Lieu et date d'écriture & en haut à droite \\ \hline
4 & d — Objet de la lettre & à gauche, avant la formule d'appel \\ \hline
5 & g — Formule d'appel & à gauche, sous l'objet \\ \hline
6 & h — Corps de la lettre & au centre de la page \\ \hline
7 & c — Formule de politesse finale & en bas à gauche \\ \hline
8 & a — Signature & en bas à droite \\ \hline
\end{tabularx}
\end{center}
\textbf{2. Justification} : la lettre se lit de haut en bas ; l'expéditeur précède toujours le destinataire ; la date est à droite ; l'objet annonce le contenu avant l'appel ; la signature clôt la lettre en bas à droite, sous la formule de politesse.
\textbf{Point de vigilance} : ne pas inverser expéditeur et destinataire, ni placer la date à gauche.
\end{corrige}

\courssub{Étape 3 : La structure interne — Introduction et conclusion}

\begin{corrige}[Exercice 2 — Vrai ou faux justifié]
\begin{enumerate}
\item \textbf{Vrai.} La lettre de motivation accompagne le CV : le CV présente le parcours de façon synthétique, la lettre explique et argumente la candidature.
\item \textbf{Faux.} Cette formulation est banale, plate et centrée sur le besoin du candidat ; une introduction efficace suit le schéma « Vous -- Moi -- Nous » et valorise d'abord l'entreprise.
\item \textbf{Vrai.} « Vous » = l'entreprise (l'annonce, ses valeurs), « Moi » = le profil du candidat (diplôme, stage), « Nous » = la collaboration future projetée.
\item \textbf{Faux.} La signature se place \emph{en bas à droite}, sous la formule de politesse ; en haut à gauche figurent les coordonnées de l'expéditeur.
\item \textbf{Faux.} La lettre de motivation exige le registre soutenu : toute abréviation de type SMS (« tkt », « svp ») est une faute de registre qui discrédite la candidature.
\end{enumerate}
\end{corrige}

\begin{corrige}[Exercice 3 — Définitions]
\begin{enumerate}
\item \textbf{La lettre de motivation} : lettre officielle par laquelle un candidat sollicite un emploi, un stage ou une formation. Elle accompagne le CV et vise à convaincre le destinataire de l'adéquation entre le profil du candidat et le poste. Elle obéit à une structure externe (présentation) et interne (introduction, corps, conclusion) codifiées.
\item \textbf{La formule d'appel} : expression de courtoisie placée au début du corps de la lettre, qui désigne le destinataire par sa qualité (« Monsieur le Directeur, », « Madame la Directrice, »). Elle ne contient jamais de nom propre et conditionne le choix de la formule de politesse finale.
\item \textbf{La gradation} : figure de style consistant à enchaîner des termes de force croissante (ou décroissante) pour amplifier une idée. \emph{Exemple : « Sérieux, rigoureux, infatigable, je m'engage pleinement dans chaque tâche. »}
\item \textbf{L'antithèse} : figure de style qui rapproche deux termes ou deux idées contraires pour créer un contraste frappant. \emph{Exemple : « Le travail bien fait n'est pas un fardeau, mais une fierté. »}
\item \textbf{La métonymie} : figure de style qui remplace un mot par un autre qui lui est lié par un rapport logique (le contenant pour le contenu, l'institution pour ses membres). \emph{Exemple : « C'est toute la Brasserie qui m'a formé » (l'entreprise pour ses formateurs).}
\end{enumerate}
\end{corrige}

\begin{corrige}[Exercice 6 — Réécrire une introduction]
\textbf{1. Deux faiblesses de l'introduction banale :}
\begin{itemize}
\item Elle est centrée sur le candidat (« je viens... vous demander ») et ignore l'entreprise : aucun élément « Vous » (annonce, réputation, poste visé).
\item Elle est vague et sans argument : « j'ai fini mes études » ne précise ni diplôme, ni compétence, ni lien avec le poste ; elle ne suscite aucun intérêt.
\end{itemize}
\textbf{2. Réécriture selon le schéma « Vous -- Moi -- Nous » :}
\begin{center}
\emph{« Votre réputation d'excellence dans le secteur brassicole et votre annonce n\textsuperscript{o} BR/26/112, relative au recrutement de techniciens de maintenance, ont retenu toute mon attention. Titulaire d'un Baccalauréat technique en Maintenance des Systèmes Automatisés, j'ai consolidé mes compétences lors d'un stage de trois mois dans un atelier de maintenance industrielle, où j'ai participé à la maintenance préventive des lignes de production. Je souhaite aujourd'hui mettre ma rigueur et mon esprit d'équipe au service de la performance de vos installations. »}
\end{center}
\textbf{Point de vigilance} : vérifier la présence des trois temps (Vous / Moi / Nous), la précision de la référence de l'annonce et le registre soutenu.
\end{corrige}

\begin{corrige}[Exercice 10 — Type Probatoire : introduction et conclusion]
\textbf{1. En-tête complet :}
\begin{center}
\begin{tabularx}{\linewidth}{X}
MBALLA Sylvie \\
Lycée technique de Garoua \\
B.P. 45 Garoua \\
Tél. : 6 99 00 00 00 \\[4pt]
\textbf{À Monsieur le Chef du Personnel} \\
SODECOTON \\
B.P. 302 Garoua \\
\hfill Garoua, le 10 juin 2026 \\[4pt]
\textbf{Objet :} demande de stage de perfectionnement, filière cotonnière (juillet--septembre 2026) \\
\end{tabularx}
\end{center}
\textbf{2. Introduction (schéma « Vous -- Moi -- Nous ») :}
\begin{center}
\emph{« Monsieur le Chef du Personnel, votre entreprise, pilier de la filière cotonnière camerounaise, offre aux élèves techniciens une opportunité précieuse de perfectionnement, annoncée dans \emph{Cameroon Tribune}. Élève en classe de Première technique au Lycée technique de Garoua, je me suis déjà familiarisé avec les machines textiles lors de nos travaux pratiques d'atelier. Je souhaite mettre ce premier savoir-faire au service de vos équipes et apprendre auprès de vos techniciens expérimentés. »}
\end{center}
\textbf{3. Conclusion :}
\begin{center}
\emph{« Je me tiens à votre disposition pour tout entretien que vous jugerez utile et pour fournir les pièces complémentaires de mon dossier. Dans l'attente d'une réponse favorable, je vous prie d'agréer, Monsieur le Chef du Personnel, l'expression de mes salutations distinguées. »}
\end{center}
\textbf{4. Justification :} le destinataire est désigné par sa qualité (« Monsieur le Chef du Personnel »), sans nom propre, conformément à l'annonce ; la formule de politesse reprend exactement ce titre, avec « salutations distinguées », formule adaptée à un chef de service (la « haute considération » étant plutôt réservée à un directeur).
\textbf{Barème indicatif :} en-tête /4 ; introduction /6 ; conclusion /6 ; langue /4.
\textbf{Points de vigilance :} objet précis (nature du stage, filière, période) ; présence des trois temps Vous--Moi--Nous ; cohérence appel/politesse.
\end{corrige}

\courssub{Étape 4 : Le corps de la lettre — Convaincre avec style}

\begin{corrige}[Exercice 7 — Figures de style]
\textbf{1. Relevé des figures :}
\begin{itemize}
\item \textbf{Gradation} : « \emph{Sérieux, rigoureux, infatigable} » — trois adjectifs de force croissante.
\item \textbf{Antithèse} : « \emph{n'est pas un fardeau, mais une fierté} » — opposition entre « fardeau » et « fierté ».
\item \textbf{Métonymie} : « \emph{toute la Brasserie qui m'a formé} » — l'institution (la Brasserie) désigne les personnes qui y travaillent (les formateurs, le collectif).
\end{itemize}
\textbf{2. Effets produits :}
\begin{itemize}
\item La \textbf{gradation} amplifie le portrait du candidat : chaque terme renforce le précédent et frappe la mémoire du lecteur.
\item L'\textbf{antithèse} crée un contraste valorisant : le travail, présenté d'abord sous son aspect négatif possible, est renversé en valeur positive, signe d'une motivation profonde.
\item La \textbf{métonymie} exprime la gratitude et l'appartenance : le candidat se présente comme le produit de l'entreprise, ce qui flatte le destinataire et annonce la fidélité (« rendre service »).
\end{itemize}
\textbf{3. Production (exemple de réponse) :}
\begin{itemize}
\item Gradation : « Ponctuel, appliqué, perfectionniste, je ne rends jamais un travail inachevé. »
\item Antithèse : « Là où d'autres voient une contrainte, je vois une chance de progresser. »
\item Métonymie : « Vos ateliers m'ont tout appris ; c'est à eux que je veux désormais apporter mon savoir-faire. »
\end{itemize}
\end{corrige}

\begin{corrige}[Exercice 8 — Corriger une lettre fautive]
\textbf{1. Relevé des huit fautes et corrections :}
\begin{center}
\begin{tabularx}{\linewidth}{|c|X|X|}\hline
\textbf{N\textsuperscript{o}} & \textbf{Faute / maladresse} & \textbf{Correction justifiée} \\ \hline
1 & « Salut le chef » : appel familier & « Monsieur le Directeur, » : formule d'appel officielle, avec la qualité du destinataire \\ \hline
2 & « je t'écris » : tutoiement & Vouvoiement obligatoire : « je vous écris » \\ \hline
3 & « pour le boulot » : registre familier et objet vague & « au poste de technicien de maintenance (réf. BR/26/112) » : objet précis, registre soutenu \\ \hline
4 & « je suis très doué et tout le monde le dit » : qualité affirmée sans preuve & Prouver par un fait : « j'ai obtenu la meilleure moyenne de ma série au Baccalauréat technique » \\ \hline
5 & « J'ai fais » : faute de conjugaison & « J'ai fait » (participe passé de \emph{faire}) \\ \hline
6 & « quelque part à Douala » : imprécision & Nommer la structure : « dans l'atelier de maintenance de la société X, à Douala » \\ \hline
7 & « Réponds-moi vite » : impératif familier, tutoiement incohérent avec « vous avez annoncé » & « Je me tiens à votre disposition pour un entretien. » \\ \hline
8 & Absence de formule de politesse & « Je vous prie d'agréer, Monsieur le Directeur, l'expression de ma haute considération. » \\ \hline
\end{tabularx}
\end{center}
\textbf{2. Réécriture complète :}
\begin{center}
\emph{« Monsieur le Directeur, je me permets de vous adresser ma candidature au poste de technicien de maintenance, annoncé sous la référence BR/26/112. Titulaire d'un Baccalauréat technique, j'ai fait mon stage de fin de formation dans l'atelier de maintenance de la société X, à Douala, où j'ai participé au suivi préventif des machines. Sérieux et rigoureux, j'ai obtenu la satisfaction de mon maître de stage pour la qualité de mes interventions. Je me tiens à votre disposition pour un entretien à votre convenance. Dans l'attente de votre réponse, je vous prie d'agréer, Monsieur le Directeur, l'expression de ma haute considération. »}
\end{center}
\textbf{Point de vigilance} : stabilité du vouvoiement, qualités prouvées par des faits, cohérence appel/politesse.
\end{corrige}

\courssub{Étape 5 : Production guidée — Rédiger une lettre de motivation complète}

\begin{corrige}[Exercice 9 — Type Baccalauréat technique : lettre complète]
\textbf{Modèle de lettre attendue :}
\begin{center}
\begin{tabularx}{\linewidth}{X}
NGOUESSI Arnaud \\
Rue des Ateliers, Akwa \\
B.P. 5120 Douala \\
Tél. : 6 90 00 00 00 \\
arnaud.ngouessi@mail.cm \\[4pt]
\textbf{À Monsieur le Directeur des Ressources Humaines} \\
Les Brasseries du Cameroun \\
B.P. 4036 Douala \\
\hfill Douala, le 15 mars 2026 \\[4pt]
\textbf{Objet :} candidature au poste de technicien de maintenance (réf. BR/26/112) \\[4pt]
Monsieur le Directeur, \\
\end{tabularx}
\end{center}
\emph{« Votre réputation d'excellence dans le secteur brassicole et votre annonce n\textsuperscript{o} BR/26/112, relative au recrutement de techniciens de maintenance, ont retenu toute mon attention. Titulaire d'un Baccalauréat technique en Maintenance des Systèmes Automatisés, j'ai consolidé ma formation par un stage de trois mois dans un atelier de maintenance industrielle. Je souhaite mettre mes compétences au service de la performance de vos lignes de production.}

\emph{Mon stage m'a permis de participer à la maintenance préventive des machines : graissage, contrôle des capteurs, remplacement de pièces usées. Sérieux, rigoureux, infatigable, j'ai appris qu'une panne évitée n'est pas une dépense, mais un investissement. Mon maître de stage m'a confié, dès le deuxième mois, la conduite seule des rondes de contrôle.}

\emph{Doté d'un réel esprit d'équipe, j'ai collaboré chaque jour avec les électromécaniciens et les conducteurs de ligne ; cette coopération m'a valu une mention élogieuse dans mon rapport de stage. Je sais écouter, transmettre et agir vite lorsque la production l'exige.}

\emph{Je me tiens à votre disposition pour un entretien à votre convenance. Dans l'attente de votre réponse, je vous prie d'agréer, Monsieur le Directeur, l'expression de ma haute considération. »}
\begin{flushright}
\emph{Signature} \\
NGOUESSI Arnaud
\end{flushright}
\textbf{Barème indicatif :}
\begin{itemize}
\item Présentation et structure externe (8 éléments bien placés) : /5
\item Contenu et argumentation (schéma Vous--Moi--Nous, 2 qualités prouvées, gradation et antithèse repérables) : /8
\item Correction de la langue (orthographe, accords, ponctuation, registre soutenu) : /7
\end{itemize}
\textbf{Points de vigilance :} la gradation « Sérieux, rigoureux, infatigable » et l'antithèse « n'est pas une dépense, mais un investissement » doivent être identifiables ; chaque qualité (rigueur, esprit d'équipe) est justifiée par une preuve concrète ; la formule de politesse reprend le titre de l'appel.
\end{corrige}

\courssub{Étape 6 : Compte-rendu, évaluation et remédiation}

\begin{corrige}[Exercice 11 — Communication par l'image et texte descriptif]
\textbf{1. Analyse de l'affiche :}
\begin{itemize}
\item \textbf{L'image} : un ouvrier portant casque, gants et lunettes. Rôle : montrer le modèle à imiter, rendre la consigne concrète et visible d'un coup d'œil.
\item \textbf{Le slogan} : « Un geste de prudence vaut mille regrets ». Rôle : frapper la mémoire, résumer le message en une formule courte et rythmée.
\item \textbf{La consigne} : « Port des équipements obligatoire ». Rôle : énoncer la règle de façon explicite et impérative ; elle transforme le message en obligation.
\end{itemize}
\textbf{2. Comment l'affiche persuade :} elle fait d'abord appel à la \emph{raison} en montrant la solution (les équipements de protection). Elle fait ensuite appel à la \emph{peur} et aux \emph{émotions} par le mot « regrets », qui évoque l'accident et ses conséquences. Enfin, elle mobilise le \emph{sens des responsabilités} : c'est au lecteur d'accomplir le « geste de prudence ».
\textbf{3. Figure de style :} le slogan contient une \cle{antithèse} : il oppose « un geste de prudence » (léger, préventif, positif) à « mille regrets » (lourds, irréparables, négatifs). Le contraste entre la facilité de la prévention et la gravité des conséquences rend le message frappant. (On peut aussi noter l'hyperbole « mille regrets », qui amplifie l'opposition.)
\textbf{4. Paragraphe descriptif (exemple) :}
\begin{center}
\emph{« L'affiche, de format vertical, est placardée à l'entrée de l'atelier, à hauteur des yeux. En son centre, un ouvrier en bleu de travail porte un casque jaune, des gants épais et des lunettes de protection ; son regard, tourné vers le lecteur, semble l'interpeller directement. En haut, sur fond rouge, la consigne « Port des équipements obligatoire » s'écrit en lettres blanches capitales. En bas, le slogan « Un geste de prudence vaut mille regrets » clôt la composition. Les couleurs vives — rouge, jaune, bleu — contrastent avec le fond clair et captent immédiatement l'attention. À la vue de cette affiche, on ressent d'abord une mise en garde, presque une inquiétude ; puis le personnage, calme et correctement équipé, rassure : la sécurité est possible, elle ne tient qu'à un geste. »}
\end{center}
\textbf{Barème indicatif :} analyse des trois composantes /6 ; identification et justification de la figure de style /4 ; paragraphe descriptif (objectif puis subjectif, 8--12 lignes, lexique précis) /10.
\textbf{Points de vigilance :} distinguer clairement description objective (éléments visuels, couleurs, disposition) et impression subjective ; justifier la figure de style par les termes exacts du slogan.
\end{corrige}

\newpage

\coursec{Fiche bilan}

\begin{popfigure}
\begin{tikzpicture}[mindmap, grow cyclic, every node/.style={concept, font=\small},
root concept/.append style={concept color=popblue, font=\bfseries},
level 1/.append style={level distance=3.6cm, sibling angle=51.4},
level 1 concept/.append style={font=\footnotesize}]
\node[root concept]{La lettre de motivation}
child[concept color=poporange]{node{Définition et objectif : convaincre un destinataire précis}}
child[concept color=popgreen]{node{Structure externe : en-tête, objet, formules, signature}}
child[concept color=poppurple]{node{Introduction : accroche et annonce de la candidature}}
child[concept color=poppink]{node{Corps : diplômes, compétences, motivations}}
child[concept color=popgold]{node{Conclusion : demande d'entretien et salutations}}
child[concept color=popblueL]{node{Style : gradation, antithèse, métonymie}}
child[concept color=popmuted]{node{Méthode : préparer, rédiger, relire, améliorer}};
\end{tikzpicture}
\legende{Carte mentale de la leçon : les sept points essentiels pour produire une lettre de motivation réussie.}
\end{popfigure}

\begin{bilan}
\textbf{Définitions clés}
\begin{itemize}
\item La \cle{lettre de motivation} : lettre officielle par laquelle un candidat sollicite un emploi, un stage ou une admission, en mettant en valeur ses qualités.
\item La \cle{structure externe} : la présentation matérielle de la lettre (coordonnées, lieu et date, objet, formules).
\item La \cle{structure interne} : l'organisation du texte (introduction, corps, conclusion).
\item La \cle{gradation} : énumération de termes de force croissante ou décroissante.
\item L'\cle{antithèse} : rapprochement de deux termes de sens opposés.
\item La \cle{métonymie} : remplacement d'un mot par un autre lié par un rapport logique (la partie pour le tout, le contenant pour le contenu).
\end{itemize}

\textbf{La structure externe à connaître par cœur}
\begin{center}
\begin{tabularx}{\linewidth}{|l|X|}
\hline
\rowcolor{popblueL} \textbf{Élément} & \textbf{Contenu et place} \\
\hline
Coordonnées de l'expéditeur & Nom, adresse, téléphone (en haut à gauche) \\
\hline
Coordonnées du destinataire & Nom, fonction, adresse (en haut à droite) \\
\hline
Lieu et date & Ville, suivi de la date (à droite) \\
\hline
Objet & Formule brève : « Objet : demande d'emploi » \\
\hline
Formule d'appel & « Monsieur le Directeur, », « Madame, » \\
\hline
Formule de politesse & « Veuillez agréer, Monsieur, l'expression de ma haute considération. » \\
\hline
Signature & Nom et signature (en bas à droite) \\
\hline
\end{tabularx}
\end{center}

\textbf{Méthode express de rédaction}
\begin{enumerate}
\item \textbf{Préparer} : analyser la situation (qui écrit ? à qui ? pour quoi ?), lister diplômes et qualités.
\item \textbf{Introduction} : accrocher le lecteur et annoncer clairement la candidature.
\item \textbf{Corps} : présenter ses diplômes, ses compétences et ses motivations en paragraphes liés.
\item \textbf{Conclusion} : solliciter un entretien et formuler les salutations.
\item \textbf{Relire et améliorer} : corriger l'orthographe, la syntaxe, la présentation.
\end{enumerate}
\end{bilan}

\begin{aretenir}[Checklist avant le Bac]
\begin{itemize}
\item $\square$ Je sais définir la lettre de motivation et son objectif.
\item $\square$ Je place correctement les coordonnées de l'expéditeur et du destinataire.
\item $\square$ Je rédige le lieu, la date et l'objet sans faute.
\item $\square$ Je choisis une formule d'appel adaptée au destinataire.
\item $\square$ Mon introduction accroche le lecteur et annonce ma candidature.
\item $\square$ Mon corps de lettre présente diplômes, compétences et motivations.
\item $\square$ Je sollicite un entretien dans la conclusion.
\item $\square$ J'emploie une formule de politesse correcte et je signe.
\item $\square$ Je sais reconnaître et employer gradation, antithèse et métonymie.
\item $\square$ Je relie mes paragraphes par des connecteurs logiques.
\item $\square$ Je relis ma lettre pour corriger orthographe et présentation.
\item $\square$ Je respecte le ton soutenu et courtois d'une lettre officielle.
\end{itemize}
\end{aretenir}

\findecours

\end{document}
